\documentclass[a4paper,11pt]{article}
\pdfoutput=1
\usepackage{jheppub}   % JHEP class file (SISSA Medialab, v1.1227), copied into this directory
\usepackage[T1]{fontenc}
\usepackage{amsmath,amssymb,amsfonts,mathtools}
\usepackage{braket}
\usepackage{booktabs}
\usepackage{array}
\allowdisplaybreaks[2]
\usepackage{xcolor}
\usepackage{tikz}
\usetikzlibrary{calc}
\usetikzlibrary{decorations.pathmorphing}
\usetikzlibrary{decorations.markings}
\usetikzlibrary{decorations.text}
\usetikzlibrary{decorations.pathreplacing}
\usepackage{tikzit}
\input{stiles.tikzstyles}
\usepackage{longtable}
\usepackage[most]{tcolorbox}

% ---------------------------------------------------------------------
% Reader guidance: purpose paragraph at the start of a section, boxed summary at its end
% ---------------------------------------------------------------------
\newcommand{\purpose}[1]{\par\smallskip\noindent\textit{What is done here and why.} \textit{#1}\par\smallskip}
\newtcolorbox{summarybox}[1][Summary of this section]{breakable,enhanced,colback=black!3,colframe=black!55,
  boxrule=0.4pt,arc=1pt,left=2mm,right=2mm,top=1mm,bottom=1mm,before skip=10pt,after skip=10pt,
  title={#1},fonttitle=\bfseries\small,coltitle=black,colbacktitle=black!12}

% ---------------------------------------------------------------------
% Notation (identical to arXiv:2607.18373 and to the calculation files)
% ---------------------------------------------------------------------
\def\mf#1{\mathbb{#1}}          % coefficients of Omega: \mf{N},\mf{A},\mf{B},\mf{C}
\def\v#1{\mathbf{#1}}
\newcommand{\Kv}{\mathcal{K}}    % Weizsaecker-Williams field  K^i(x)= x^i/x^2
\newcommand{\MS}{\overline{\rm MS}}
\newcommand{\Tr}{{\rm Tr}}
\newcommand{\NLOfirst}{\mathrm{NLO}}
\newcommand{\Lam}{\Lambda}
\newcommand{\NOTE}[1]{{\color{red}\textbf{NOTE:} #1}}
\newcommand{\PV}{\mathrm{P}}
% upper cut-off of the soft modes is denoted \vee, as in 2607.18373 and 1610.03453 (\vee = \Lambda e^{\delta Y})

\title{Single inclusive gluon production at mid rapidity at next-to-leading order in the Color Glass Condensate}

\author{AUTHORS}
\affiliation{AFFILIATIONS}
% \emailAdd{}

\abstract{We compute the single inclusive gluon spectrum in the collision of a dilute projectile with a dense target at next-to-leading order in the Color Glass Condensate framework, for a gluon produced at mid rapidity, in pure Yang--Mills theory. The calculation is performed in light-cone perturbation theory: the projectile is described by its soft gluon light-cone wave function, generated from the valence colour charge density $\rho$ by a unitary dressing operator $\Omega$, and the target by eikonal adjoint Wilson lines. The coefficient of $\Omega$ that creates one gluon at order $g^3$ is computed here; the others were obtained previously. The valence charges are treated as non-commuting operators throughout. We organise the $O(g^4)$ cross section into interference terms with one gluon in the final state and into contributions with two gluons in the final state, one of which is integrated over. All terms enhanced by the logarithm of the longitudinal cut-off ratio have the structure required by high-energy factorisation: the logarithm splits at the longitudinal momentum $k^+$ of the produced gluon, and the coefficient of $\log(\vee/k^+)$ is exactly the leading-order BFKL evolution of the projectile colour charge correlator, while the coefficient of $\log(k^+/\Lambda)$ is exactly the JIMWLK evolution of the target Wilson-line operator of the leading-order cross section. The terms quartic and (in the symmetrised sense) cubic in $\rho$ that accompany these logarithms cancel. The remaining contributions are a collinear term, which is the final-state DGLAP fragmentation of the produced gluon with the $\overline{\rm MS}$ splitting function, a running-coupling term, which sets the scale of the leading-order coupling, and finite terms. We account for all poles of dimensional regularisation. We give the finite terms row by row, after the cancellations that are exact, with every product of three charges written as its symmetrised product plus two-charge contact terms, for an arbitrary projectile and, averaged over the projectile, in colour-singlet form; the colour structures are written through differences of Wilson lines and the longitudinal integrals with plus prescriptions in a momentum fraction. The large-distance tails of individual terms cancel in pairs of virtual and real terms, without any condition on the target at large distance. The terms cubic in $\rho$ do not vanish after the average over the projectile: their symmetrised part, odd in the transverse momentum, is proportional to the symmetrised ($d^{abc}$) three-point function of the projectile, for which no model is assumed, and to C-odd target structures, while the commutators of the charges leave contact terms of the same order as the quadratic terms. Energy denominators that vanish inside the integration region are defined as principal values. In the dilute limit the real-emission part reproduces exactly the two-gluon production cross section in quasi-multi-Regge kinematics. The points that remain open are listed.
}

\begin{document}
\maketitle
\flushbottom

% ======================================================================
\section{Introduction}\label{sec:intro}
% ======================================================================
\NOTE{Section 1 (Introduction): to be written.}

% ======================================================================
\section{Background and the setup}\label{sec:setup}
% ======================================================================

In this section we fix the conventions, introduce the light-cone Hamiltonian of pure Yang--Mills theory and its eikonal form in the presence of fast (valence) modes, and review the light-cone perturbation theory that is used to construct the soft gluon wave function of the projectile. At the end of the section the unitary operator $\Omega$, which generates the soft gluon wave function from the valence state, is defined and parametrised by a set of coefficient functions. These coefficients are computed in Section~\ref{sec:omega}.

\subsection{Light-cone coordinates and conventions}\label{sec:conventions}

We review the basic setup of QCD on the light cone~\cite{Lepage:1980fj,Brodsky:1997de}. In light-cone coordinates a four-vector $x^\mu$ has the components
\begin{align}
x^\mu=(x^+,x^-,\v x)\,,\qquad x^\pm=\frac{1}{\sqrt2}\,(x^0\pm x^3)\,,\qquad \v x=(x^1,x^2)\,.
\end{align}
Quantisation is carried out on the surfaces of fixed $x^+$. Hence $x^+$ plays the role of the light-cone time, along which states evolve, and $x^-$ is the longitudinal coordinate. Momenta are written in the same way,
\begin{align}
p^\mu=(p^+,p^-,\v p)\,,\qquad p^\pm=\frac{1}{\sqrt2}\,(p^0\pm p^3)\,,\qquad p\cdot x=p^+x^-+p^-x^+-\v p\cdot\v x\,.
\end{align}
The projectile moves in the $+$ direction. The longitudinal momentum $p^+$ is positive for a physical particle, and $p^-$ is the light-cone energy, the generator of translations in $x^+$. For an on-shell massless particle $p^2=2p^+p^--\v p^2=0$, so that
\begin{align}\label{eq:energy}
p^-=\frac{\v p^2}{2p^+}\,.
\end{align}
All energy denominators of light-cone perturbation theory are differences of such light-cone energies. We use the abbreviations
\begin{align}\label{eq:measure}
&d^3p\equiv dp^+\,d^2p\,,\qquad \delta^{(3)}(k-p)\equiv\delta(k^+-p^+)\,\delta^{(2)}(\v k-\v p)\,,\qquad \int_x\equiv\int d^2x\,,\notag\\
&E_g(p)\equiv\frac{\v p^2}{2p^+}\,,\qquad E_{gg}(p,q)\equiv\frac{\v p^2}{2p^+}+\frac{\v q^2}{2q^+}\,,
\end{align}
where a momentum argument without boldface, as in $E_g(p)$ or $a(p)$, stands for the pair $(p^+,\v p)$. Where no confusion is possible, transverse vectors are also written without boldface, and $\v p^2$ always denotes the square of the transverse vector. Transverse indices are $i,j,k,l,m=1,2$; since the transverse metric is Euclidean, they are written as upper or lower indices without distinction, and repeated indices are summed. Colour indices in the adjoint representation are $a,b,c,d,e$, and $f^{abc}$ are the structure constants of $SU(N_c)$. The natural logarithm is denoted by $\log$.

We work in the light-cone gauge $A^+=0$. In this gauge the transverse components $A^a_i$, $i=1,2$, are the only dynamical degrees of freedom; $A^-$ is not dynamical and is expressed through the transverse components by the equations of motion. The gluon creation and annihilation operators satisfy
\begin{align}\label{eq:aadagger}
\big[a^a_i(k),a^{\dagger b}_j(p)\big]=(2\pi)^3\,\delta^{ab}\,\delta_{ij}\,\delta^{(3)}(k-p)\,,
\end{align}
with all other commutators vanishing. For the Fourier transforms in the transverse plane we use
\begin{align}\label{eq:FTconventions}
&a(k^+,\v k)=\int d^2z\,e^{-i\v k\cdot\v z}\,a(k^+,\v z)\,,\qquad
a^\dagger(k^+,\v k)=\int d^2z\,e^{i\v k\cdot\v z}\,a^\dagger(k^+,\v z)\,,\notag\\
&f(\v k)=\int d^2x\,e^{i\v k\cdot\v x}\,f(\v x)\,,\qquad f(\v x)=\int\frac{d^2k}{(2\pi)^2}\,e^{-i\v k\cdot\v x}\,f(\v k)\,,
\end{align}
where $f$ stands for any function or operator other than $a$ and $a^\dagger$, in particular for the colour charge density and for the coefficient functions introduced below. With these conventions
\begin{align}\label{eq:aadaggerX}
\big[a^a_i(k^+,\v z),a^{\dagger b}_j(p^+,\v z')\big]=2\pi\,\delta^{ab}\delta_{ij}\,\delta(k^+-p^+)\,\delta^{(2)}(\v z-\v z')\,.
\end{align}
We shall frequently need the Weizs\"acker--Williams field of a point charge and its Fourier transform,
\begin{align}\label{eq:kernels}
\Kv^i(\v x)\equiv\frac{x^i}{\v x^2}\,,\qquad \int_x e^{i\v k\cdot\v x}\,\Kv^i(\v x)=2\pi i\,\frac{k^i}{\v k^2}\,,\qquad
K(\v x,\v y,\v z)\equiv\Kv(\v x-\v z)\cdot\Kv(\v y-\v z)\,.
\end{align}

\subsection{Light-cone QCD Hamiltonian}\label{sec:hamiltonian}

The derivation of the QCD Hamiltonian in the light-cone gauge is standard~\cite{Brodsky:1997de,Lublinsky:2016meo}. We consider pure Yang--Mills theory, $N_f=0$. This is sufficient for the soft gluon sector at high energy; quarks can be added along the lines of~\cite{Lublinsky:2016meo}. After the elimination of $A^-$ the Hamiltonian is a functional of the transverse fields only and is the sum of an electric and a magnetic part,
\begin{align}\label{eq:HLC}
H=\int dx^-\,d^2x\,\Big[\frac12\,\Pi^a\,\Pi^a+\frac14\,F^a_{ij}\,F^a_{ij}\Big]\,,
\end{align}
with
\begin{align}\label{eq:PiF}
\Pi^a\equiv-\frac{1}{\partial^+}\big(D^{ab}_i\,\partial^+A^b_i\big)\,,\qquad
F^a_{ij}\equiv\partial_iA^a_j-\partial_jA^a_i-gf^{abc}A^b_iA^c_j\,,\qquad
D^{ab}_i\equiv\delta^{ab}\partial_i-gf^{abc}A^c_i\,.
\end{align}
Here $\partial^+\equiv\partial/\partial x^-$, $\partial_i\equiv\partial/\partial x^i$ and $g$ is the coupling. The expansion of Eq.~\eqref{eq:HLC} in powers of $g$, $H=H_0+H_{\rm int}$, and the quantisation of the transverse field are given in Appendix~\ref{app:hamiltonian}. The eigenstates of the free Hamiltonian $H_0$ are the Fock states of gluons with the energies~\eqref{eq:energy}; $H_{\rm int}$ contains the three- and four-gluon vertices and the instantaneous interactions characteristic of light-cone quantisation, Eq.~\eqref{eq:Hint}.

\subsection{Born--Oppenheimer formalism and the eikonal approximation}\label{sec:BO}

\paragraph{Valence and soft modes.} In the CGC description of a fast projectile~\cite{McLerran:1993ni,McLerran:1993ka,Jalilian-Marian:1997jx,Kovner:2005en} the gluon modes are divided according to their longitudinal momentum. The \emph{valence} modes have $k^+>\vee$, the \emph{soft} modes have
\begin{align}\label{eq:softwindow}
\Lambda<k^+<\vee\,,\qquad \vee\equiv\Lambda\,e^{\delta Y}\,.
\end{align}
Thus $\vee$ is the upper and $\Lambda$ the lower cut-off on the longitudinal momentum of the soft modes, and $\delta Y=\log(\vee/\Lambda)$ is the rapidity interval that they occupy; the notation is that of~\cite{Lublinsky:2016meo}. Modes with $k^+<\Lambda$ are not part of the projectile wave function considered here. The position of the boundary $\vee$ between valence and soft modes is arbitrary: a change of $\vee$ moves modes from one class to the other, and observables do not depend on it once the valence state is evolved accordingly. Accordingly the field is decomposed as
\begin{align}\label{eq:Asplit}
A^a_i(x)=\underline A^a_i(x)+\overline A^a_i(x)\,,
\end{align}
where the soft field $\underline A$ is given by Eq.~\eqref{eq:modeexp} with the $k^+$ integral restricted to $\Lambda<k^+<\vee$, and the valence field $\overline A$ by the same expression with $k^+>\vee$.

\paragraph{Born--Oppenheimer approximation.} The valence modes have much larger longitudinal momenta than the soft ones and, by Eq.~\eqref{eq:energy}, much smaller light-cone energies at a given transverse momentum. On the time scale of the soft modes they are frozen. The wave function of the soft modes is therefore found at a fixed configuration of the valence modes, which act as a static source~\cite{Kovner:2005en,Lublinsky:2016meo}. The only property of the valence modes that the soft modes are sensitive to, in the approximation described below, is their colour charge density integrated over $x^-$,
\begin{align}\label{eq:rhodef}
\rho^a(-\v p)=-if^{abc}\int_\vee^\infty\frac{dk^+}{2\pi}\int\frac{d^2k}{(2\pi)^2}\,a^{\dagger b}_j(k^+,\v k)\,a^c_j(k^+,\v k+\v p)\,.
\end{align}
It is an operator on the Hilbert space of the valence modes and commutes with the soft operators. It is Hermitian in coordinate space, so that $[\rho^a(\v p)]^\dagger=\rho^a(-\v p)$. From Eqs.~\eqref{eq:rhodef} and~\eqref{eq:aadagger} it follows that the charges at different colours do not commute but form a local $SU(N_c)$ algebra,
\begin{align}\label{eq:rhoalgebra}
\big[\rho^a(\v k),\rho^b(\v p)\big]=if^{abc}\,\rho^c(\v k+\v p)\,,\qquad
\big[\rho^a(\v x),\rho^b(\v y)\big]=if^{abc}\,\rho^c(\v x)\,\delta^{(2)}(\v x-\v y)\,.
\end{align}
We keep the charges as non-commuting operators throughout the paper. Products of charges therefore appear in a definite order, which has to be kept, and the commutator of two charges is one charge: a term with $n$ charges can be traded for a term with $n-1$ charges at the same order in the coupling.

\paragraph{Eikonal approximation.} When the decomposition~\eqref{eq:Asplit} is inserted into $H_{\rm int}$, terms with soft fields only, with valence fields only, and mixed terms are generated. We keep the leading power in the ratio of soft to valence longitudinal momenta, $k^+_{\rm soft}/k^+_{\rm valence}$. In practice this means that every term in which an inverse derivative $1/\partial^+$ acts on a product containing a valence field together with a soft field is dropped, and that the valence fields enter only through the bilinear $\overline A^b_j\,\partial^+\overline A^c_j$, whose integral over $x^-$ is the charge density~\eqref{eq:rhodef}. The terms that contain valence fields only do not act on the soft modes and are omitted. The interaction then is
\begin{align}\label{eq:Hdecomp}
H_{\rm int}\to H_g+H_{ggg}+H_{gg\text{-inst}}+H_{gggg}\,,
\end{align}
as in~\cite{Lublinsky:2016meo}; the four terms are given in Eqs.~\eqref{eq:Hpieces}--\eqref{eq:Hinst} of Appendix~\ref{app:hamiltonian}. $H_{ggg}$ is the three-gluon vertex of the soft modes, $H_{gg\text{-inst}}$ the instantaneous interaction of two soft gluons with the valence charge, and $H_{gggg}$ the four-gluon interaction of the soft modes, which is not needed in this paper. The term that drives the soft gluon wave function is the coupling of one soft gluon to the valence charge,
\begin{align}\label{eq:Hg}
H_g=\int_\Lambda^\vee\frac{dk^+}{2\pi}\int\frac{d^2k}{(2\pi)^2}\,\frac{g\,k^i}{\sqrt2\,(k^+)^{3/2}}\Big[a^{\dagger a}_i(k^+,\v k)\,\rho^a(-\v k)+a^a_i(k^+,\v k)\,\rho^a(\v k)\Big]\,.
\end{align}

\paragraph{Range of validity.} The expressions above are the leading terms of an expansion in $k^+_{\rm soft}/k^+_{\rm valence}$; power corrections in this ratio (non-eikonal corrections) are neglected. The charge density is not assumed to be small in this subsection. In the calculation of the wave function below we expand in powers of $g$ at fixed $\rho$, i.e.\ we consider a dilute projectile, $g\rho\ll1$; the order of a term in $g$ and the number of charges that it contains are then two separate labels.

\subsection{Light-cone quantization and Fock space expansion}\label{sec:fock}

\paragraph{States.} The soft vacuum $\ket0$ is annihilated by all soft annihilation operators. It is a state of the soft modes only: the state of the projectile is a vector in $\mathcal H_{\rm soft}\otimes\mathcal H_{\rm val}$, and $\ket0$ stands for $\ket0\otimes\ket v$ with $\ket v$ the unspecified state of the valence modes. The states with one and two soft gluons are $\ket{g^a_i(k)}=(2\pi)^{-3/2}a^{\dagger a}_i(k)\ket0$ and $\ket{g^a_i(k)\,g^b_j(p)}=(2\pi)^{-3}a^{\dagger a}_i(k)a^{\dagger b}_j(p)\ket0$, normalised to $\delta$-functions without factors of $2\pi$; in transverse coordinate space $\ket{g^a_i(k^+,\v z)}=(2\pi)^{-1/2}a^{\dagger a}_i(k^+,\v z)\ket0$. The complete definitions are collected in Appendix~\ref{app:states}.

\paragraph{Fock space expansion.} A state of the projectile is expanded in the Fock basis of the soft modes,
\begin{align}\label{eq:fock}
\ket\Psi=\psi_0\ket0+\int d^3k\ \psi^a_i(k)\ket{g^a_i(k)}+\frac12\int d^3k\,d^3p\ \psi^{ab}_{ij}(k,p)\ket{g^a_i(k)\,g^b_j(p)}+\dots
\end{align}
The expansion is in the soft modes only. The valence state is not expanded, and the coefficients $\psi_0,\psi^a_i,\dots$ are operators on $\mathcal H_{\rm val}$: they are functionals of the charge density $\rho$ that act on $\ket v$. In a product of such coefficients the order of the charges is kept, as required by Eq.~\eqref{eq:rhoalgebra}. A component of the wave function is characterised by two independent labels, the number $n$ of soft gluons and the number $m$ of charges that multiply it. At a given order in $g$ a component with a given $n$ receives contributions with different $m$; for instance the two-gluon component at order $g^2$ contains a term with one charge (a gluon emitted by the valence charge splits) and a term with two charges (two gluons are emitted). We label the components accordingly, e.g.\ $\ket{\Psi_{gg\rho}}$ and $\ket{\Psi_{gg\rho\rho}}$. Since by Eq.~\eqref{eq:rhoalgebra} a commutator of two charges is one charge, the label $m$ is defined only after a convention for the products of charges has been chosen. We use the decomposition
\begin{align}\label{eq:decomp}
\rho^a(\v k)\,\rho^b(\v p)=\frac12\big\{\rho^a(\v k),\rho^b(\v p)\big\}+\frac{i}{2}\,f^{abc}\,\rho^c(\v k+\v p)\,,
\end{align}
and count the anticommutator as a term with two charges and the commutator as a term with one charge. The anticommutator describes two emitters at different transverse positions; the commutator is local in coordinate space and is equivalent to one emitter. We refer to Eq.~\eqref{eq:decomp} as the decomposition procedure.

\subsection{Light-cone perturbation theory}\label{sec:lcpt}

Light-cone perturbation theory (LCPT) is time-independent (old-fashioned) perturbation theory for the Hamiltonian $H=H_0+V$, $V\equiv H_{\rm int}$, in the basis of the eigenstates of $H_0$~\cite{Lepage:1980fj,Brodsky:1997de}. The interaction acts on a free Fock state and generates the other Fock components; each insertion of $V$ is accompanied by an energy denominator, the difference between the energy of the initial state and the energy of the intermediate state that follows the insertion. In the present problem there is one difference to the textbook case. The matrix elements of $V$ between soft Fock states,
\begin{align}
V_{nm}\equiv\braket{n|H_{\rm int}|m}\,,
\end{align}
are not numbers but operators on the valence Hilbert space, since $H_g$ and $H_{gg\text{-inst}}$ contain $\rho$. They do not commute with each other. The energies $E_n$, on the other hand, are numbers, because the valence modes do not contribute to the energy in the approximation of Section~\ref{sec:BO}. The perturbative series therefore has the same form as for a number-valued perturbation, provided that the order of the matrix elements is kept.

\paragraph{Perturbation theory with non-commuting matrix elements.} Let $\ket i$ be an eigenstate of $H_0$ with energy $E_i$, and assume $V_{ii}=0$, which holds for all cases considered below. The state that develops from $\ket i$ when the interaction is switched on is, to second order,
\begin{align}\label{eq:pt2}
\ket{\Psi_i}=\ &\ket i\,\Big[1-\frac12\sum_{m\neq i}\frac{V_{im}V_{mi}}{(E_i-E_m)^2}\Big]
+\sum_{n\neq i}\ket n\,\frac{V_{ni}}{E_i-E_n}\notag\\
&+\sum_{n,m\neq i}\ket n\,\frac{V_{nm}\,V_{mi}}{(E_i-E_n)(E_i-E_m)}+O(V^3)\,.
\end{align}
The sums run over the free Fock states of the soft modes, and the later interaction stands to the left. The second and third terms are the Rayleigh--Schr\"odinger series; the term in the square bracket normalises the state. The corresponding series to third order for the soft vacuum, Eq.~\eqref{eq:pt3} of Appendix~\ref{app:pt}, which underlies the results quoted in Section~\ref{sec:vacin}, is the one of Appendix~E of~\cite{Lublinsky:2016meo}. The states are defined uniquely by two properties: $H$ does not connect dressed states of different energy, and the dressed states are orthonormal with a Hermitian component along the initial state (canonical normalisation). The second property fixes the order of the factors in the normalisation terms. We have verified the series, including the order of the factors, against an exact numerical construction (Appendix~\ref{app:pt}).

\paragraph{Energy denominators.} For the dressing of the soft vacuum the denominators are $E_0-E_n=-E_n<0$ and never vanish. For the dressing of a one-gluon state $\ket{g^a_i(k)}$ the denominators are $E_g(k)-E_n$. If the state $\ket n$ contains one gluon of momentum $p$, the denominator $E_g(k)-E_g(p)$ vanishes inside the integration region, on the surface $p^+\v k^2=k^+\v p^2$; if $\ket n$ contains two gluons that originate from a splitting of the gluon $k$, it is negative definite. We define the vanishing denominators as principal values,
\begin{align}\label{eq:PV}
\frac{1}{E_g(k)-E_g(p)}\to\PV\,\frac{1}{E_g(k)-E_g(p)}\,.
\end{align}
Eq.~\eqref{eq:PV} is part of the definition of the dressed states used in this paper, and all results below refer to it. It is the Hermitian part of the prescription $1/(E_i-E_n+i\varepsilon)$. The adiabatic prescription $E_i-E_n\to E_i-E_n+i\varepsilon$, which is used in~\cite{Radhakrishnan:2026soft}, differs from the principal value by terms proportional to $\delta(E_i-E_n)$, which have support on the energy shell; these terms are not included here. Vanishing denominators occur in one coefficient only, $\mf B_3$ of Section~\ref{sec:oneout}, and therefore in one row of the cross section (Section~\ref{sec:nuance_pv}).

\paragraph{Kinematic restrictions.} Every soft gluon, in the final state or in an intermediate state, has $\Lambda<p^+<\vee$. In addition, the longitudinal momentum that the last term of $H_{gg\text{-inst}}$ transfers between a soft gluon and the valence modes is bounded from below by $\Lambda$, $|p^+-q^+|>\Lambda$ in Eq.~\eqref{eq:Hinst}; the same restriction applies to the longitudinal momentum of a gluon that is emitted or absorbed by the valence charge through $H_g$, since this gluon is itself a soft mode. These restrictions fix the limits of all longitudinal integrals.

\paragraph{Matrix elements.} The matrix elements of $H_g$, $H_{ggg}$ and $H_{gg\text{-inst}}$ between states with up to two gluons, Eqs.~\eqref{eq:m1}--\eqref{eq:m5}, are listed in Appendix~\ref{app:matrixelements}. They coincide with those of~\cite{Lublinsky:2016meo}. All of them except the three-gluon vertex are operators on $\mathcal H_{\rm val}$.

\subsection{Soft gluon wave function in the CGC and the evolution operator \texorpdfstring{$\Omega$}{Omega}}\label{sec:omega_def}

\paragraph{Definition.} The states~\eqref{eq:pt2} are the dressed states of the soft modes in the presence of the valence charge. They are obtained from the free Fock states by an operator $\Omega$,
\begin{align}\label{eq:Omegadef}
\ket{\Psi_n}=\Omega\ket n\,,\qquad \Omega^\dagger\Omega=\Omega\,\Omega^\dagger=1\,,
\end{align}
which acts on $\mathcal H_{\rm soft}\otimes\mathcal H_{\rm val}$ and is unitary because it maps one orthonormal basis onto another. By the first of the two properties stated in Section~\ref{sec:lcpt}, $\Omega^\dagger H\Omega$ has no matrix elements between free soft Fock states of different energy: $\Omega$ diagonalises the Hamiltonian of the soft modes order by order in $g$~\cite{Kovner:2005en,Radhakrishnan:2026soft}. In particular
\begin{align}
\ket{\Psi_0}=\Omega\ket0
\end{align}
is the soft gluon wave function of the projectile: the cloud of soft gluons that accompanies the valence charges. When the rapidity interval of the soft modes is increased, $\Omega$ generates the gluons that populate the new interval, which is why it is called the evolution operator.

At the lowest order in $g$ the only interaction is $H_g$, and $\Omega=1+\int\frac{d^3p}{(2\pi)^3}\big[\mf A^b_j(p)\,a^{\dagger b}_j(p)-\mf A^{\dagger b}_j(p)\,a^b_j(p)\big]+O(g^2)$ with $\mf A^b_j(p)\propto g\,p^j\rho^b(-\v p)/\v p^2$ the Weizs\"acker--Williams field of the valence charge, Eq.~\eqref{eq:A1mom} below. For commuting charges this is the expansion of a coherent operator, which dresses the valence state with independently emitted gluons~\cite{Kovner:2005en}. Beyond the lowest order $\Omega$ contains the non-Abelian vertices and the commutators of the charges, and it is not a coherent operator.

\paragraph{Normal-ordered expansion.} Since any operator on the soft Fock space can be written as a sum of normal-ordered products of $a^\dagger$ and $a$, we parametrise
\begin{align}\label{eq:OmegaExp}
\Omega = \ &\mf{N}
+\int\frac{d^3p}{(2\pi)^3}\Big[\mf{A}^b_{1j}(p)\,a^b_j(p)+\mf{A}^b_{2j}(p)\,a^{\dagger b}_j(p)\Big]
\notag\\&
+\int\frac{d^3p}{(2\pi)^3}\frac{d^3q}{(2\pi)^3}\Big[\mf{B}^{cb}_{1kj}(p,q)\,a^b_j(p)a^c_k(q)+\mf{B}^{cb}_{2kj}(p,q)\,a^{\dagger b}_j(p)a^{\dagger c}_k(q)
\notag\\&\hspace{3.4cm}+\mf{B}^{cb}_{3kj}(p,q)\,a^{\dagger c}_k(q)a^b_j(p)\Big]
\notag\\&
+\int\frac{d^3p}{(2\pi)^3}\frac{d^3q}{(2\pi)^3}\Big[\mf{C}^{dcb}_{1lkj}(p,q)\,a^{\dagger d}_l(q)a^{\dagger c}_k(p-q)a^b_j(p)
\notag\\&\hspace{3.4cm}+\mf{C}^{dcb}_{2lkj}(p,q)\,a^{\dagger b}_j(p)a^c_k(p-q)a^d_l(q)\Big]+\dots
\end{align}
The coefficients $\mf N$, $\mf A_{1,2}$, $\mf B_{1,2,3}$, $\mf C_{1,2}$ are operators on $\mathcal H_{\rm val}$, i.e.\ functionals of $\rho$, and do not commute with each other; they commute with $a$ and $a^\dagger$. In each term the first index pair and the first momentum argument are tied to each other as indicated; note that in $\mf B^{cb}_{kj}(p,q)$ the \emph{second} colour and polarisation indices $(b,j)$ belong to the \emph{first} argument $p$. In the terms with three operators we have used that all vertices conserve the longitudinal and the transverse momentum of the soft gluons when $\rho$ is not involved; this is sufficient here because $\mf C_{1,2}$ are needed only at order $g$, where they are independent of $\rho$. The dots stand for terms with three operators that depend on $\rho$ and for terms with four or more operators. They are of order $g^2$ or higher and do not contribute to the single inclusive gluon cross section at the order considered in this paper (Section~\ref{sec:threeout}).

The products of operators in Eq.~\eqref{eq:OmegaExp} are symmetric under the exchange of two creation or of two annihilation operators, so that only the symmetric parts of the coefficients are defined. We impose
\begin{align}\label{eq:symm}
\mf B^{cb}_{1kj}(p,q)=\mf B^{bc}_{1jk}(q,p)\,,\qquad \mf B^{cb}_{2kj}(p,q)=\mf B^{bc}_{2jk}(q,p)\,,\qquad
\mf C^{dcb}_{1lkj}(p,q)=\mf C^{cdb}_{1klj}(p,p-q)\,,
\end{align}
and the same relation for $\mf C_2$.

\paragraph{Orders in the coupling.} In the expansion in $g$ at fixed $\rho$ the coefficients start at the following orders,
\begin{align}\label{eq:orders}
\mf N=1+\mf N_2\,,\quad \mf N_2=O(g^2)\,;\qquad
\mf A_{1,2}=\mf A^{(1)}_{1,2}+\mf A^{(3)}_{1,2}\,;\qquad
\mf B_{1,2,3}=O(g^2)\,;\qquad \mf C_{1,2}=O(g)\,,
\end{align}
where $\mf A^{(n)}=O(g^n)$. There are no terms of order $g^2$ in $\mf A_{1,2}$ and no terms of order $g$ in $\mf N$, because the interaction changes the number of gluons by one at order $g$ and the states are normalised canonically. When no confusion can arise we write $\mf A\equiv\mf A^{(1)}_2$ and $\mf C\equiv\mf C_1$.

\paragraph{Action on the vacuum and on a one-gluon state.} Acting with Eq.~\eqref{eq:OmegaExp} on the soft vacuum, only the terms without annihilation operators survive. With the normalisation~\eqref{eq:states},
\begin{align}\label{eq:Omegavac}
\Omega\ket0=\ &\mf N\ket0+(2\pi)^{3/2}\int\frac{d^3p}{(2\pi)^3}\,\mf A^b_{2j}(p)\ket{g^b_j(p)}\notag\\&+(2\pi)^3\int\frac{d^3p}{(2\pi)^3}\frac{d^3q}{(2\pi)^3}\,\mf B^{cb}_{2kj}(p,q)\ket{g^b_j(p)\,g^c_k(q)}+\dots
\end{align}
Acting on a one-gluon state, each annihilation operator has to be contracted with the gluon of the state, by Eq.~\eqref{eq:aadagger}, and one obtains
\begin{align}\label{eq:Omegag}
\Omega\ket{g^a_i(k)}=\ &\frac{1}{(2\pi)^{3/2}}\,\mf A^a_{1i}(k)\ket0+\mf N\ket{g^a_i(k)}+\int\frac{d^3p}{(2\pi)^3}\,\mf B^{ba}_{3ji}(k,p)\ket{g^b_j(p)}
\notag\\&
+(2\pi)^{3/2}\int\frac{d^3p}{(2\pi)^3}\,\mf A^b_{2j}(p)\ket{g^a_i(k)\,g^b_j(p)}
\notag\\&+(2\pi)^{3/2}\int\frac{d^3p}{(2\pi)^3}\,\mf C^{bca}_{1jki}(k,p)\ket{g^b_j(p)\,g^c_k(k-p)}+\dots,
\end{align}
where the dots are states with three or more gluons and terms of higher order. The terms have a simple interpretation. $\mf A_1$ absorbs the gluon; $\mf N$ and $\mf B_3$ leave the number of gluons unchanged, $\mf B_3$ describing the scattering of the gluon on the valence charges; $\mf A_2$ adds a gluon emitted by the valence charges to the incoming one, which is a spectator; $\mf C_1$ splits the gluon into two.

\paragraph{How the coefficients are determined.} The coefficients are fixed by two requirements:
\begin{itemize}
\item \emph{Matching.} $\Omega\ket n$ equals the state $\ket{\Psi_n}$ computed in LCPT, Eqs.~\eqref{eq:pt2} and~\eqref{eq:pt3}. Comparing Eq.~\eqref{eq:Omegavac} with $\ket{\Psi_0}$, Fock sector by Fock sector, fixes $\mf N$, $\mf A_2$ and $\mf B_2$. Comparing Eq.~\eqref{eq:Omegag} with the dressed one-gluon state fixes $\mf A_1$, $\mf B_3$ and $\mf C_1$, and reproduces $\mf N$ and $\mf A_2$.
\item \emph{Unitarity.} $\Omega^\dagger\Omega=1$, imposed order by order in $g$, relates the coefficients to each other. It determines the coefficients that multiply annihilation operators only, $\mf A_1$ and $\mf B_1$, and $\mf C_2$, in terms of those fixed by the matching, and it gives relations among the latter which serve as checks.
\end{itemize}
Each coefficient is probed by the transition between the Fock sectors that correspond to its annihilation and creation operators. It is therefore sufficient to consider the vacuum and the one-gluon state as initial states: all coefficients of Eq.~\eqref{eq:OmegaExp} have at most one annihilation operator, except $\mf B_1$ and $\mf C_2$, which follow from unitarity. The vacuum initial state was treated in~\cite{Lublinsky:2016meo}, and we quote the results in Section~\ref{sec:vacin}. The one-gluon initial state is treated in Section~\ref{sec:onegin}, and the unitarity relations are derived in Section~\ref{sec:unitarity}.

% ======================================================================
\section{Perturbative expansion of the evolution operator \texorpdfstring{$\Omega$}{Omega}}\label{sec:omega}
% ======================================================================

In this section we determine the coefficients of the expansion~\eqref{eq:OmegaExp} to the order needed for the single inclusive gluon cross section at next-to-leading order. We first collect the results for the soft vacuum as initial state, which are known~\cite{Lublinsky:2016meo}. We then compute the dressed one-gluon state in LCPT, component by component, and read off the coefficients $\mf A_1$, $\mf B_3$ and $\mf C_1$. Finally we impose the unitarity of $\Omega$, derive the relations among the coefficients order by order in $g$, and summarise the coefficients in momentum and in coordinate space.

In the diagrams of this section the straight horizontal line represents the valence charges, curly lines are soft gluons, and a vertical dashed line marks an intermediate state whose energy enters a denominator. A component of the dressed one-gluon state is denoted by $\ket{(\Psi^a_i(k))_{\dots}}$, where $(k,a,i)$ are the momentum, colour and polarisation of the incoming gluon and the subscript lists the outgoing gluons and the charges, as explained below Eq.~\eqref{eq:fock}. Outgoing and intermediate momenta are integrated over.

\subsection{The vacuum incoming state}\label{sec:vacin}

The dressing of the soft vacuum, $\ket{\Psi_0}=\Omega\ket0$, was computed to order $g^3$ in~\cite{Lublinsky:2016meo}, including quarks. We do not repeat the derivation and quote the results that are needed, restricted to pure Yang--Mills theory and written in the conventions of Section~\ref{sec:setup}. They follow from the third-order series~\eqref{eq:pt3} (Appendix~\ref{app:pt}) with the matrix elements~\eqref{eq:m1}--\eqref{eq:m5} (Appendix~\ref{app:matrixelements}); the coefficients are read off by comparison with Eq.~\eqref{eq:Omegavac}.

\paragraph{Order $g$.} The valence charge emits one soft gluon,
\begin{align}\label{eq:PsiLO}
\ket{\Psi^{\rm LO}_{g\rho}}=-\int_\Lambda^\vee dk^+\!\int d^2k\,\ket{g^a_i(k)}\frac{\braket{g^a_i(k)|H_g|0}}{E_g(k)}
=-\int_\Lambda^\vee dk^+\!\int d^2k\,\frac{g\,k^i\rho^a(-\v k)}{2\pi^{3/2}\sqrt{k^+}\,\v k^2}\ket{g^a_i(k)}\,,
\end{align}
so that, by Eq.~\eqref{eq:Omegavac},
\begin{align}\label{eq:A1mom}
\mf A^{(1)a}_{2i}(k)=-\frac{\sqrt2\,g\,k^i\rho^a(-\v k)}{\sqrt{k^+}\,\v k^2}\,,\qquad \Lambda<k^+<\vee\,.
\end{align}
This is the Weizs\"acker--Williams field of the valence charge. The coefficient of the vacuum component at order $g^2$ is the normalisation term of Eq.~\eqref{eq:pt3},
\begin{align}\label{eq:N2mom}
\mf N_2=-\frac12\braket{\Psi^{\rm LO}_{g\rho}|\Psi^{\rm LO}_{g\rho}}=-\frac{g^2}{8\pi^3}\,\log\frac{\vee}{\Lambda}\int d^2p\,\frac{\rho^a(-\v p)\,\rho^a(\v p)}{\v p^2}\,.
\end{align}
It is Hermitian, and it does not commute with a single charge.

\paragraph{Order $g^2$: two gluons.} The two-gluon component arises from the emission of two gluons by the valence charge, from the emission of one gluon that splits, and from the instantaneous emission of a pair. The product of two charges in the first contribution is decomposed with Eq.~\eqref{eq:decomp}. The resulting coefficient $\mf B_2$~\cite{Lublinsky:2016meo} is given in Eq.~\eqref{eq:B2mom} of Appendix~\ref{app:coefficients}. It is the sum of a part linear in $\rho$ (commutator of the two emissions, splitting, instantaneous emission), with support $k^++p^+<\vee$ because the parent gluon is a soft mode, and of the anticommutator part $\propto\{\rho^b(-\v p),\rho^a(-\v k)\}$, with support $\Lambda<k^+,p^+<\vee$: two independently emitted gluons are each soft modes, with no condition on $k^++p^+$.

\NOTE{The commutator part of the two-emission diagram is written with the support $k^++p^+<\vee$; by its kinematics it could also be the full square $\Lambda<k^+,p^+<\vee$. To be decided by the authors.}

\paragraph{Order $g^3$: one gluon.} The one-gluon component at order $g^3$ defines $\mf A^{(3)}_2$. It contains terms with one, two and three charges,
\begin{align}
\mf A^{(3)}_2=\big[\mf A^{(3)}_2\big]_\rho+\big[\mf A^{(3)}_2\big]_{\rho\rho}+\big[\mf A^{(3)}_2\big]_{\rho\rho\rho}\,,
\end{align}
which originate from the diagrams with one, two and three valence vertices and from the two normalisation terms of Eq.~\eqref{eq:pt3}, after the decomposition~\eqref{eq:decomp}. The transverse loop integral of the self-energy of the emitted gluon is ultraviolet divergent; it is done in $d=2-\epsilon$ transverse dimensions, and $\mu_{\MS}$ is the $\MS$ scale. In the limit $\Lambda\to0$ wherever it can be taken, the part linear in $\rho$ is
\begin{align}\label{eq:A3onerho}
\Big[\mf A^{(3)a}_{2i}(k)\Big]_{\rho}=\ &\frac{g^3N_c\,\rho^a(-\v k)\,k^i}{8\sqrt{2k^+}\,\pi^2\,\v k^2}\Bigg\{\Big[-\frac2\epsilon+\log\frac{\v k^2}{\mu^2_{\MS}}\Big]
\Big[\frac{11}{3}-2\log\frac{k^+}{\Lambda}+\log\Big(1-\frac{k^+}{\vee}\Big)\Big]
\notag\\&\hspace{1.5cm}
+\frac{\pi^2}{3}-\frac{67}{9}+\log^2\Big(1-\frac{k^+}{\vee}\Big)+\log\Big(1-\frac{k^+}{\vee}\Big)\log\frac{\vee}{k^+}\Bigg\}\,.
\end{align}
The part with two charges is proportional to $f^{abc}\{\rho^b(-\v p),\rho^c(-\v k+\v p)\}$ and is given in Eq.~\eqref{eq:A3tworho} of Appendix~\ref{app:coefficients}; the part with three charges is
\begin{align}\label{eq:A3threerho}
\Big[\mf A^{(3)a}_{2i}(k)\Big]_{\rho\rho\rho}=\log\frac{\vee}{\Lambda}\,\frac{\sqrt2\,g^3\,k^i}{8\pi^3\sqrt{k^+}\,\v k^2}\,\rho^a(-\v k)\int d^2p\,\frac{\rho^b(\v p)\,\rho^b(-\v p)}{\v p^2}
=\mf A^{(1)a}_{2i}(k)\,\mf N_2\,.
\end{align}
The terms proportional to $1/\Lambda$ and $1/\vee$ of the individual diagrams cancel in Eq.~\eqref{eq:A3threerho}. The order of the charges in Eq.~\eqref{eq:A3threerho}, with $\rho^a(-\v k)$ to the left of the norm factor, follows from the order of the factors in the last term of Eq.~\eqref{eq:pt3}. It matters, because $\mf N_2$ and $\mf A^{(1)}_2$ do not commute; see Section~\ref{sec:unitarity}.

Eqs.~\eqref{eq:A3onerho}--\eqref{eq:A3threerho} correspond to Eqs.~(3.97)--(3.103) of~\cite{Lublinsky:2016meo} (arXiv version~2), with one difference: in Eq.~(3.102) there the three charges are written in the order $\rho^b\rho^b\rho^a$, whereas the order that follows from Eq.~\eqref{eq:pt3} is that of Eq.~\eqref{eq:A3threerho}.

\subsection{Calculation of the single-gluon incoming state}\label{sec:onegin}

We now compute the state $\ket{\Psi^a_i(k)}=\Omega\ket{g^a_i(k)}$ that develops from a soft gluon of momentum $k=(k^+,\v k)$, colour $a$ and polarisation $i$, with $\Lambda<k^+<\vee$. We use Eq.~\eqref{eq:pt2} with $\ket i=\ket{g^a_i(k)}$ and $E_i=E_g(k)$. The possible final states are the vacuum, one gluon, two gluons and three gluons. For each of them we give the LCPT expression, evaluate it with the matrix elements of Appendix~\ref{app:matrixelements}, and give the result in momentum space.

The long intermediate expressions, the Fourier transforms and the forms of all components in transverse coordinate space are collected in Appendix~\ref{app:onegluon}; the cancellation of the double pole, the algebra of the terms with two valence vertices and the normalisation are in Appendices~\ref{app:doublepole}, \ref{app:tworho} and~\ref{app:norm}.

\subsubsection{Vacuum outgoing state}\label{sec:vacout}

At order $g$ the incoming gluon can be absorbed by the valence charge (Figure~\ref{fig:g_vac}). This is the first-order term of Eq.~\eqref{eq:pt2} with $\ket n=\ket0$ and the denominator $E_g(k)-E_0=E_g(k)$,
\begin{align}\label{eq:psi_rho}
\ket{(\Psi^a_i(k))_\rho}=\ket0\,\frac{\braket{0|H_g|g^a_i(k)}}{E_g(k)}=\frac{g\,k^i\rho^a(\v k)}{2\pi^{3/2}\sqrt{k^+}\,\v k^2}\ket0\,,
\end{align}
where the Hermitian conjugate of Eq.~\eqref{eq:m1} was used. Comparing with the first term of Eq.~\eqref{eq:Omegag} gives
\begin{align}\label{eq:A11mom}
\mf A^{(1)a}_{1i}(k)=\frac{\sqrt2\,g\,k^i\rho^a(\v k)}{\sqrt{k^+}\,\v k^2}=-\big[\mf A^{(1)a}_{2i}(k)\big]^\dagger\,,
\end{align}
where the last equality refers to Eq.~\eqref{eq:A1mom}. The absorption amplitude is minus the conjugate of the emission amplitude, as required by unitarity (Section~\ref{sec:unitarity}). The vacuum component of order $g^3$, $\mf A^{(3)}_1$, will not be computed diagrammatically; it follows from unitarity, Eq.~\eqref{eq:A31unit}.

\begin{figure}[t]
\centering
\tikzfig{g_tovacum}
\caption{Absorption of the incoming gluon by the valence charge, Eq.~\eqref{eq:psi_rho}.}
\label{fig:g_vac}
\end{figure}

\subsubsection{One outgoing gluon}\label{sec:oneout}

At order $g^2$ the incoming gluon can turn into a gluon of different momentum $p=(p^+,\v p)$, colour $b$ and polarisation $j$ by exchanging momentum with the valence charges. The sum of the contributions defines $\mf B_3$ through the third term of Eq.~\eqref{eq:Omegag}. We distinguish the terms with one valence vertex, which are linear in $\rho$, from the terms with two valence vertices. Since the momentum exchanged with the valence modes is not zero, the outgoing gluon is never identical to the incoming one in these terms; the component along $\ket{g^a_i(k)}$ itself is the normalisation, which is discussed in Section~\ref{sec:summarywf}.

\paragraph{Single $\rho$ part: instantaneous interaction.} The instantaneous interaction~\eqref{eq:m5} contributes at first order of Eq.~\eqref{eq:pt2} (Figure~\ref{fig:g_g_one}a),
\begin{align}\label{eq:psi_inst_def}
\ket{(\Psi^a_i(k))_{g\rho\,{\rm inst}}}=\int dp^+\!\int d^2p\ \ket{g^b_j(p)}\,\frac{\braket{g^b_j(p)|H_{gg\text{-inst}}|g^a_i(k)}}{E_g(k)-E_g(p)}\,.
\end{align}
With $E_g(k)-E_g(p)=(p^+\v k^2-k^+\v p^2)/(2k^+p^+)$,
\begin{align}\label{eq:psi_inst}
\ket{(\Psi^a_i(k))_{g\rho\,{\rm inst}}}=\frac{ig^2f^{abc}}{(2\pi)^3}\int dp^+\!\int d^2p\ \frac{\sqrt{p^+k^+}\,(k^++p^+)}{(p^+-k^+)^2}\,\frac{\rho^c(\v k-\v p)}{p^+\v k^2-k^+\v p^2}\,\ket{g^b_i(p)}\,.
\end{align}
The longitudinal momentum $p^+-k^+$ is exchanged with the valence modes. The range of $p^+$ consists of two intervals (Figure~\ref{fig:kinem}),
\begin{align}\label{eq:ranges}
k^++\Lambda<p^+<\vee\quad(\text{region }1)\,,\qquad\qquad \Lambda<p^+<k^+-\Lambda\quad(\text{region }2)\,.
\end{align}
In region~1 the valence modes transfer longitudinal momentum to the gluon, in region~2 the gluon transfers longitudinal momentum to the valence modes. We denote the two contributions by $\ket{(\Psi^a_i(k))^{1}_{g\rho\,{\rm inst}}}$ and $\ket{(\Psi^a_i(k))^{2}_{g\rho\,{\rm inst}}}$; they have the same integrand. The denominator $p^+\v k^2-k^+\v p^2$ vanishes inside the region of integration and is a principal value, Eq.~\eqref{eq:PV}. \begin{figure}[t]
\centering
\tikzfig{kinem}
\caption{Longitudinal momenta allowed in the transition of a soft gluon of momentum $k^+$ into a soft gluon of momentum $p^+$ on the valence charges. One shaded region is the condition $\Lambda<|p^+-k^+|<\vee$ on the exchanged momentum, the other the conditions $\Lambda<k^+<\vee$ and $\Lambda<p^+<\vee$ on the two gluons; the allowed region is their intersection, Eq.~\eqref{eq:ranges}.}
\label{fig:kinem}
\end{figure}

\begin{figure}[t]
\centering
\begin{tabular}{ccc}
\scalebox{0.7}{\tikzfig{gtog_inst}}&\scalebox{0.7}{\tikzfig{gandgtog}}&\scalebox{0.7}{\tikzfig{gto3g_torho}}\\[2mm]
(a)&(b)&(c)
\end{tabular}
\caption{Transition of the incoming gluon into a gluon of different momentum, terms linear in $\rho$: (a) instantaneous interaction, Eq.~\eqref{eq:psi_inst_def}; (b) emission of a gluon by the valence charge followed by the merging of the two gluons, Eq.~\eqref{eq:psi1_def}; (c) splitting of the incoming gluon followed by the absorption of one of the two gluons, Eq.~\eqref{eq:psi2_def}.}
\label{fig:g_g_one}
\end{figure}

\paragraph{Single $\rho$ part: three-gluon vertex.} At second order of Eq.~\eqref{eq:pt2} there are two contributions linear in $\rho$, which contain the three-gluon vertex~\eqref{eq:m3} and one valence vertex. In the first one (Figure~\ref{fig:g_g_one}b) the valence charge emits a gluon of momentum $r=p-k$, and the two gluons merge. The intermediate state contains the gluons $k$ and $r$, so that the first denominator is $E_g(k)-E_{gg}(k,r)=-E_g(r)$:
\begin{align}\label{eq:psi1_def}
\ket{(\Psi^a_i(k))^1_{g\rho}}=\frac12\int d^3p\,d^3q\,d^3r\ \ket{g^b_j(p)}\,
\frac{\braket{g^b_j(p)|H_{ggg}|g^c_l(q)g^d_m(r)}\braket{g^c_l(q)g^d_m(r)|H_g|g^a_i(k)}}{\big(E_g(k)-E_{gg}(q,r)\big)\big(E_g(k)-E_g(p)\big)}\,.
\end{align}
The factor $\frac12$ is the weight of the two-gluon intermediate states (Section~\ref{sec:fock}); the two terms of Eq.~\eqref{eq:m2} give equal contributions and cancel it. The first matrix element is the complex conjugate of Eq.~\eqref{eq:m3}. Since the emitted gluon is a soft mode, $r^+=p^+-k^+>\Lambda$: this contribution exists in region~1 of Eq.~\eqref{eq:ranges} only. Inserting the matrix elements (Appendix~\ref{app:onegluon}, Eq.~\eqref{eq:psi1}) we find, with the momentum fraction $\xi\equiv p^+/k^+$,
\begin{align}\label{eq:psi1_xi}
\ket{(\Psi^a_i(k))^1_{g\rho}}=\ &-\frac{ig^2f^{abc}}{4\pi^3}\int_{1+\Lambda/k^+}^{\vee/k^+}d\xi\int d^2p\ \frac{\sqrt\xi}{1-\xi}\,\frac{(p-k)^m\,(p-\xi k)^l}{(\v p-\v k)^2\,(\v p^2-\xi\v k^2)}\,\rho^c(\v k-\v p)\notag\\
&\times\Big(\delta_{il}\delta_{jm}-\frac1\xi\,\delta_{jl}\delta_{im}-\frac{1}{1-\xi}\,\delta_{ij}\delta_{lm}\Big)\ket{g^b_j(\xi k^+,\v p)}\,.
\end{align}
In the second contribution (Figure~\ref{fig:g_g_one}c) the incoming gluon splits into the gluons $p$ and $r=k-p$, and the gluon $r$ is absorbed by the valence charge. The first denominator is $E_g(k)-E_{gg}(p,r)$, which is negative definite:
\begin{align}\label{eq:psi2_def}
\ket{(\Psi^a_i(k))^2_{g\rho}}=\frac12\int d^3p\,d^3q\,d^3r\ \ket{g^b_j(p)}\,
\frac{\braket{g^b_j(p)|H_g|g^c_l(q)g^d_m(r)}\braket{g^c_l(q)g^d_m(r)|H_{ggg}|g^a_i(k)}}{\big(E_g(k)-E_{gg}(q,r)\big)\big(E_g(k)-E_g(p)\big)}\,.
\end{align}
Here $r^+=k^+-p^+>\Lambda$, i.e.\ region~2. The result is (Appendix~\ref{app:onegluon}, Eq.~\eqref{eq:psi2})
\begin{align}\label{eq:psi2_xi}
\ket{(\Psi^a_i(k))^2_{g\rho}}=\ &\frac{ig^2f^{abc}}{4\pi^3}\int_{\Lambda/k^+}^{1-\Lambda/k^+}d\xi\ \frac{\xi^{3/2}}{1-\xi}\int d^2p\ \frac{(k-p)^m\,(p-\xi k)^l}{(\v p^2-\xi\v k^2)\,(\v p-\xi\v k)^2}\,\rho^c(\v k-\v p)\notag\\
&\times\Big(\delta_{il}\delta_{jm}-\frac1\xi\,\delta_{jl}\delta_{im}-\frac{1}{1-\xi}\,\delta_{lm}\delta_{ij}\Big)\ket{g^b_j(\xi k^+,\v p)}\,.
\end{align}
The overall signs of Eqs.~\eqref{eq:psi1} and~\eqref{eq:psi2} are fixed by the matrix elements~\eqref{eq:m2} and~\eqref{eq:m3}; we have checked them numerically by evaluating Eqs.~\eqref{eq:psi1_def} and~\eqref{eq:psi2_def} directly, and through the unitarity relation~\eqref{eq:B3unit}.

Each of the four terms linear in $\rho$ has a double pole at $p^+=k^+$, which alone would give a contribution proportional to $k^+/\Lambda$. It cancels between the instantaneous and the non-instantaneous term in each of the two regions~\eqref{eq:ranges} (Appendix~\ref{app:doublepole}).

\paragraph{Two $\rho$ part.} At second order of Eq.~\eqref{eq:pt2} there are two contributions with two valence vertices (Figure~\ref{fig:g_g_two}). They differ by the order of the two vertices in light-cone time, and since the charges do not commute, also by the order of the charges. In the first one the incoming gluon is absorbed and a gluon is emitted afterwards; the intermediate state is the soft vacuum, with the denominator $E_g(k)$,
\begin{align}\label{eq:psi1rr}
\ket{(\Psi^a_i(k))^1_{g\rho\rho}}&=\int_\Lambda^\vee dp^+\!\int d^2p\ \ket{g^b_j(p)}\,\frac{\braket{g^b_j(p)|H_g|0}\braket{0|H_g|g^a_i(k)}}{E_g(k)\,\big(E_g(k)-E_g(p)\big)}\notag\\
&=\frac{g^2}{4\pi^3}\int_\Lambda^\vee dp^+\!\int d^2p\ \sqrt{\frac{k^+}{p^+}}\ \frac{k^ip^j\,\rho^b(-\v p)\,\rho^a(\v k)}{\v k^2\,(p^+\v k^2-k^+\v p^2)}\,\ket{g^b_j(p)}\,.
\end{align}
In the second one a gluon is emitted first and the incoming gluon is absorbed afterwards; the intermediate state contains the two gluons $k$ and $p$, with the denominator $E_g(k)-E_{gg}(k,p)=-E_g(p)$,
\begin{align}\label{eq:psi2rr}
\ket{(\Psi^a_i(k))^2_{g\rho\rho}}&=\frac12\int d^3p\,d^3q\,d^3r\ \ket{g^b_j(p)}\,\frac{\braket{g^b_j(p)|H_g|g^c_l(q)g^d_m(r)}\braket{g^c_l(q)g^d_m(r)|H_g|g^a_i(k)}}{\big(E_g(k)-E_{gg}(q,r)\big)\big(E_g(k)-E_g(p)\big)}\notag\\
&=-\frac{g^2}{4\pi^3}\int_\Lambda^\vee dp^+\!\int d^2p\ \sqrt{\frac{p^+}{k^+}}\ \frac{k^ip^j\,\rho^a(\v k)\,\rho^b(-\v p)}{\v p^2\,(p^+\v k^2-k^+\v p^2)}\,\ket{g^b_j(p)}\,.
\end{align}
In both expressions the later vertex stands to the left: $\rho^b(-\v p)$ of the emission to the left of $\rho^a(\v k)$ of the absorption in Eq.~\eqref{eq:psi1rr}, and the opposite order in Eq.~\eqref{eq:psi2rr}. In Eq.~\eqref{eq:psi2rr} we have dropped the disconnected term in which the emitted gluon is reabsorbed and the incoming gluon is a spectator; it is proportional to $\ket{g^a_i(k)}$ and is excluded from the sum over $\ket n\neq\ket i$ in Eq.~\eqref{eq:pt2}. Both gluons are soft modes, $\Lambda<p^+<\vee$, without a restriction on $p^+-k^+$. With the decomposition~\eqref{eq:decomp} the energy denominator cancels in the anticommutator part (Appendix~\ref{app:tworho}), and
\begin{align}\label{eq:rrsum}
\ket{(\Psi^a_i(k))^1_{g\rho\rho}}+\ket{(\Psi^a_i(k))^2_{g\rho\rho}}=\ket{(\Psi^a_i(k))_{g\{\rho,\rho\}}}+\ket{(\Psi^a_i(k))_{g[\rho,\rho]}}\,,
\end{align}
with
\begin{align}
\ket{(\Psi^a_i(k))_{g\{\rho,\rho\}}}&=-\frac{g^2}{8\pi^3}\int_\Lambda^\vee\frac{dp^+}{\sqrt{k^+p^+}}\int d^2p\ \frac{k^ip^j}{\v k^2\v p^2}\,\big\{\rho^a(\v k),\rho^b(-\v p)\big\}\ket{g^b_j(p)}\,,\label{eq:psi_anti}\\
\ket{(\Psi^a_i(k))_{g[\rho,\rho]}}&=\frac{ig^2f^{abc}}{8\pi^3}\int_\Lambda^\vee\frac{dp^+}{\sqrt{k^+p^+}}\int d^2p\ \frac{k^ip^j}{\v k^2\v p^2}\,\frac{k^+\v p^2+p^+\v k^2}{k^+\v p^2-p^+\v k^2}\,\rho^c(\v k-\v p)\ket{g^b_j(p)}\,.\label{eq:psi_comm}
\end{align}
The anticommutator part is the product of the amplitude for the absorption of the incoming gluon and the amplitude for the emission of a Weizs\"acker--Williams gluon, symmetrised in the order of the two; it is regular at $E_g(p)=E_g(k)$. The commutator part is linear in $\rho$ and has the principal-value denominator. \begin{figure}[t]
\centering
\begin{tabular}{cc}
\scalebox{0.8}{\tikzfig{gtovactog_tworho}}&\scalebox{0.8}{\tikzfig{gg_tworho}}\\[2mm]
(a)&(b)
\end{tabular}
\caption{Transition of the incoming gluon into a gluon of different momentum with two valence vertices: (a) absorption followed by emission, Eq.~\eqref{eq:psi1rr}; (b) emission followed by absorption, Eq.~\eqref{eq:psi2rr}.}
\label{fig:g_g_two}
\end{figure}

\paragraph{The coefficient $\mf B_3$.} Comparing the sum of Eqs.~\eqref{eq:psi_inst}, \eqref{eq:psi1}, \eqref{eq:psi2}, \eqref{eq:psi_anti} and~\eqref{eq:psi_comm} with the third term of Eq.~\eqref{eq:Omegag}, i.e.\ multiplying the coefficient of $\ket{g^b_j(p)}$ by $(2\pi)^3$, we obtain
\begin{align}\label{eq:B3mom}
\mf B^{ba}_{3ji}(k,p)=\ &\frac{ig^2f^{abc}\sqrt{p^+k^+}\,(k^++p^+)}{(p^+-k^+)^2}\,\frac{\rho^c(\v k-\v p)\,\delta_{ij}}{p^+\v k^2-k^+\v p^2}\,\big[\Theta_>+\Theta_<\big]
\notag\\&
+\frac{2ig^2f^{abc}\rho^c(\v k-\v p)\,(k-p)^l\sqrt{k^+p^+}}{(k^+-p^+)^2}\,\frac{k^mp^+-k^+p^m}{\v k^2p^+-\v p^2k^+}
\notag\\&\quad\times\Big[\frac{\delta_{jm}\delta_{il}}{p^+}-\frac{\delta_{im}\delta_{jl}}{k^+}-\frac{\delta_{lm}\delta_{ij}}{p^+-k^+}\Big]
\notag\\&\quad\times
\Bigg[\frac{p^+-k^+}{(\v k-\v p)^2}\,\Theta_>+\frac{\Theta_<}{\frac{\v k^2}{k^+}-\frac{(\v k-\v p)^2}{k^+-p^+}-\frac{\v p^2}{p^+}}\Bigg]
\notag\\&
-\frac{g^2k^ip^j}{\v k^2\v p^2\sqrt{k^+p^+}}\Big[\{\rho^a(\v k),\rho^b(-\v p)\}-if^{abc}\rho^c(\v k-\v p)\,\frac{k^+\v p^2+p^+\v k^2}{k^+\v p^2-p^+\v k^2}\Big]\,,
\end{align}
with
\begin{align}
\Theta_>\equiv\Theta(p^+-k^+-\Lambda)\,,\qquad \Theta_<\equiv\Theta(k^+-\Lambda-p^+)\,,
\end{align}
and $\Lambda<k^+,p^+<\vee$. The first argument of $\mf B_3$ is the momentum of the incoming gluon, with the indices $(a,i)$, and the second that of the outgoing gluon, with the indices $(b,j)$. The first line is the instantaneous interaction, the second to fourth lines are the terms with the three-gluon vertex, and the last line contains the two terms with two valence vertices. The denominators $p^+\v k^2-k^+\v p^2$ are principal values.

\subsubsection{Two outgoing gluons}\label{sec:twoout}

At order $g$ there are two contributions with two gluons in the final state (Figure~\ref{fig:g_gg}).

\paragraph{$\rho$ independent part.} The incoming gluon splits into two soft gluons through the three-gluon vertex. This is the first-order term of Eq.~\eqref{eq:pt2} with the two-gluon states as $\ket n$,
\begin{align}\label{eq:psigg_def}
\ket{(\Psi^a_i(k))_{gg}}=\frac12\int d^3p\,d^3q\ \ket{g^c_l(q)\,g^b_j(p)}\,\frac{\braket{g^c_l(q)\,g^b_j(p)|H_{ggg}|g^a_i(k)}}{E_g(k)-E_{gg}(p,q)}\,.
\end{align}
The $\delta$-function of Eq.~\eqref{eq:m3} sets $q=k-p$. With $\xi=p^+/k^+$ the three brackets of Eq.~\eqref{eq:m3} become $2(p-\xi k)^i$, $-\frac2\xi(p-\xi k)^j$ and $-\frac{2}{1-\xi}(p-\xi k)^l$, and the energy denominator is $E_g(k)-E_{gg}(p,k-p)=-(\v p-\xi\v k)^2/[2\xi(1-\xi)k^+]$, which never vanishes. Hence
\begin{align}\label{eq:psigg}
\ket{(\Psi^a_i(k))_{gg}}=\ &-\frac{igf^{abc}\sqrt{k^+}}{4\pi^{3/2}}\int_{\Lambda/k^+}^{1-\Lambda/k^+}d\xi\,\sqrt{\xi(1-\xi)}\int d^2p\ \frac{p^m-\xi k^m}{(\v p-\xi\v k)^2}\notag\\
&\times\Big[\delta_{im}\delta_{jl}-\frac1\xi\,\delta_{jm}\delta_{il}-\frac{1}{1-\xi}\,\delta_{lm}\delta_{ij}\Big]\ket{g^c_l\big((1-\xi)k^+,\v k-\v p\big)\,g^b_j(\xi k^+,\v p)}\,.
\end{align}
Both gluons are soft modes, which gives the limits of the $\xi$ integral. The integrand is symmetric under the exchange of the two gluons, $(b,j,\xi,\v p)\leftrightarrow(c,l,1-\xi,\v k-\v p)$. Comparing with the last term of Eq.~\eqref{eq:Omegag},
\begin{align}\label{eq:C1mom}
\mf C^{bca}_{1jki}(k,p)=\ &-\frac{igf^{abc}\,(p^mk^+-p^+k^m)\sqrt{k^+p^+(k^+-p^+)}}{\sqrt2\,(k^+\v p-p^+\v k)^2}\notag\\&\times\Big[\frac{\delta_{im}\delta_{jk}}{k^+}-\frac{\delta_{ik}\delta_{mj}}{p^+}-\frac{\delta_{km}\delta_{ij}}{k^+-p^+}\Big]\,\Theta(k^+-\Lambda-p^+)\,,
\end{align}
with $p^+>\Lambda$. In this coefficient $k$ is the momentum of the incoming gluon $(a,i)$, $p$ the momentum of the outgoing gluon $(b,j)$, and the second outgoing gluon $(c,k)$ carries $k-p$; in the bracket the polarisation index $k$ is not to be confused with the momentum. $\mf C_1$ is independent of $\rho$ and obeys the symmetry relation~\eqref{eq:symm}.

\paragraph{One $\rho$ part.} The valence charge emits a gluon, and the incoming gluon is a spectator. With Eq.~\eqref{eq:m2},
\begin{align}\label{eq:psiggrho}
\ket{(\Psi^a_i(k))_{gg\rho}}&=\frac12\int d^3p\,d^3q\ \ket{g^c_l(q)\,g^b_j(p)}\,\frac{\braket{g^c_l(q)\,g^b_j(p)|H_g|g^a_i(k)}}{E_g(k)-E_{gg}(p,q)}\notag\\
&=-\frac{g}{2\pi^{3/2}}\int_\Lambda^\vee\frac{dp^+}{\sqrt{p^+}}\int d^2p\ \frac{p^j}{\v p^2}\,\rho^b(-\v p)\,\ket{g^a_i(k)\,g^b_j(p)}\,,
\end{align}
where the two terms of Eq.~\eqref{eq:m2} are equal and the denominator is $E_g(k)-E_{gg}(k,p)=-E_g(p)$. This is the leading-order wave function~\eqref{eq:PsiLO} in the presence of the incoming gluon. Comparing with the fourth term of Eq.~\eqref{eq:Omegag} reproduces $\mf A^{(1)}_2$ of Eq.~\eqref{eq:A1mom}, as it must: the coefficient of $a^\dagger$ in $\Omega$ does not depend on the state on which $\Omega$ acts.

\begin{figure}[t]
\centering
\begin{tabular}{cc}
\scalebox{0.8}{\tikzfig{g_togg}}&\scalebox{0.8}{\tikzfig{gtog_plusg}}\\[2mm]
(a)&(b)
\end{tabular}
\caption{Two gluons in the final state at order $g$: (a) splitting of the incoming gluon, Eq.~\eqref{eq:psigg_def}; (b) emission of a gluon by the valence charge with the incoming gluon as a spectator, Eq.~\eqref{eq:psiggrho}.}
\label{fig:g_gg}
\end{figure}

\subsubsection{Three outgoing gluons}\label{sec:threeout}

Three gluons appear in the dressed one-gluon state at order $g^2$: through the four-gluon vertex, through two successive splittings, through a splitting accompanied by an emission from the valence charge in the two time orderings, and through two emissions from the valence charge. In the expansion~\eqref{eq:OmegaExp} they correspond to the term $\mf B_2\,a^\dagger a^\dagger$, with the incoming gluon as a spectator, and to terms with three creation operators and one annihilation operator, which are among the terms denoted by dots. These components were computed in~\cite{Radhakrishnan:2026soft}. They are not needed here, and we do not reproduce them. The reason is a counting of orders. As shown in Section~\ref{sec:observable_def}, the cross section at order $g^4$ requires the overlap of $\Omega\ket0$ after the scattering with a dressed one-gluon state to order $g^3$ and with a dressed two-gluon state to order $g^2$. The three-gluon component of $\Omega\ket0$ is of order $g^3$, and the three-gluon component of $\Omega\ket{g^a_i(k)}$ is of order $g^2$; their product is of order $g^5$. In the same way the three-gluon components of a dressed two-gluon state, of order $g$, multiply the three-gluon component of $\Omega\ket0$ and contribute to the amplitude at order $g^4$. For the same reason neither the two-gluon component of $\Omega\ket{g^a_i(k)}$ at order $g^3$ nor the four-gluon vertex $H_{gggg}$ is needed: all components of the dressed one-gluon state that enter are those computed in Sections~\ref{sec:vacout}--\ref{sec:twoout}, and $\mf A^{(3)}_1$.

\NOTE{The three-gluon components are not derived here; it is assumed, as in the calculation file, that the normal-ordered $H_{gggg}$ gives no contribution to $\mf A^{(3)}_2$. To be confirmed by the authors.}

\subsection{Summary of the wave function}\label{sec:summarywf}

Collecting the components, the dressed one-gluon state to order $g^2$, restricted to at most two gluons, is
\begin{align}\label{eq:psi_sum}
\ket{\Psi^a_i(k)}=\ &\mf N\,\ket{g^a_i(k)}+\int\frac{d^3p}{(2\pi)^3}\,\Big[\mf B^{ba}_{3ji}(k,p)\Big]_{\rho^0}\ket{g^b_j(p)}\notag\\
&+\ket{(\Psi^a_i(k))_\rho}+\ket{(\Psi^a_i(k))_{gg}}+\ket{(\Psi^a_i(k))_{gg\rho}}\notag\\
&+\ket{(\Psi^a_i(k))^1_{g\rho\,{\rm inst}}}+\ket{(\Psi^a_i(k))^2_{g\rho\,{\rm inst}}}+\ket{(\Psi^a_i(k))^1_{g\rho}}+\ket{(\Psi^a_i(k))^2_{g\rho}}\notag\\
&+\ket{(\Psi^a_i(k))_{g\{\rho,\rho\}}}+\ket{(\Psi^a_i(k))_{g[\rho,\rho]}}+\dots
\end{align}
The second line contains the components of order $g$, Eqs.~\eqref{eq:psi_rho}, \eqref{eq:psigg} and~\eqref{eq:psiggrho}; the third and fourth lines the components of order $g^2$ with one gluon, Eqs.~\eqref{eq:psi_inst}, \eqref{eq:psi1}, \eqref{eq:psi2}, \eqref{eq:psi_anti} and~\eqref{eq:psi_comm}. The first line is the component along the incoming gluon itself, the square bracket of Eq.~\eqref{eq:pt2}. It is fixed by the normalisation
\begin{align}\label{eq:normcond}
\braket{\Psi^{a'}_{i'}(k')|\Psi^a_i(k)}=\delta^{aa'}\delta_{ii'}\,\delta^{(3)}(k-k')\,,
\end{align}
which we now verify to order $g^2$. The terms of order $g^2$ in the overlap are the overlaps of two components of order $g$ and the overlaps of $\ket{g}$ with the components of order $g^2$.

The verification is carried out in Appendix~\ref{app:norm}; the result is the following. (i) The terms with two charges that are not proportional to $\delta^{(3)}(k-k')$ cancel among the overlaps of $\ket{(\Psi)_\rho}$, $\ket{(\Psi)_{gg\rho}}$ and $\ket{(\Psi)_{g\{\rho,\rho\}}}$, for non-commuting charges; the remaining term fixes the coefficient of $\ket{g^a_i(k)}$ to be $\mf N=1+\mf N_2$ with the same $\mf N_2$ as in the vacuum sector, Eq.~\eqref{eq:N2mom}. (ii) The terms linear in $\rho$ cancel; for the terms with the three-gluon vertex this requires the overlaps of the two-gluon components, Eq.~\eqref{eq:ov4}. (iii) The $\rho$-independent probability for the incoming gluon to split, Eq.~\eqref{eq:ov5}, is compensated by a term of $\mf B_3$ proportional to $\delta^{(3)}(k-p)$, Eq.~\eqref{eq:B3rho0}, the wave-function renormalisation of the soft gluon. It multiplies the scaleless integral $\int d^2\tilde p/\tilde{\v p}^2$ and vanishes in dimensional regularisation.

\NOTE{Setting the term~\eqref{eq:B3rho0} to zero relies on $\int d^2\tilde p/\tilde{\v p}^2=0$. If ultraviolet and collinear divergences are treated separately in the cross section, it has to be kept. To be decided by the authors.}

\subsection{Unitarity of \texorpdfstring{$\Omega$}{Omega} and the relations between the coefficients}\label{sec:unitarity}

The coefficients of Eq.~\eqref{eq:OmegaExp} are not independent. Since $\Omega$ maps an orthonormal basis onto an orthonormal basis,
\begin{align}\label{eq:unit}
\Omega^\dagger\,\Omega=1\,.
\end{align}
We impose this condition order by order in $g$: $\Omega^\dagger\Omega$ is brought to normal order with the commutator~\eqref{eq:aadagger}, keeping the order of the coefficients, and the coefficient of each normal-ordered product of $a^\dagger$ and $a$ is set to zero. The algebra is given in Appendix~\ref{app:unitarity}. At order $g$ and $g^2$ one finds
\begin{align}\label{eq:unit_summary}
&\mf A^{(1)\dagger b}_{1j}(p)=-\mf A^{(1)b}_{2j}(p)\,,\qquad
\mf C^{\dagger dcb}_{2lkj}(p,q)=-\mf C^{dcb}_{1lkj}(p,q)\,,\notag\\
&\mf N_2=-\frac12\int\frac{d^3p}{(2\pi)^3}\,\mf A^{(1)\dagger b}_{2j}(p)\,\mf A^{(1)b}_{2j}(p)\,,
\end{align}
together with Eq.~\eqref{eq:B1unit}, which expresses $\mf B_1$ through $\mf B_2$, $\mf C_1\mf A^{(1)}_2$ and the symmetrised product of two $\mf A^{(1)}_2$, and Eq.~\eqref{eq:B3unit}, a relation between $\mf B_3$, $\mf B_3^\dagger$ and products of $\mf A^{(1)}_2$ and $\mf C_1$. At order $g^3$ the coefficient of $a^\dagger$ gives
\begin{align}\label{eq:A31unit}
\mf A^{(3)\dagger b}_{1j}(p)=\ &-\mf A^{(3)b}_{2j}(p)-\big[\mf N_2,\mf A^{(1)b}_{2j}(p)\big]
+\int\frac{d^3q}{(2\pi)^3}\Big[2\,\mf A^{(1)c}_{1k}(q)\,\mf B^{cb}_{2kj}(p,q)\notag\\
&-\mf B^{\dagger cb}_{3kj}(p,q)\,\mf A^{(1)c}_{2k}(q)
-2\,\mf C^{\dagger dcb}_{1lkj}(p,q)\,\mf B^{cd}_{2kl}(q,p-q)\Big]\,.
\end{align}
The first two relations of Eq.~\eqref{eq:unit_summary} are satisfied by Eqs.~\eqref{eq:A1mom} and~\eqref{eq:A11mom} and determine $\mf C_2$; the third reproduces Eq.~\eqref{eq:N2mom}. Eq.~\eqref{eq:B3unit} is satisfied by Eq.~\eqref{eq:B3mom} and fixes the relative sign of its term $\propto\Theta_>$ and $\mf C_1$.

\paragraph{The order of the operators.} For commuting charges the commutator in Eq.~\eqref{eq:A31unit} is absent. Here $\mf N_2$ is bilinear in the charges at two different points and does not commute with $\mf A^{(1)}_2$; the commutator, Eq.~\eqref{eq:commN2A}, is a term with two charges $\propto g^3f^{abc}\{\rho^b,\rho^c\}\log(\vee/\Lambda)$, equal to the term $-2k^i\log(\vee/\Lambda)$ of $[\mf A^{(3)}_2]_{\rho\rho}$. The order of the charges in the three-$\rho$ term~\eqref{eq:A3threerho} is therefore not a matter of convention, and with that order
\begin{align}\label{eq:N2Acomb}
-\big[\mf A^{(3)}_2\big]_{\rho\rho\rho}-\big[\mf N_2,\mf A^{(1)}_2\big]=-\mf N_2\,\mf A^{(1)}_2\,.
\end{align}

\paragraph{What is fixed by what.} The diagrammatic calculation fixes $\mf N$, $\mf A_2=\mf A^{(1)}_2+\mf A^{(3)}_2$ and $\mf B_2$ (vacuum initial state) and $\mf A^{(1)}_1$, $\mf B_3$ and $\mf C_1$ (one-gluon initial state). Unitarity fixes $\mf C_2$, $\mf B_1$ and $\mf A^{(3)}_1$ through Eqs.~\eqref{eq:C2unit}, \eqref{eq:B1unit} and~\eqref{eq:A31unit}. The relations~\eqref{eq:A11unit}, \eqref{eq:N2unit} and~\eqref{eq:B3unit} involve only coefficients fixed by the diagrammatic calculation and are checks. We have verified the relations of this subsection in two ways: by a symbolic normal ordering of $\Omega^\dagger\Omega$ with non-commuting coefficients for a discrete set of modes, and, for the perturbative series of Section~\ref{sec:lcpt} on which they rest, by a numerical construction of $\Omega$ with matrices in place of the charges.

\subsection{Summary of the coefficients}\label{sec:coefficients}

Table~\ref{tab:coeffs} lists the coefficients of $\Omega$, the order at which they are needed, what determines them, and the equations in which they are given. After the relations~\eqref{eq:A11unit} and~\eqref{eq:C2unit} have been used, the expansion~\eqref{eq:OmegaExp} reads
\begin{align}\label{eq:OmegaFinal}
\Omega = \ &\mf{N}
+\int\frac{d^3p}{(2\pi)^3}\Big[\mf{A}^b_{j}(p)\,a^{\dagger b}_j(p)-\mf{A}^{\dagger b}_{j}(p)\,a^{b}_j(p)+\mf{A}^{(3)b}_{2j}(p)\,a^{\dagger b}_j(p)+\mf{A}^{(3)b}_{1j}(p)\,a^{b}_j(p)\Big]
\notag\\&
+\int\frac{d^3p}{(2\pi)^3}\frac{d^3q}{(2\pi)^3}\Big[\mf{B}^{cb}_{1kj}(p,q)\,a^b_j(p)a^c_k(q)+\mf{B}^{cb}_{2kj}(p,q)\,a^{\dagger b}_j(p)a^{\dagger c}_k(q)
\notag\\&\hspace{3.4cm}+\mf{B}^{cb}_{3kj}(p,q)\,a^{\dagger c}_k(q)a^b_j(p)\Big]
\notag\\&
+\int\frac{d^3p}{(2\pi)^3}\frac{d^3q}{(2\pi)^3}\Big[\mf{C}^{dcb}_{lkj}(p,q)\,a^{\dagger d}_l(q)a^{\dagger c}_k(p-q)a^b_j(p)
\notag\\&\hspace{3.4cm}-\mf{C}^{\dagger dcb}_{lkj}(p,q)\,a^{\dagger b}_j(p)a^c_k(p-q)a^d_l(q)\Big]+\dots,
\end{align}
with $\mf A\equiv\mf A^{(1)}_2$ and $\mf C\equiv\mf C_1$.

\begin{table}[t]
\centering
\renewcommand{\arraystretch}{1.3}
\footnotesize\setlength{\tabcolsep}{4pt}
\begin{tabular}{@{}l c >{\raggedright\arraybackslash}p{0.37\textwidth} c c@{}}
\toprule
coefficient & order & fixed by & momentum space & coordinate space\\
\midrule
$\mf A^{(1)}_2\equiv\mf A$ & $g$ & $\Omega\ket0$, one gluon; first order in $H_g$ & \eqref{eq:A1mom} & \eqref{eq:A1X}\\
$\mf A^{(1)}_1$ & $g$ & $\Omega\ket g$, vacuum; check: unitarity, Eq.~\eqref{eq:A11unit} & \eqref{eq:A11mom} & \eqref{eq:A1X}\\
$\mf C_1\equiv\mf C$ & $g$ & $\Omega\ket g$, two gluons; first order in $H_{ggg}$ & \eqref{eq:C1mom} & \eqref{eq:C1X}\\
$\mf C_2$ & $g$ & unitarity, Eq.~\eqref{eq:C2unit} & \eqref{eq:C2unit} & \eqref{eq:unitX}\\
$\mf N_2$ & $g^2$ & norm of $\Omega\ket0$ and of $\Omega\ket g$; unitarity, Eq.~\eqref{eq:N2unit} & \eqref{eq:N2mom} & \eqref{eq:N2X}\\
$\mf B_2$ & $g^2$ & $\Omega\ket0$, two gluons; second order~\cite{Lublinsky:2016meo} & \eqref{eq:B2mom} & \eqref{eq:B2X}\\
$\mf B_3$ & $g^2$ & $\Omega\ket g$, one gluon; second order, denominators $E_g(k)-E_n$; check: Eq.~\eqref{eq:B3unit} & \eqref{eq:B3mom} & \eqref{eq:B3X}\\
$\mf B_1$ & $g^2$ & unitarity, Eq.~\eqref{eq:B1unit} & \eqref{eq:B1unit} & \eqref{eq:unitX}\\
$\mf A^{(3)}_2$ & $g^3$ & $\Omega\ket0$, one gluon; third order~\cite{Lublinsky:2016meo} & \eqref{eq:A3onerho}--\eqref{eq:A3threerho} & \eqref{eq:A3onerhoX}--\eqref{eq:A3threerhoX}\\
$\mf A^{(3)}_1$ & $g^3$ & unitarity, Eq.~\eqref{eq:A31unit} & \eqref{eq:A31unit} & \eqref{eq:unitX}\\
\bottomrule
\end{tabular}
\caption{The coefficients of the evolution operator $\Omega$, Eq.~\eqref{eq:OmegaExp}, that are needed for the single inclusive gluon cross section at next-to-leading order. $\Omega\ket0$ and $\Omega\ket g$ denote the matching to the dressed vacuum (Section~\ref{sec:vacin}) and to the dressed one-gluon state (Section~\ref{sec:onegin}).}
\label{tab:coeffs}
\end{table}

The expansion of $\Omega$ in transverse coordinate space, Eq.~\eqref{eq:OmegaX}, and the coefficients in coordinate space, in which they are used in the cross section, are given in Appendix~\ref{app:coefficients}.

With the coefficients of Table~\ref{tab:coeffs} the operator $\Omega$ is known to the accuracy required. In Section~\ref{sec:cross} we use it to set up the calculation of the single inclusive gluon cross section.

% ======================================================================
\section{Single inclusive gluon production: setting up the cross section}\label{sec:cross}
% ======================================================================

We now use the operator $\Omega$ to write the single inclusive gluon spectrum in the collision of the dilute projectile with a dense target. In this section the spectrum is expressed through the coefficients of $\Omega$ and the Wilson lines of the target, to order $g^4$; the coefficients are not yet inserted. The result is organised in three groups of terms, which are the starting point of the evaluation in the following sections.

\subsection{The observable}\label{sec:observable_def}

\paragraph{Eikonal scattering.} Before the collision the projectile is in the state $\Omega\ket0$, i.e.\ $\Omega\ket0\otimes\ket v$ with the valence state $\ket v$. The target is described by the eikonal $S$-matrix~\cite{Kovner:2005en,Lublinsky:2016meo}: at high energy a soft gluon and the valence charges keep their transverse positions and longitudinal momenta while they cross the target, and their colour is rotated by the adjoint Wilson line $U(\v x)$ of the target field at their transverse position,
\begin{align}\label{eq:Srotation}
S^\dagger\rho^a(\v x)\,S&=U^{ab}(\v x)\,\rho^b(\v x)\,,\notag\\
S^\dagger a^a_i(k^+,\v x)\,S&=U^{ab}(\v x)\,a^b_i(k^+,\v x)\,,\qquad
S^\dagger a^{\dagger a}_i(k^+,\v x)\,S=U^{ab}(\v x)\,a^{\dagger b}_i(k^+,\v x)\,.
\end{align}
The matrix $U$ is real and orthogonal, $U^\dagger=U^T$, $U^TU=1$, and $S\ket0=\ket0$ for the soft vacuum. After the collision the state is $S\,\Omega\ket0$.

\paragraph{The number of produced gluons.} A gluon observed in the final state is a dressed gluon, i.e.\ a state $\Omega\,a^\dagger\ket0$ and not a bare quantum of the cloud of the projectile. The number of dressed gluons is counted by the number operator rotated by $\Omega$, $\Omega\,a^\dagger a\,\Omega^\dagger$~\cite{Kovner:2006wr}. The single inclusive spectrum is therefore
\begin{align}\label{eq:observable}
\frac{d^3N}{d^3k}=\frac{1}{(2\pi)^3}\int_{w,w'}e^{-i\v k\cdot(\v w-\v w')}\,
\Big\langle\bra0\,\Omega^\dagger S^\dagger\Omega\;a^{\dagger a}_i(k^+,\v w')\,a^a_i(k^+,\v w)\;\Omega^\dagger S\,\Omega\ket0\Big\rangle\,,
\end{align}
where $\langle\dots\rangle$ is the average over the valence state, taken last. The factor $(2\pi)^{-3}$ and the phase follow from Eqs.~\eqref{eq:aadagger} and~\eqref{eq:FTconventions}: the density of gluons in $d^3k$ is $(2\pi)^{-3}a^{\dagger a}_i(k)a^a_i(k)$, and $a^{\dagger}(k^+,\v k)\,a(k^+,\v k)=\int_{w,w'}e^{-i\v k\cdot(\v w-\v w')}a^{\dagger}(k^+,\v w')\,a(k^+,\v w)$. The amplitude stands at the unprimed point $\v w$ and the conjugate amplitude at the primed point $\v w'$. We use this convention, the phase $e^{-i\v k\cdot(\text{unprimed}-\text{primed})}$ with the conjugated factor to the left, in all expressions below. Without a target, $S=1$, the spectrum vanishes to all orders because $\Omega^\dagger\Omega=1$.

\paragraph{Barred coefficients.} We insert $SS^\dagger=1$ to the left of the number operator and between $\Omega^\dagger$ and $S$, and define
\begin{align}\label{eq:Omegabar}
\bar\Omega\equiv S^\dagger\,\Omega\,S\,.
\end{align}
By Eq.~\eqref{eq:Srotation}, $\bar\Omega$ is given by the expansion~\eqref{eq:OmegaX} with every charge and every gluon operator rotated. A bar on a coefficient denotes the replacement of every charge in it by the rotated charge,
\begin{align}\label{eq:bar}
\bar{\mf A}\equiv\mf A\big|_{\rho^a(\v u)\to U^{ab}(\v u)\rho^b(\v u)}\,,\qquad\text{and the same for }\ \bar{\mf N},\ \bar{\mf B}_{1,2,3}\,.
\end{align}
The coefficients $\mf C_{1,2}$ do not depend on $\rho$, and $\bar{\mf C}_{1}=\mf C_1$; we keep the bar as a reminder that the coefficient originates from $\bar\Omega$. Since the barred charges obey the same algebra~\eqref{eq:rhoalgebra}, the barred coefficients obey the unitarity relations of Section~\ref{sec:unitarity}. With $S^\dagger a^\dagger a\,S=Ua^\dagger\,Ua$ the spectrum becomes
\begin{align}\label{eq:observable2}
\frac{d^3N}{d^3k}=\frac{1}{(2\pi)^3}\int_{w,w'}e^{-i\v k\cdot(\v w-\v w')}\,\Big\langle\bra0\,\Omega^\dagger\bar\Omega\;U^{ab'}(\v w')\,a^{\dagger b'}_i(k^+,\v w')\;U^{ab}(\v w)\,a^b_i(k^+,\v w)\;\bar\Omega^\dagger\Omega\ket0\Big\rangle\,.
\end{align}

\paragraph{Amplitudes.} Between the two gluon operators we insert the complete set of soft Fock states. The state $a(k^+,\v w)\,\bar\Omega^\dagger\Omega\ket0$ is projected on the vacuum, on one-gluon states, and so on; these projections are the overlaps of $\bar\Omega^\dagger\Omega\ket0$ with states with one, two, \dots\ gluons, one of which is the measured gluon. With the normalisation~\eqref{eq:states},
\begin{align}\label{eq:N_amplitudes}
\frac{d^3N}{d^3k}=\frac{1}{(2\pi)^3}\Big\langle\mathcal A^{a\dagger}_i(k)\,\mathcal A^a_i(k)+\int dq^+\!\int_{v}\mathcal A^{ab\dagger}_{ij}(k;q^+,\v v)\,\mathcal A^{ab}_{ij}(k;q^+,\v v)\Big\rangle+\dots,
\end{align}
with the one-gluon and two-gluon amplitudes
\begin{align}\label{eq:amplitudes_def}
\mathcal A^a_i(k)&=\sqrt{2\pi}\int_{w}e^{-i\v k\cdot\v w}\,U^{ab}(\v w)\,\bra{g^b_i(k^+,\v w)}\bar\Omega^\dagger\Omega\ket0\,,\notag\\
\mathcal A^{ab}_{ij}(k;q^+,\v v)&=\sqrt{2\pi}\int_{w}e^{-i\v k\cdot\v w}\,U^{ac}(\v w)\,\bra{g^b_j(q^+,\v v)\,g^c_i(k^+,\v w)}\bar\Omega^\dagger\Omega\ket0\,.
\end{align}
The amplitudes are operators on $\mathcal H_{\rm val}$. The second gluon of $\mathcal A^{ab}_{ij}$, with longitudinal momentum $q^+$ at the transverse position $\v v$, is not measured and is integrated over. The dots in Eq.~\eqref{eq:N_amplitudes} are the contributions of final states with three or more gluons; the amplitude for $n$ gluons is of order $g^n$, so that they are of order $g^6$. The amplitudes are evaluated with the expansion~\eqref{eq:OmegaX}: $\Omega$ acting on $\ket0$ involves only $\mf N$, $\mf A_2$ and $\mf B_2$, Eq.~\eqref{eq:Omegavac}, and $\bar\Omega^\dagger$ then acts on states with up to two gluons. This is where the coefficients $\mf B_3$ and $\mf C_1$ of the dressed one-gluon state enter, and it confirms the counting of Section~\ref{sec:threeout}: the one-gluon amplitude is needed to order $g^3$ and the two-gluon amplitude to order $g^2$, for which the coefficients of Table~\ref{tab:coeffs} are sufficient.

From here on $\mf A\equiv\mf A^{(1)}_2$, $\mf A^{(3)}\equiv\mf A^{(3)}_2$, and the coefficients $\mf A_1$, $\mf B_1$, $\mf C_2$ of $\bar\Omega^\dagger$ are eliminated with the unitarity relations~\eqref{eq:unitX}.

\subsection{Expansion in the coupling and the leading order}\label{sec:LO}

The amplitudes are expanded in $g$,
\begin{align}\label{eq:ampexp}
\mathcal A^a_i(k)=\mathcal A^{a}_{{\rm LO}\,i}(k)+\mathcal A^a_{(3)\,i}(k)+O(g^5)\,,\qquad \mathcal A^{ab}_{ij}=O(g^2)\,,
\end{align}
with $\mathcal A_{\rm LO}=O(g)$ and $\mathcal A_{(3)}=O(g^3)$. At order $g$, $\Omega=1+\int\frac{dp^+}{2\pi}\int_x\big[\mf A^b_j(p^+,-\v x)\,a^{\dagger b}_j(p^+,\v x)-\mf A^{\dagger b}_j(p^+,-\v x)\,a^b_j(p^+,\v x)\big]$, and acting with $\bar\Omega^\dagger\Omega$ on the vacuum only the creation parts survive,
\begin{align}\label{eq:LOstate}
\bar\Omega^\dagger\Omega\ket0\Big|_{O(g)}=\int\frac{dp^+}{2\pi}\int_x\Big[\mf A^c_j(p^+,-\v x)-\bar{\mf A}^b_j(p^+,-\v x)\,U^{bc}(\v x)\Big]a^{\dagger c}_j(p^+,\v x)\ket0\,.
\end{align}
The first term is the gluon created by $\Omega$ before the collision, the second the gluon removed by $\bar\Omega^\dagger$, with the amplitude of the rotated charges. Projecting on $\bra{g^b_i(k^+,\v w)}$ and using $U^{ab}U^{b'b}=\delta^{ab'}$,
\begin{align}\label{eq:ALO}
\mathcal A^{a}_{{\rm LO}\,i}(k)&=\int_{w}e^{-i\v k\cdot\v w}\Big[U^{ab}(\v w)\,\mf A^b_i(k^+,-\v w)-\bar{\mf A}^a_i(k^+,-\v w)\Big]\,.
\end{align}
The leading-order spectrum in terms of the coefficients is therefore
\begin{align}\label{eq:LOcoeff}
\frac{d^3N_{\rm LO}}{d^3k}=\ &\frac{1}{(2\pi)^3}\int_{w,w'}e^{-i\v k\cdot(\v w-\v w')}\notag\\&\times\Big\langle\Big[\mf A^{\dagger b'}_i(k^+,-\v w')\,U^{ab'}(\v w')-\bar{\mf A}^{\dagger a}_i(k^+,-\v w')\Big]\notag\\
&\times\Big[U^{ab}(\v w)\,\mf A^b_i(k^+,-\v w)-\bar{\mf A}^a_i(k^+,-\v w)\Big]\Big\rangle\,.
\end{align}
As a check we insert the Weizs\"acker--Williams amplitude~\eqref{eq:A1X}. Then
\begin{align}\label{eq:ALOexpl}
\mathcal A^{a}_{{\rm LO}\,i}(k)=\frac{ig}{\sqrt2\,\pi\sqrt{k^+}}\int_{w,x}e^{-i\v k\cdot\v w}\,\Kv^i(\v x-\v w)\,\big[U(\v w)-U(\v x)\big]^{ab}\rho^b(\v x)\,,
\end{align}
and
\begin{align}\label{eq:LO}
\frac{d^3N_{\rm LO}}{d^3k}=\ &\frac{1}{(2\pi)^3}\,\frac{g^2}{2\pi^2}\,\frac{1}{k^+}\int_{x,x',w,w'}e^{-i\v k\cdot(\v w-\v w')}\,
\Kv(\v x-\v w)\cdot\Kv(\v x'-\v w')
\notag\\&\times
\Big\langle\big[U(\v x')-U(\v w')\big]^{ab'}\rho^{b'}(\v x')\,\big[U(\v x)-U(\v w)\big]^{ab}\rho^{b}(\v x)\Big\rangle\,.
\end{align}
The amplitude contains the difference of the Wilson lines at the position of the gluon and at the position of the charge that emits it: a gluon is produced only if the target rotates it differently from its source. For a projectile with $\langle\rho^a(\v x)\rho^b(\v y)\rangle=\delta^{ab}n(\v x,\v y)/(N_c^2-1)$, Eq.~\eqref{eq:LO} is the known leading-order result~\cite{Kovchegov:2001sc,Kovner:2006wr} with the target operator ${\rm Tr}\{[U^\dagger(\v x')-U^\dagger(\v w')][U(\v x)-U(\v w)]\}$ and the prefactor $\alpha_s/[4\pi^4(N_c^2-1)]$, $\alpha_s=g^2/4\pi$.

\subsection{Next-to-leading order: the amplitudes}\label{sec:NLOamplitudes}

\paragraph{One gluon in the final state.} At order $g^3$ we collect all terms of $\bar\Omega^\dagger\Omega$ that connect the vacuum with a one-gluon state. $\Omega\ket0$ contains the vacuum at order $1$ and $g^2$, one gluon at order $g$ and $g^3$, and two gluons at order $g^2$, Eq.~\eqref{eq:Omegavac}; $\bar\Omega^\dagger$ maps them onto a one-gluon state with $\bar{\mf A}^\dagger_1$, with $\bar{\mf N}$ and $\bar{\mf B}_3^\dagger$, and with $\bar{\mf A}_2^\dagger$ and $\bar{\mf C}_1^\dagger$, respectively. After the elimination of $\bar{\mf A}^{(3)}_1$ with Eq.~\eqref{eq:A31unit}, the terms that this relation generates combine with the direct ones, and the result has five parts,
\begin{align}\label{eq:A3amp}
\mathcal A^a_{(3)\,i}(k)=\int_w e^{-i\v k\cdot\v w}\Big\{\mathcal A_{N}+\mathcal A_{A^{(3)}}+\mathcal A_{B_3}+\mathcal A_{AB_2}+\mathcal A_{CB_2}\Big\}^a_i(\v w)\,,
\end{align}
with
\begin{align}\label{eq:A3parts}
\mathcal A_{N}&=-\bar{\mf A}^a_i(k^+,-\v w)\,\mf N_2+U^{ab}(\v w)\,\bar{\mf N}_2\,\mf A^b_i(k^+,-\v w)-\big[\bar{\mf N}_2,\bar{\mf A}^a_i(k^+,-\v w)\big]\,,
\notag\\
\mathcal A_{A^{(3)}}&=U^{ab}(\v w)\,\mf A^{(3)b}_i(k^+,-\v w)-\bar{\mf A}^{(3)a}_i(k^+,-\v w)\,,
\notag\\
\mathcal A_{B_3}&=\int\frac{dp^+}{2\pi}\int_x\bar{\mf B}^{\dagger ba}_{3ji}(k^+,\v w;p^+,-\v x)\Big[U^{bc}(\v x)\,\mf A^c_j(p^+,-\v x)-\bar{\mf A}^b_j(p^+,-\v x)\Big]\,,
\notag\\
\mathcal A_{AB_2}&=2\int\frac{dp^+}{2\pi}\int_x\Big[U^{ad}(\v w)U^{bc}(\v x)\,\bar{\mf A}^{\dagger b}_j(p^+,-\v x)\,\mf B^{cd}_{2ji}(k^+,-\v w;p^+,-\v x)
\notag\\&\hspace{2.2cm}
-\bar{\mf A}^{\dagger b}_j(p^+,-\v x)\,\bar{\mf B}^{ba}_{2ji}(k^+,-\v w;p^+,-\v x)\Big]\,,
\notag\\
\mathcal A_{CB_2}&=2\int\frac{dp^+}{2\pi}\int_{x,z}\Big[\bar{\mf C}^{\dagger dca}_{1jki}(k^+,\v w-\v z;p^+,\v z-\v x)\,U^{db}(\v x)U^{ce}(\v z)\,\mf B^{be}_{2jk}(k^+-p^+,-\v z;p^+,-\v x)
\notag\\&\hspace{2.2cm}
-\bar{\mf C}^{\dagger bca}_{1jki}(k^+,\v w-\v z;p^+,\v z-\v x)\,\bar{\mf B}^{bc}_{2jk}(k^+-p^+,-\v z;p^+,-\v x)\Big]\,.
\end{align}
$\mathcal A_N$ contains the normalisation of the state before and after the collision and the commutator that Eq.~\eqref{eq:A31unit} produces. $\mathcal A_{A^{(3)}}$ is the one-gluon component of order $g^3$ of the wave function, scattered on the target or removed after it. In $\mathcal A_{B_3}$ a gluon of momentum $p^+$ emitted at leading order is turned into the measured gluon by its interaction with the rotated valence charges after the collision. In $\mathcal A_{AB_2}$ one of the two gluons of the two-gluon component is absorbed by the valence charges after the collision, and in $\mathcal A_{CB_2}$ the two gluons merge after the collision. In each part the two terms differ by the position of the target relative to the last interaction, and their sum vanishes for $U=1$.

The commutator in $\mathcal A_N$ is the place where the order of the operators discussed in Section~\ref{sec:unitarity} enters the cross section. With $[\mf A^{(3)}]_{\rho\rho\rho}=\mf A\,\mf N_2$, Eq.~\eqref{eq:A3threerho}, and Eq.~\eqref{eq:N2Acomb} for the barred coefficients, the terms of $\mathcal A_N$ and the three-$\rho$ part of $\mathcal A_{A^{(3)}}$ add up to
\begin{align}\label{eq:Nterms}
\mathcal A_N+\big[\mathcal A_{A^{(3)}}\big]_{\rho\rho\rho}=-\bar{\mf A}\,\mf N_2+U\,\bar{\mf N}_2\,\mf A+U\,\mf A\,\mf N_2-\bar{\mf N}_2\,\bar{\mf A}\,,
\end{align}
in which no commutator is left and which vanishes for $U=1$.

\paragraph{Two gluons in the final state.} At order $g^2$ the two-gluon state is reached from the two-gluon component of $\Omega\ket0$ with $\bar{\mf N}=1$, from its one-gluon component with $\bar{\mf A}^\dagger_1$ (a second gluon is added) and $\bar{\mf C}_2^\dagger$ (the gluon splits), and from the vacuum with $\bar{\mf B}_1^\dagger$, which is eliminated with Eq.~\eqref{eq:B1unit}. The result is
\begin{align}\label{eq:A2gamp}
\mathcal A^{ab}_{ij}(k;q^+,\v v)=\frac{2}{\sqrt{2\pi}}\int_w e^{-i\v k\cdot\v w}\Big\{\mathcal A_B+\mathcal A^{(1)}_{AA}+\mathcal A^{(2)}_{AA}-\mathcal A_{CA}\Big\}^{ab}_{ij}(\v w;q^+,\v v)\,,
\end{align}
with
\begin{align}\label{eq:A2parts}
\mathcal A_B&=U^{ac}(\v w)\,\mf B^{bc}_{2ji}(k^+,-\v w;q^+,-\v v)-U^{\dagger bc}(\v v)\,\bar{\mf B}^{ca}_{2ji}(k^+,-\v w;q^+,-\v v)\,,
\notag\\
\mathcal A^{(1)}_{AA}&=\tfrac12\Big[U^{\dagger bc}(\v v)\,\bar{\mf A}^c_j(q^+,-\v v)\,\bar{\mf A}^a_i(k^+,-\v w)\notag\\&\hspace{1.2cm}-U^{ac}(\v w)\,\bar{\mf A}^d_j(q^+,-\v v)\,U^{\dagger bd}(\v v)\,\mf A^c_i(k^+,-\v w)\Big]\,,
\notag\\
\mathcal A^{(2)}_{AA}&=\tfrac12\Big[U^{\dagger bc}(\v v)\,\bar{\mf A}^a_i(k^+,-\v w)\,\bar{\mf A}^c_j(q^+,-\v v)-\bar{\mf A}^a_i(k^+,-\v w)\,\mf A^b_j(q^+,-\v v)\Big]\,,
\notag\\
\mathcal A_{CA}&=\int_z\Big[\bar{\mf C}^{dac}_{1jik}(q^++k^+,\v z-\v w;q^+,\v w-\v v)\,U^{\dagger bd}(\v v)\,U^{ce}(\v z)\,\mf A^e_k(q^++k^+,-\v z)
\notag\\&\hspace{0.6cm}
-U^{\dagger bd}(\v v)\,\bar{\mf C}^{dac}_{1jik}(q^++k^+,\v z-\v w;q^+,\v w-\v v)\,\bar{\mf A}^c_k(q^++k^+,-\v z)\Big]\,,
\end{align}
where $U^{\dagger bc}=U^{cb}$. $\mathcal A_B$ is the two-gluon component of the wave function, scattered on the target or removed after it. $\mathcal A^{(1)}_{AA}$ and $\mathcal A^{(2)}_{AA}$ are products of two single emissions in the two orders that $\bar\Omega^\dagger$ generates; they do not combine into one term because the coefficients do not commute. $\mathcal A_{CA}$ is the splitting, after the collision, of a gluon of momentum $q^++k^+$ into the measured gluon and the second gluon. The colour index $b$ of the second gluon is that of the bare state; its rotation by $U(\v v)$ would drop out of the sum over final states.

Eqs.~\eqref{eq:A3amp}--\eqref{eq:A2parts} were verified by an independent symbolic expansion, described in Appendix~\ref{app:rows_check}.

\subsection{The cross section at order \texorpdfstring{$g^4$}{g4}: Groups and rows}\label{sec:groups}

Inserting Eqs.~\eqref{eq:ampexp} and~\eqref{eq:A2gamp} into Eq.~\eqref{eq:N_amplitudes},
\begin{align}\label{eq:NLOdecomp}
\frac{d^3N}{d^3k}=\frac{d^3N_{\rm LO}}{d^3k}+\frac{d^3N^{\rm(I)}}{d^3k}+\frac{d^3N^{\rm(II)}}{d^3k}+\frac{d^3N^{\rm(III)}}{d^3k}+O(g^6)\,,
\end{align}
where
\begin{align}\label{eq:Groups}
\frac{d^3N^{\rm(I)}}{d^3k}&=\frac{1}{(2\pi)^3}\Big\langle\mathcal A^{a\dagger}_{{\rm LO}\,i}(k)\,\mathcal A^a_{(3)\,i}(k)\Big\rangle+{\rm c.c.}\,,\notag\\
\frac{d^3N^{\rm(II)}}{d^3k}&=\frac{1}{(2\pi)^3}\int dp^+\!\int_z\Big\langle\big|\mathcal A_B\big|^2+\big|\mathcal A^{(1)}_{AA}\big|^2+\big|\mathcal A^{(2)}_{AA}\big|^2+\big|\mathcal A_{CA}\big|^2\Big\rangle\,,\notag\\
\frac{d^3N^{\rm(III)}}{d^3k}&=\frac{1}{(2\pi)^3}\int dp^+\!\int_z\Big\langle\mathcal A^{(1)\dagger}_{AA}\mathcal A_B+\mathcal A^{(2)\dagger}_{AA}\mathcal A_B-\mathcal A^{\dagger}_{CA}\mathcal A_B
\notag\\&\hspace{3.3cm}+\mathcal A^{(2)\dagger}_{AA}\mathcal A^{(1)}_{AA}-\mathcal A^{\dagger}_{CA}\mathcal A^{(1)}_{AA}-\mathcal A^{\dagger}_{CA}\mathcal A^{(2)}_{AA}\Big\rangle+{\rm c.c.}
\end{align}
Here the momentum and the position of the second gluon have been renamed, $(q^+,\v v)\to(p^+,\v z)$, and the products of the parts of the two-gluon amplitude stand for
\begin{align}\label{eq:XYdef}
\mathcal X^\dagger\mathcal Y\equiv\frac{2}{\pi}\int_{w,w'}e^{-i\v k\cdot(\v w-\v w')}\,\Big[\mathcal X^{ab}_{ij}(\v w';p^+,\v z)\Big]^\dagger\,\mathcal Y^{ab}_{ij}(\v w;p^+,\v z)\,,\qquad|\mathcal X|^2\equiv\mathcal X^\dagger\mathcal X\,,
\end{align}
where the factor $2/\pi$ is the square of the prefactor of Eq.~\eqref{eq:A2gamp}. Group~I is the interference of the leading-order amplitude with the one-gluon amplitude of order $g^3$; the second gluon is virtual. Groups~II and~III have two gluons in the final state: Group~II contains the squares of the four parts of the two-gluon amplitude, Group~III their six interferences. Groups~I and~III are written with their complex conjugates, i.e.\ the Hermitian conjugate of the operator inside the average; Group~II is Hermitian by itself. The square of $\mathcal A_{(3)}$ is of order $g^6$.

We call the individual products of coefficients \emph{rows}. Group~I has five rows, one for each part of Eq.~\eqref{eq:A3parts}; Group~II has four and Group~III six. They are listed in Table~\ref{tab:rows}. As a representative of each group we display one row in terms of the coefficients; the complete list is given in Appendix~\ref{app:rows}.

\paragraph{Group I, Row I.III.} The interference of the leading-order amplitude with $\mathcal A_{B_3}$ is
\begin{align}\label{eq:rowI_III}
\frac{d^3N^{\rm(I.III)}}{d^3k}=\ &\frac{1}{(2\pi)^3}\int_{w,w'}e^{-i\v k\cdot(\v w-\v w')}\int\frac{dp^+}{2\pi}\int_x\notag\\&\times\Big\langle\Big[\mf A^{\dagger b'}_i(k^+,-\v w')\,U^{ab'}(\v w')-\bar{\mf A}^{\dagger a}_i(k^+,-\v w')\Big]\notag\\
&\times\bar{\mf B}^{\dagger ba}_{3ji}(k^+,\v w;p^+,-\v x)\Big[U^{bc}(\v x)\,\mf A^c_j(p^+,-\v x)-\bar{\mf A}^b_j(p^+,-\v x)\Big]\Big\rangle+{\rm c.c.}
\end{align}

\paragraph{Group II, Row II.I.} The square of $\mathcal A_B$ is
\begin{align}\label{eq:rowII_I}
\frac{d^3N^{\rm(II.I)}}{d^3k}=\ &\frac{4}{(2\pi)^3}\int\frac{dp^+}{2\pi}\int_{z,w,w'}e^{-i\v k\cdot(\v w-\v w')}\notag\\
&\times\Big\langle\Big[\mf B^{\dagger bc'}_{2ji}(k^+,-\v w';p^+,-\v z)\,U^{ac'}(\v w')-\bar{\mf B}^{\dagger c'a}_{2ji}(k^+,-\v w';p^+,-\v z)\,U^{c'b}(\v z)\Big]\notag\\
&\quad\times\Big[U^{ac}(\v w)\,\mf B^{bc}_{2ji}(k^+,-\v w;p^+,-\v z)-U^{cb}(\v z)\,\bar{\mf B}^{ca}_{2ji}(k^+,-\v w;p^+,-\v z)\Big]\Big\rangle\,.
\end{align}

\paragraph{Group III, Row III.I.} The interference of $\mathcal A^{(1)}_{AA}$ with $\mathcal A_B$ is
\begin{align}\label{eq:rowIII_I}
\frac{d^3N^{\rm(III.I)}}{d^3k}=\ &\frac{2}{(2\pi)^3}\int\frac{dp^+}{2\pi}\int_{z,w,w'}e^{-i\v k\cdot(\v w-\v w')}\notag\\
&\times\Big\langle\Big[\bar{\mf A}^{\dagger a}_i(k^+,-\v w')\,\bar{\mf A}^{\dagger c'}_j(p^+,-\v z)\,U^{c'b}(\v z)\notag\\&\qquad\quad-\mf A^{\dagger c'}_i(k^+,-\v w')\,U^{d'b}(\v z)\,\bar{\mf A}^{\dagger d'}_j(p^+,-\v z)\,U^{ac'}(\v w')\Big]\notag\\
&\quad\times\Big[U^{ac}(\v w)\,\mf B^{bc}_{2ji}(k^+,-\v w;p^+,-\v z)-U^{cb}(\v z)\,\bar{\mf B}^{ca}_{2ji}(k^+,-\v w;p^+,-\v z)\Big]\Big\rangle+{\rm c.c.}
\end{align}
In these expressions the conjugated factor stands to the left, at the primed point, and the order of the coefficients is the order in which they act.

\begin{table}[t]
\centering
\renewcommand{\arraystretch}{1.25}
\footnotesize
\begin{tabular}{@{}l l l c l l@{}}
\toprule
Row & product & coefficients & $\#\rho$ & type & range of $p^+$\\
\midrule
I.I & $\mathcal A_{\rm LO}^\dagger\times$ Eq.~\eqref{eq:Nterms} & $\mf A^\dagger\,\mf A\,\mf N_2$ & 4 & virtual & ---\\
I.II & $\mathcal A_{\rm LO}^\dagger\,\mathcal A_{A^{(3)}}$ (one, two $\rho$) & $\mf A^\dagger\,\mf A^{(3)}$ & 2, 3 & virtual & ---\\
I.III & $\mathcal A_{\rm LO}^\dagger\,\mathcal A_{B_3}$ & $\mf A^\dagger\,\bar{\mf B}_3^\dagger\,\mf A$ & 3, 4 & virtual & $\Lambda<p^+<\vee$, $|p^+-k^+|>\Lambda$\\
I.IV & $\mathcal A_{\rm LO}^\dagger\,\mathcal A_{AB_2}$ & $\mf A^\dagger\,\bar{\mf A}^\dagger\,\mf B_2$ & 3, 4 & virtual & $\Lambda<p^+<\vee-k^+$ $[\vee]$\\
I.V & $\mathcal A_{\rm LO}^\dagger\,\mathcal A_{CB_2}$ & $\mf A^\dagger\,\bar{\mf C}_1^\dagger\,\mf B_2$ & 2, 3 & virtual & $\Lambda<p^+<k^+-\Lambda$\\
\midrule
II.I & $|\mathcal A_B|^2$ & $\mf B_2^\dagger\,\mf B_2$ & 2--4 & real & $\Lambda<p^+<\vee-k^+$ $[\vee]$\\
II.II & $|\mathcal A^{(1)}_{AA}|^2$ & $\mf A^\dagger\mf A^\dagger\,\mf A\,\mf A$ & 4 & real & $\Lambda<p^+<\vee$\\
II.III & $|\mathcal A^{(2)}_{AA}|^2$ & $\mf A^\dagger\mf A^\dagger\,\mf A\,\mf A$ & 4 & real & $\Lambda<p^+<\vee$\\
II.IV & $|\mathcal A_{CA}|^2$ & $\mf A^\dagger\,\bar{\mf C}_1^\dagger\,\bar{\mf C}_1\,\mf A$ & 2 & real & $\Lambda<p^+<\vee-k^+$\\
\midrule
III.I & $\mathcal A^{(1)\dagger}_{AA}\,\mathcal A_B$ & $\mf A^\dagger\mf A^\dagger\,\mf B_2$ & 3, 4 & real & $\Lambda<p^+<\vee-k^+$ $[\vee]$\\
III.II & $\mathcal A^{(2)\dagger}_{AA}\,\mathcal A_B$ & $\mf A^\dagger\mf A^\dagger\,\mf B_2$ & 3, 4 & real & $\Lambda<p^+<\vee-k^+$ $[\vee]$\\
III.III & $-\mathcal A^{\dagger}_{CA}\,\mathcal A_B$ & $\mf A^\dagger\,\bar{\mf C}_1^\dagger\,\mf B_2$ & 2, 3 & real & $\Lambda<p^+<\vee-k^+$\\
III.IV & $\mathcal A^{(2)\dagger}_{AA}\,\mathcal A^{(1)}_{AA}$ & $\mf A^\dagger\mf A^\dagger\,\mf A\,\mf A$ & 4 & real & $\Lambda<p^+<\vee$\\
III.V & $-\mathcal A^{\dagger}_{CA}\,\mathcal A^{(1)}_{AA}$ & $\mf A^\dagger\,\bar{\mf C}_1^\dagger\,\mf A\,\mf A$ & 3 & real & $\Lambda<p^+<\vee-k^+$\\
III.VI & $-\mathcal A^{\dagger}_{CA}\,\mathcal A^{(2)}_{AA}$ & $\mf A^\dagger\,\bar{\mf C}_1^\dagger\,\mf A\,\mf A$ & 3 & real & $\Lambda<p^+<\vee-k^+$\\
\bottomrule
\end{tabular}
\caption{The rows of the cross section at order $g^4$. ``Coefficients'' lists the types of coefficients in the product, barred or unbarred. $\#\rho$ is the number of charges before any reordering, with $\mf B_2$ and $\mf B_3$ containing parts with one and with two charges. ``Virtual'': one gluon in the final state; ``real'': two gluons in the final state, the second integrated over. $p^+$ is the longitudinal momentum of the second gluon; where two ranges are given, the first applies to the part of $\mf B_2$ linear in $\rho$ and the bracket $[\vee]$, i.e.\ $\Lambda<p^+<\vee$, to its $\{\rho,\rho\}$ part. Groups~I and~III are to be taken with the complex conjugate. The three-$\rho$ part of $\mathcal A_{A^{(3)}}$ is included in Row~I.I.}
\label{tab:rows}
\end{table}

\paragraph{Number of charges and the four-$\rho$ terms.} The number of charges in a row follows from the coefficients: $\mf A$ and the linear parts of $\mf B_2$, $\mf B_3$ contain one charge, $\mf N_2$ and the $\{\rho,\rho\}$ parts of $\mf B_2$, $\mf B_3$ two, $\mf C_1$ none, and $\mf A^{(3)}$ one, two or three. Terms with four charges, the maximal number at order $g^4$, occur in ten rows: I.I, I.III, I.IV, II.I, II.II, II.III, III.I, III.II and III.IV, and in Row~I.II through the three-$\rho$ part of $\mf A^{(3)}$, which we have included in Row~I.I, Eq.~\eqref{eq:Nterms}. Rows~I.I, II.II, II.III and III.IV consist of four-$\rho$ terms only. Rows~I.V, II.IV and III.III are the only ones with terms quadratic in $\rho$ besides I.II and II.I; Rows~III.V and III.VI are cubic. The ranges of the longitudinal integrals in Table~\ref{tab:rows} follow from the supports of the coefficients, i.e.\ from the kinematic restrictions of Section~\ref{sec:lcpt}.

\NOTE{The assignment of the three-$\rho$ part of $\mathcal A_{A^{(3)}}$ to Row~I.I follows the current version of the paper; the calculation file keeps it in the second row. The sum is the same.}

Eqs.~\eqref{eq:Groups}, with the amplitudes~\eqref{eq:A3parts} and~\eqref{eq:A2parts} and the coefficients of Table~\ref{tab:coeffs}, are the complete cross section at order $g^4$ in terms of the coefficients of $\Omega$. Its evaluation is the subject of the following sections.


\appendix
\sloppy
% ======================================================================
\section{From the light-cone Hamiltonian to its eikonal form}\label{app:hamiltonian}
% ======================================================================

This appendix gives the steps between the light-cone Hamiltonian~\eqref{eq:HLC} and the eikonal Hamiltonian of the soft modes used in the main text. It follows Appendices~A and~B of~\cite{Lublinsky:2016meo}, restricted to pure Yang--Mills theory.

\subsection{Expansion in the coupling and quantisation}

In Eq.~\eqref{eq:PiF} all fields are taken at the same light-cone time. Inserting $D^{ab}_i=\delta^{ab}\partial_i-gf^{abc}A^c_i$ gives $\Pi^a=-\partial_iA^a_i+gf^{abc}\frac{1}{\partial^+}\big(A^c_i\,\partial^+A^b_i\big)$, and $\frac12\Pi^a\Pi^a$ contains the terms
\begin{align}
\frac12(\partial_iA^a_i)^2\,,\qquad -gf^{abc}(\partial_iA^a_i)\frac{1}{\partial^+}(A^b_j\partial^+A^c_j)\,,\qquad \frac{g^2}{2}\Big[f^{abc}\frac{1}{\partial^+}(A^b_i\partial^+A^c_i)\Big]^2\,.
\end{align} The magnetic part $\frac14F^a_{ij}F^a_{ij}$ contains $\frac12(\partial_iA^a_j)^2-\frac12(\partial_iA^a_j)(\partial_jA^a_i)$, the cubic term $-gf^{abc}(\partial_iA^a_j)A^b_iA^c_j$ and the quartic term $\frac{g^2}{4}f^{abc}f^{ade}A^b_iA^c_jA^d_iA^e_j$. After an integration by parts in the transverse coordinates, $(\partial_iA^a_j)(\partial_jA^a_i)\to(\partial_iA^a_i)^2$, the terms quadratic in the field combine into the free part. Expanding in powers of $g$ gives $H=H_0+H_{\rm int}$. The free part is
\begin{align}\label{eq:H0x}
H_0=\int dx^-\,d^2x\ \frac12\,\big(\partial_iA^a_j\big)\big(\partial_iA^a_j\big)\,,
\end{align}
and the interaction is
\begin{align}\label{eq:Hint}
H_{\rm int}=\int dx^-\,d^2x\,\Big[&-gf^{abc}\,\big(\partial_iA^a_j\big)A^b_iA^c_j+\frac{g^2}{4}\,f^{abc}f^{ade}\,A^b_iA^c_jA^d_iA^e_j\notag\\
&-gf^{abc}\,\big(\partial_iA^a_i\big)\frac{1}{\partial^+}\big(A^b_j\,\partial^+A^c_j\big)\notag\\
&+\frac{g^2}{2}\,f^{abc}f^{ade}\,\frac{1}{\partial^+}\big(A^b_i\,\partial^+A^c_i\big)\frac{1}{\partial^+}\big(A^d_j\,\partial^+A^e_j\big)\Big]\,.
\end{align}
The first two terms are the three- and four-gluon vertices that originate from the magnetic part. The third and fourth terms originate from the electric part; they contain the inverse of $\partial^+$ and are the instantaneous interactions characteristic of light-cone quantisation. Eq.~\eqref{eq:Hint} follows from Eqs.~\eqref{eq:HLC} and~\eqref{eq:PiF} by direct expansion, with integrations by parts in the transverse coordinates in the free part.

The theory is quantised by expanding the transverse field at $x^+=0$ in creation and annihilation operators,
\begin{align}\label{eq:modeexp}
A^a_i(x)=\int_0^\infty\frac{dk^+}{2\pi}\int\frac{d^2k}{(2\pi)^2}\,\frac{1}{\sqrt{2k^+}}\Big(a^a_i(k^+,\v k)\,e^{-ik\cdot x}+a^{\dagger a}_i(k^+,\v k)\,e^{ik\cdot x}\Big)\,,
\end{align}
with $k\cdot x=k^+x^--\v k\cdot\v x$ at $x^+=0$ and the commutator~\eqref{eq:aadagger}. Inserted into Eq.~\eqref{eq:H0x} this gives, after normal ordering,
\begin{align}\label{eq:H0}
H_0=\int_0^\infty\frac{dk^+}{2\pi}\int\frac{d^2k}{(2\pi)^2}\,\frac{\v k^2}{2k^+}\,a^{\dagger a}_i(k^+,\v k)\,a^a_i(k^+,\v k)\,.
\end{align}
The eigenstates of $H_0$ are the free Fock states, obtained by acting with creation operators on the perturbative vacuum; the energy of a state is the sum of the energies~\eqref{eq:energy} of its gluons. The interaction $H_{\rm int}$ generates transitions between Fock states. Its form in terms of creation and annihilation operators is needed only after the separation of the modes, which we discuss next.


\subsection{Soft and valence modes; the eikonal terms}

With the decomposition~\eqref{eq:Asplit} and the eikonal approximation described in Section~\ref{sec:BO}, the interaction~\eqref{eq:Hint} reduces to the four terms of Eq.~\eqref{eq:Hdecomp},
\begin{align}\label{eq:Hpieces}
H_g&=-g\int dx^-d^2x\,\big(\partial_i\underline A^a_i\big)\,f^{abc}\,\frac{1}{\partial^+}\big(\overline A^b_j\,\partial^+\overline A^c_j\big)\,,\notag\\
H_{ggg}&=-gf^{abc}\int dx^-d^2x\,\Big[\big(\partial_i\underline A^a_j\big)\underline A^b_i\underline A^c_j+\big(\partial_i\underline A^a_i\big)\frac{1}{\partial^+}\big(\underline A^b_j\,\partial^+\underline A^c_j\big)\Big]\,,\notag\\
H_{gg\text{-inst}}&=g^2f^{abc}f^{ade}\int dx^-d^2x\,\frac{1}{\partial^+}\big(\underline A^b_i\,\partial^+\underline A^c_i\big)\,\frac{1}{\partial^+}\big(\overline A^d_j\,\partial^+\overline A^e_j\big)\,,\notag\\
H_{gggg}&=\frac{g^2}{4}f^{abc}f^{ade}\int dx^-d^2x\,\Big[\underline A^b_i\underline A^c_j\underline A^d_i\underline A^e_j+2\,\frac{1}{\partial^+}\big(\underline A^b_i\,\partial^+\underline A^c_i\big)\frac{1}{\partial^+}\big(\underline A^d_j\,\partial^+\underline A^e_j\big)\Big]\,.
\end{align}
$H_g$ is the coupling of one soft gluon to the valence charge; $H_{ggg}$ is the three-gluon vertex of the soft modes, including its instantaneous part; $H_{gg\text{-inst}}$ is the instantaneous interaction of two soft gluons with the valence charge; $H_{gggg}$ contains the four-gluon vertex of the soft modes and the instantaneous interaction among four soft gluons. $H_{gggg}$ is listed for completeness: it does not enter any of the coefficients needed in this paper (Section~\ref{sec:threeout}).

In terms of the soft creation and annihilation operators the free Hamiltonian of the soft modes and the instantaneous term that contains the valence charge are (the term $H_g$ is given in Eq.~\eqref{eq:Hg})
\begin{align}
H_0&=\int_\Lambda^\vee\frac{dk^+}{2\pi}\int\frac{d^2k}{(2\pi)^2}\,\frac{\v k^2}{2k^+}\,a^{\dagger a}_i(k^+,\v k)\,a^a_i(k^+,\v k)\,,\label{eq:H0soft}\\
H_{gg\text{-inst}}&=\int_\Lambda^\vee\frac{dp^+}{2\pi}\frac{dq^+}{2\pi}\int\frac{d^2p}{(2\pi)^2}\frac{d^2q}{(2\pi)^2}\,\frac{ig^2f^{abc}}{\sqrt{p^+q^+}}
\Bigg(\frac{q^+\,\rho^a(-\v p-\v q)}{2(p^++q^+)^2}\,a^{\dagger b}_i(p)\,a^{\dagger c}_i(q)-{\rm h.c.}\notag\\
&\hspace{6.2cm}-\frac{(p^++q^+)\,\rho^a(-\v p+\v q)}{2(p^+-q^+)^2}\,a^{\dagger b}_i(p)\,a^c_i(q)\Bigg)\,.\label{eq:Hinst}
\end{align}
The first two terms of $H_{gg\text{-inst}}$ create or annihilate a pair of soft gluons. The last term transfers the longitudinal momentum $p^+-q^+$ and the transverse momentum $\v p-\v q$ between a soft gluon and the valence modes. For $H_{ggg}$ we shall need only the matrix element between a one-gluon and a two-gluon state, which is given in Eq.~\eqref{eq:m3} below.

The forms~\eqref{eq:Hpieces} and~\eqref{eq:Hinst} are those of~\cite{Lublinsky:2016meo}. The four-gluon interaction $H_{gggg}$ is not used in this paper (Section~\ref{sec:threeout}); the sign of its instantaneous term is written as it follows from $\frac12\Pi^a\Pi^a$.

% ======================================================================
\section{States, normalisation and matrix elements}\label{app:states}
% ======================================================================

\subsection{States and normalisation}

The soft vacuum $\ket0$ is the state annihilated by all soft annihilation operators, $a^a_i(k^+,\v k)\ket0=0$ for $\Lambda<k^+<\vee$. It is a state of the soft modes only. The complete state of the projectile is a vector in the product of the soft and the valence Hilbert spaces, $\mathcal H_{\rm soft}\otimes\mathcal H_{\rm val}$, and $\ket0$ is to be read as $\ket0\otimes\ket v$ with $\ket v$ the state of the valence modes, which is left unspecified. The soft vacuum is therefore not the vacuum of the theory. Its energy is set to zero, $E_0=0$. The one- and two-gluon states are
\begin{align}\label{eq:states}
\ket{g^a_i(k)}&=\frac{a^{\dagger a}_i(k)}{(2\pi)^{3/2}}\ket0\,,&
\ket{g^a_i(k)\,g^b_j(p)}&=\frac{a^{\dagger a}_i(k)\,a^{\dagger b}_j(p)}{(2\pi)^3}\ket0\,,\notag\\
\ket{g^a_i(k^+,\v z)}&=\frac{a^{\dagger a}_i(k^+,\v z)}{(2\pi)^{1/2}}\ket0\,,&
\ket{g^a_i(k^+,\v z)\,g^b_j(p^+,\v z')}&=\frac{a^{\dagger a}_i(k^+,\v z)\,a^{\dagger b}_j(p^+,\v z')}{2\pi}\ket0\,,
\end{align}
and states with more gluons are defined in the same way, with a factor $(2\pi)^{-3/2}$ for each gluon in momentum space. They are eigenstates of $H_0$ with the energies $E_g(k)$ and $E_{gg}(k,p)$ of Eq.~\eqref{eq:measure}, and they are normalised to $\delta$-functions without factors of $2\pi$,
\begin{align}\label{eq:norms}
\braket{g^b_j(p)|g^a_i(k)}&=\delta^{ab}\delta_{ij}\,\delta^{(3)}(k-p)\,,\notag\\
\braket{g^b_j(p^+,\v z')|g^a_i(k^+,\v z)}&=\delta^{ab}\delta_{ij}\,\delta(k^+-p^+)\,\delta^{(2)}(\v z-\v z')\,,\notag\\
\braket{g^{a'}_{i'}(k')\,g^{b'}_{j'}(p')|g^a_i(k)\,g^b_j(p)}&=\delta^{aa'}\delta_{ii'}\delta^{(3)}(k-k')\,\delta^{bb'}\delta_{jj'}\delta^{(3)}(p-p')\notag\\
&\quad+\delta^{ab'}\delta_{ij'}\delta^{(3)}(k-p')\,\delta^{ba'}\delta_{ji'}\delta^{(3)}(p-k')\,.
\end{align}
The states in transverse coordinate space are related to those in momentum space by
\begin{align}\label{eq:stateFT}
\ket{g^a_i(k^+,\v z)}=\int\frac{d^2k}{2\pi}\,e^{-i\v k\cdot\v z}\ket{g^a_i(k^+,\v k)}\,,\qquad
\ket{g^a_i(k^+,\v k)}=\int\frac{d^2z}{2\pi}\,e^{i\v k\cdot\v z}\ket{g^a_i(k^+,\v z)}\,.
\end{align}
A two-gluon state is symmetric under the exchange of the labels of the two gluons. An expansion $\int d^3p\,d^3q\,F^{bc}_{jk}(p,q)\ket{g^b_j(p)\,g^c_k(q)}$ therefore determines only the part of $F$ that is symmetric under $(b,j,p)\leftrightarrow(c,k,q)$, and the sum over the two-gluon states with unit weight is $\frac12\int d^3p\,d^3q\ket{g^b_j(p)g^c_k(q)}\bra{g^b_j(p)g^c_k(q)}$.


\subsection{Matrix elements}\label{app:matrixelements}

With the states~\eqref{eq:states}, the matrix elements needed in this paper are
\begin{align}
\braket{g^a_i(k)|H_g|0}&=\frac{g\,k^i\rho^a(-\v k)}{4\pi^{3/2}(k^+)^{3/2}}\,,\label{eq:m1}\\
\braket{g^c_l(q)\,g^b_j(p)|H_g|g^a_i(k)}&=\delta^{ac}\delta_{il}\,\delta^{(3)}(k-q)\,\frac{g\,p^j\rho^b(-\v p)}{4\pi^{3/2}(p^+)^{3/2}}
+\delta^{ab}\delta_{ij}\,\delta^{(3)}(k-p)\,\frac{g\,q^l\rho^c(-\v q)}{4\pi^{3/2}(q^+)^{3/2}}\,,\label{eq:m2}\\
\braket{g^c_l(q)\,g^b_j(p)|H_{ggg}|g^a_i(k)}&=\frac{igf^{abc}\,\delta^{(3)}(k-p-q)}{8\pi^{3/2}\sqrt{k^+p^+q^+}}\Bigg[\Big(p^i-q^i+\frac{q^+-p^+}{k^+}k^i\Big)\delta_{jl}\notag\\
&\hspace{0.6cm}+\Big(k^j+q^j-\frac{k^++q^+}{p^+}p^j\Big)\delta_{il}
+\Big(\frac{k^++p^+}{q^+}q^l-p^l-k^l\Big)\delta_{ij}\Bigg]\,,\label{eq:m3}\\
\braket{g^c_j(q)\,g^b_i(p)|H_{gg\text{-inst}}|0}&
=\frac{ig^2f^{abc}\,(q^+-p^+)\,\delta_{ij}\,\rho^a(-\v p-\v q)}{2(2\pi)^3\sqrt{p^+q^+}\,(p^++q^+)^2}\,,\label{eq:m4}\\
\braket{g^c_j(q)|H_{gg\text{-inst}}|g^b_i(p)}&
=\frac{ig^2f^{abc}\,(p^++q^+)\,\delta_{ij}\,\rho^a(\v p-\v q)}{2(2\pi)^3\sqrt{p^+q^+}\,(p^+-q^+)^2}\,.\label{eq:m5}
\end{align}
Eq.~\eqref{eq:m1} is the emission of a soft gluon by the valence charge; Eq.~\eqref{eq:m2} the emission of a second gluon in the presence of a first one, which is a spectator; Eq.~\eqref{eq:m3} the splitting of a soft gluon into two; Eq.~\eqref{eq:m4} the instantaneous emission of two gluons; Eq.~\eqref{eq:m5} the instantaneous interaction of a soft gluon with the valence charge. The matrix elements of the opposite transitions follow by Hermitian conjugation, with $[\rho^a(\v p)]^\dagger=\rho^a(-\v p)$; for instance $\braket{0|H_g|g^a_i(k)}=g\,k^i\rho^a(\v k)/[4\pi^{3/2}(k^+)^{3/2}]$. Eqs.~\eqref{eq:m1}, \eqref{eq:m2}, \eqref{eq:m4} and~\eqref{eq:m5} follow directly from Eqs.~\eqref{eq:Hg} and~\eqref{eq:Hinst}; all of them are operators on $\mathcal H_{\rm val}$ except Eq.~\eqref{eq:m3}.

As an illustration, Eq.~\eqref{eq:m1} follows from Eq.~\eqref{eq:Hg} with $\bra{g^a_i(k)}a^{\dagger b}_j(p)\ket0=(2\pi)^{3/2}\delta^{ab}\delta_{ij}\delta^{(3)}(k-p)$: the factor is $(2\pi)^{3/2}/(2\pi)^3\times g/\sqrt2=g/(4\pi^{3/2})$. In Eq.~\eqref{eq:m2} the gluon that is not emitted is a spectator, and the two terms correspond to the two ways of identifying it with one of the outgoing gluons. Eqs.~\eqref{eq:m4} and~\eqref{eq:m5} follow in the same way from the first and the last term of Eq.~\eqref{eq:Hinst}; in Eq.~\eqref{eq:m4} the two contractions of the pair of creation operators combine into the factor $q^+-p^+$ by the antisymmetry of $f^{abc}$. Eq.~\eqref{eq:m3} is taken from~\cite{Lublinsky:2016meo}.

% ======================================================================
\section{Perturbation theory with non-commuting matrix elements}\label{app:pt}
% ======================================================================

In this appendix we give the perturbative series for the dressed states when the matrix elements $V_{nm}=\braket{n|H_{\rm int}|m}$ are operators on $\mathcal H_{\rm val}$, and describe the numerical test of the order of the factors. The series was given in Appendix~E of~\cite{Lublinsky:2016meo}.

\subsection{The series}

The derivation is that of Rayleigh--Schr\"odinger perturbation theory. One writes $\ket{\Psi_i}=\ket i\,Z_i+\sum_{n\ne i}\ket n\,c_{ni}$ with operator-valued $Z_i$ and $c_{ni}$ and requires that $H\ket{\Psi_i}$ has no component along the states $\ket{\Psi_j}$ with $E_j\neq E_i$. Since the energies are numbers, the equations for the $c_{ni}$ are solved order by order as for a number-valued perturbation, with each new insertion of $V$ multiplying from the left. The factor $Z_i$ is fixed by the orthonormality of the dressed states, $\braket{\Psi_j|\Psi_i}=\delta_{ji}$, and by the condition that it is Hermitian. To second order this gives Eq.~\eqref{eq:pt2}. For the soft vacuum, $\ket i=\ket0$ with $E_0=0$, we shall also quote results of third order, for which
\begin{align}\label{eq:pt3}
\ket{\Psi_0}=\ &\ket0-\sum_n\ket n\frac{V_{n0}}{E_n}+\sum_{n,m}\ket n\frac{V_{nm}V_{m0}}{E_nE_m}
-\sum_{n,m,l}\ket n\frac{V_{nm}V_{ml}V_{l0}}{E_nE_mE_l}
\notag\\
&+\sum_{n,m}\ket n\frac{V_{n0}\,V_{0m}V_{m0}}{E_n^2E_m}-\frac12\sum_m\frac{V_{0m}V_{m0}}{E_m^2}\ket0
+\frac12\sum_{n,m}\ket n\frac{V_{n0}\,V_{0m}V_{m0}}{E_nE_m^2}+\dots,
\end{align}
with all sums over states different from $\ket0$. The components along the states $\ket n\neq\ket0$ are given to third order. The component along $\ket0$ is given to second order; its third-order term, $\propto V_{0n}V_{nm}V_{m0}$, is of order $g^4$ for the interactions of Section~\ref{sec:BO} and is not needed. Eqs.~\eqref{eq:pt2} and~\eqref{eq:pt3} hold for operator-valued $V_{nm}$ in the order written. Two properties define the states uniquely. First, $H$ does not connect $\ket{\Psi_i}$ and $\ket{\Psi_j}$ for $E_i\neq E_j$; more precisely, $\braket{\Psi_j|H|\Psi_i}$ is an operator on $\mathcal H_{\rm val}$ which vanishes for $E_i\ne E_j$ to the order considered. Second, the states are orthonormal, $\braket{\Psi_j|\Psi_i}=\delta_{ji}$ as operators on $\mathcal H_{\rm val}$, with the component along the initial state Hermitian (canonical normalisation). The second property fixes the order of the factors in the normalisation terms: in the last term of each line of Eq.~\eqref{eq:pt3} the single matrix element $V_{n0}$ stands to the left of the norm factor $V_{0m}V_{m0}$. We have verified Eqs.~\eqref{eq:pt2} and~\eqref{eq:pt3}, including the order of the factors, numerically: for a model with a few soft levels and matrices in place of the $V_{nm}$, the unitary transformation with the two properties above was constructed exactly and compared with the series (Section~\ref{sec:unitarity}).


\subsection{Numerical test}\label{app:pt_num}

The series was tested in a model with $L=5$ soft levels with energies $E_n=(0,1,2.3,3.1,4.7)$ and a valence space of dimension $d=3$, so that the $V_{nm}$ are $3\times3$ blocks of a random Hermitian $15\times15$ matrix (with $V_{00}=0$, and all diagonal blocks zero in the test of Eq.~\eqref{eq:pt2}). For $H=H_0+gV$ the exact unitary $\Omega$ with the two properties stated above was constructed from the spectral projectors $Q_n$ of $H$ on the group of eigenvalues that evolves from $E_n$: $\Omega P_n=Q_nP_n(P_nQ_nP_n)^{-1/2}$, with $P_n$ the projector on the level $n$ of $H_0$. The blocks $\braket{n|\Omega|i}$ were compared with the series. The results are (script \texttt{scripts/check\_pt\_ordering.py}):
\begin{itemize}
\item Eq.~\eqref{eq:pt2} for the level $i=2$: the maximal deviation is $2.4\times10^{-3}$, $3.0\times10^{-4}$, $3.7\times10^{-5}$ for $g=0.02,\,0.01,\,0.005$, i.e.\ it scales as $g^3$.
\item Eq.~\eqref{eq:pt3}, components along $\ket n\ne\ket0$: the maximal deviation is $7.6\times10^{-5}$, $4.3\times10^{-6}$, $2.5\times10^{-7}$ for the same couplings, i.e.\ it scales as $g^4$.
\item Eq.~\eqref{eq:pt3} with the norm factor $V_{0m}V_{m0}$ to the \emph{left} of $V_{n0}$ in the two normalisation terms: the deviation is $4.6\times10^{-4}$, $6.2\times10^{-5}$, $8.0\times10^{-6}$, i.e.\ of order $g^3$. This order of the factors is wrong.
\end{itemize}

% ======================================================================
\section{Details of the single-gluon incoming state}\label{app:onegluon}
% ======================================================================

This appendix contains the intermediate steps of Section~\ref{sec:onegin}: the Fourier transforms, the explicit forms of the terms with the three-gluon vertex, the cancellation of the double pole, the algebra of the terms with two valence vertices, the coordinate-space forms of all components, and the verification of the normalisation. The organisation follows that of the supplementary Appendix~G of~\cite{Lublinsky:2016meo} for the vacuum initial state.

\subsection{Fourier transforms}\label{app:fourier}

For the coordinate-space forms we define the state of an incoming gluon at the transverse position $\v z'$ by Eq.~\eqref{eq:stateFT},
\begin{align}
\ket{\Psi^a_i(k^+,\v z')}\equiv\int\frac{d^2k}{2\pi}\,e^{-i\v k\cdot\v z'}\ket{\Psi^a_i(k^+,\v k)}\,,
\end{align}
and use the abbreviations
\begin{align}\label{eq:XY}
\v X\equiv\v x-\v z\,,\qquad \v X'\equiv\v x-\v z'\,,\qquad \v Y\equiv\v y-\v z\,,
\end{align}
where $\v x$, $\v y$ are the positions of the charges and $\v z$ is the position of an outgoing gluon.
The Fourier transforms that are needed are, for real $\xi>0$ and real $a,b$,
\begin{align}
&\int\frac{d^2k}{2\pi}\,\frac{k^i}{\v k^2}\,e^{i\v k\cdot\v X}=\frac{iX^i}{\v X^2}\,,\label{eq:FT_K}\\
&\int\frac{d^2k}{2\pi}\frac{d^2p}{2\pi}\,\frac{k^ip^j}{\v k^2\,(a\v k^2+b\v p^2)}\,e^{i\v k\cdot\v X+i\v p\cdot\v Z}=-\frac{X^iZ^j}{\v Z^2\,(b\v X^2+a\v Z^2)}\,,\label{eq:FT_i2}\\
&\int\frac{d^2k}{2\pi}\frac{d^2p}{2\pi}\,\frac{1}{\xi\v k^2-\v p^2}\,e^{i\v k\cdot\v X+i\v p\cdot\v Z}=\frac{1}{\xi\v Z^2-\v X^2}\,,\label{eq:FT_i3}\\
&\int\frac{d^2k}{2\pi}\frac{d^2p}{2\pi}\,\frac{k^m-p^m}{(\v k-\v p)^2}\,\frac{p^l-\xi k^l}{\v p^2-\xi\v k^2}\,e^{i\v k\cdot\v X+i\v p\cdot\v Z}=-\frac{X^m+\xi Z^m}{\v X^2-\xi\v Z^2}\,\frac{X^l+Z^l}{(\v X+\v Z)^2}\,,\label{eq:FT_i4}\\
&\int\frac{d^2k}{2\pi}\frac{d^2p}{2\pi}\,\frac{k^m-p^m}{\v p^2-\xi\v k^2}\,\frac{p^l-\xi k^l}{(\v p-\xi\v k)^2}\,e^{i\v k\cdot\v X+i\v p\cdot\v Z}=-\frac{X^l+Z^l}{\v X^2-\xi\v Z^2}\,\frac{X^m+\xi Z^m}{(\v X+\xi\v Z)^2}\,,\label{eq:FT_i5}
\end{align}
with principal values where a denominator vanishes~\cite{Lublinsky:2016meo,Radhakrishnan:2026soft}. In all cases below they are applied with $\v X\to\v X'$ and $\v Z\to-\v X$ or $\v Z\to-\v Y$, because the incoming gluon contributes the phase $e^{-i\v k\cdot\v z'}$, the outgoing gluon the phase $e^{i\v p\cdot\v z}$ and a charge $\rho(\v q)=\int_xe^{i\v q\cdot\v x}\rho(\v x)$.

Eq.~\eqref{eq:FT_K} follows from $\int_0^{2\pi}d\theta\,\cos\theta\,e^{ikX\cos\theta}=2\pi iJ_1(kX)$ and $\int_0^\infty dt\,J_1(t)=1$. Eqs.~\eqref{eq:FT_i2} and~\eqref{eq:FT_i3} follow from Eq.~\eqref{eq:FT_K} and from the transform of $1/(a\v k^2+b\v p^2)$ in four dimensions.

\NOTE{Eqs.~\eqref{eq:FT_i4} and~\eqref{eq:FT_i5} are quoted from the sources and were not re-derived here.}

\subsection{Vacuum outgoing state}

For Eq.~\eqref{eq:psi_rho}: $\braket{0|H_g|g^a_i(k)}=g\,k^i\rho^a(\v k)/[4\pi^{3/2}(k^+)^{3/2}]$ and $E_g(k)=\v k^2/2k^+$. In coordinate space, with Eq.~\eqref{eq:FT_K},
\begin{align}\label{eq:psi_rhoX}
\ket{(\Psi^a_i(k^+,\v z'))_\rho}=\frac{ig}{2\pi^{3/2}\sqrt{k^+}}\int_x\frac{X'^i}{\v X'^2}\,\rho^a(\v x)\ket0\,.
\end{align}

\subsection{One outgoing gluon: instantaneous term}

The matrix element in Eq.~\eqref{eq:psi_inst_def} is Eq.~\eqref{eq:m5} with $(p,b,i)\to(k,a,i)$ and $(q,c,j)\to(p,b,j)$, i.e.\ $ig^2f^{abc}(k^++p^+)\delta_{ij}\rho^c(\v k-\v p)/[2(2\pi)^3\sqrt{k^+p^+}(k^+-p^+)^2]$ after renaming the colour index of the charge. Dividing by $E_g(k)-E_g(p)$ gives Eq.~\eqref{eq:psi_inst}. With the momentum fraction $\xi\equiv p^+/k^+$ and Eq.~\eqref{eq:FT_i3},
\begin{align}\label{eq:psi_instX}
\ket{(\Psi^a_i(k^+,\v z'))_{g\rho\,{\rm inst}}}=\frac{ig^2f^{abc}}{(2\pi)^3}\int d\xi\int_{x,z}\frac{\sqrt\xi\,(1+\xi)}{(1-\xi)^2}\,\frac{\rho^c(\v x)}{\xi\v X^2-\v X'^2}\,\ket{g^b_i(\xi k^+,\v z)}\,,
\end{align}
where the $\xi$ integral runs over $1+\Lambda/k^+<\xi<\vee/k^+$ in region~1 and over $\Lambda/k^+<\xi<1-\Lambda/k^+$ in region~2. Note that $\xi>1$ in region~1.


\subsection{One outgoing gluon: terms with the three-gluon vertex}

In Eq.~\eqref{eq:psi1_def} the matrix element of $H_g$, Eq.~\eqref{eq:m2}, sets one of the intermediate gluons equal to the incoming one, $q=k$, $(c,l)=(a,i)$, and the other to the emitted gluon, $r=p-k$; the conjugate of Eq.~\eqref{eq:m3} with the parent $(p,b,j)$ is $-igf^{bda}\,[\dots]/[8\pi^{3/2}\sqrt{p^+r^+k^+}]$, and the denominators are $-E_g(r)$ and $E_g(k)-E_g(p)$. Inserting the matrix elements we find
\begin{align}\label{eq:psi1}
\ket{(\Psi^a_i(k))^1_{g\rho}}=\ &\frac{ig^2f^{abc}}{8\pi^3}\int_{k^++\Lambda}^{\vee}dp^+\!\int d^2p\ \frac{\rho^c(\v k-\v p)\,(p-k)^l}{\sqrt{p^+k^+}\,(p^+-k^+)^2}\,
\frac{1}{\Big(\frac{\v k^2}{k^+}-\frac{\v p^2}{p^+}\Big)\frac{(\v k-\v p)^2}{k^+-p^+}}\notag\\
&\times\Bigg[\Big(2k^j-p^j+\frac{p^+-2k^+}{p^+}\,p^j\Big)\delta_{il}+\Big(2p^i-k^i-\frac{2p^+-k^+}{k^+}\,k^i\Big)\delta_{jl}\notag\\
&\qquad+\Big(\frac{p^++k^+}{p^+-k^+}\,(p-k)^l-k^l-p^l\Big)\delta_{ij}\Bigg]\ket{g^b_j(p)}\,.
\end{align}
The three brackets are $-\frac2\xi(p-\xi k)^j$, $2(p-\xi k)^i$ and $-\frac{2}{1-\xi}(p-\xi k)^l$ with $\xi=p^+/k^+$, so that
Eq.~\eqref{eq:psi1_xi} follows. With Eq.~\eqref{eq:FT_i4},
\begin{align}\label{eq:psi1X}
\ket{(\Psi^a_i(k^+,\v z'))^1_{g\rho}}=\ &-\frac{ig^2f^{abc}}{4\pi^3}\int_{1+\Lambda/k^+}^{\vee/k^+}d\xi\int_{x,z}\frac{\sqrt\xi}{1-\xi}\,\frac{X'^m-\xi X^m}{\v X'^2-\xi\v X^2}\,\frac{X'^l-X^l}{(\v X'-\v X)^2}\,\rho^c(\v x)\notag\\
&\times\Big(\delta_{il}\delta_{jm}-\frac1\xi\,\delta_{jl}\delta_{im}-\frac{1}{1-\xi}\,\delta_{ij}\delta_{lm}\Big)\ket{g^b_j(\xi k^+,\v z)}\,.
\end{align}
In Eq.~\eqref{eq:psi2_def} the matrix element of $H_{ggg}$ is Eq.~\eqref{eq:m3} with the outgoing gluons $p$ and $r=k-p$, and the matrix element of $H_g$ is the conjugate of Eq.~\eqref{eq:m2}, which absorbs the gluon $r$ with the factor $g\,r^m\rho^d(\v r)/[4\pi^{3/2}(r^+)^{3/2}]$. The result is
\begin{align}\label{eq:psi2}
\ket{(\Psi^a_i(k))^2_{g\rho}}=\ &\frac{ig^2f^{abc}}{8\pi^3}\int_{\Lambda}^{k^+-\Lambda}dp^+\!\int d^2p\ \frac{\rho^c(\v k-\v p)\,(k-p)^l}{\sqrt{k^+p^+}\,(k^+-p^+)^2}\notag\\
&\times\frac{1}{\Big(\frac{\v k^2}{k^+}-\frac{(\v k-\v p)^2}{k^+-p^+}-\frac{\v p^2}{p^+}\Big)\Big(\frac{\v k^2}{k^+}-\frac{\v p^2}{p^+}\Big)}\notag\\
&\times\Bigg[\Big(2p^i-k^i+\frac{k^+-2p^+}{k^+}\,k^i\Big)\delta_{jl}+\Big(2k^j-p^j-\frac{2k^+-p^+}{p^+}\,p^j\Big)\delta_{il}\notag\\
&\qquad+\Big(\frac{k^++p^+}{k^+-p^+}\,(k-p)^l-p^l-k^l\Big)\delta_{ij}\Bigg]\ket{g^b_j(p)}\,,
\end{align}
or, with $\frac{\v k^2}{k^+}-\frac{(\v k-\v p)^2}{k^+-p^+}-\frac{\v p^2}{p^+}=-\frac{(\v p-\xi\v k)^2}{\xi(1-\xi)k^+}$,
Eq.~\eqref{eq:psi2_xi} follows. In coordinate space, with Eq.~\eqref{eq:FT_i5},
\begin{align}\label{eq:psi2X}
\ket{(\Psi^a_i(k^+,\v z'))^2_{g\rho}}=\ &-\frac{ig^2f^{abc}}{4\pi^3}\int_{\Lambda/k^+}^{1-\Lambda/k^+}d\xi\ \frac{\xi^{3/2}}{1-\xi}\int_{x,z}\frac{X'^l-X^l}{\v X'^2-\xi\v X^2}\,\frac{X'^m-\xi X^m}{(\v X'-\xi\v X)^2}\,\rho^c(\v x)\notag\\
&\times\Big(\delta_{il}\delta_{jm}-\frac1\xi\,\delta_{jl}\delta_{im}-\frac{1}{1-\xi}\,\delta_{lm}\delta_{ij}\Big)\ket{g^b_j(\xi k^+,\v z)}\,.
\end{align}

\subsection{Cancellation of the double pole}\label{app:doublepole}

Each of the four terms linear in $\rho$ has a double pole at $p^+=k^+$, i.e.\ a factor $1/(1-\xi)^2$. The pole lies outside the regions~\eqref{eq:ranges}, but the limits of integration approach it from both sides at the distance $\Lambda$, and a single term gives a contribution proportional to $k^+/\Lambda$. In the sum of the instantaneous and the non-instantaneous terms the double pole cancels in each region separately: for $p^+\to k^+$ the last bracket of Eq.~\eqref{eq:psi1} is dominated by its third term, and one finds that the integrand of Eq.~\eqref{eq:psi1} approaches minus the integrand of Eq.~\eqref{eq:psi_inst}; the same holds for Eq.~\eqref{eq:psi2}. No term proportional to $k^+/\Lambda$ remains.

Numerically, for a generic configuration of the transverse momenta, $(p^+-k^+)^2$ times the sum of the instantaneous and the non-instantaneous term of Eq.~\eqref{eq:B3mom} is $1.7\times10^{-3}$, $1.7\times10^{-4}$, $1.7\times10^{-5}$ for $|p^+-k^+|/k^+=10^{-2},10^{-3},10^{-4}$ on either side, while the instantaneous term alone gives $1.46$ (script \texttt{scripts/check\_onegluon\_state.py}). A simple pole remains; it is odd in $p^+-k^+$ and gives no term $\propto k^+/\Lambda$.

\subsection{Terms with two valence vertices}\label{app:tworho}

In Eq.~\eqref{eq:psi1rr} the matrix elements are the conjugate of Eq.~\eqref{eq:m1} for the absorption and Eq.~\eqref{eq:m1} for the emission, and $E_g(k)\,(E_g(k)-E_g(p))=\v k^2(p^+\v k^2-k^+\v p^2)/(4k^{+2}p^+)$. In Eq.~\eqref{eq:psi2rr} the two terms of Eq.~\eqref{eq:m2} and of its conjugate give two equal connected contributions, which cancel the factor $\frac12$, and the disconnected one discussed in the main text. The sum of the two contributions is
\begin{align}
\frac{g^2\,k^ip^j}{4\pi^3\sqrt{k^+p^+}\,\v k^2\v p^2}\ \frac{k^+\v p^2\,\rho^b(-\v p)\rho^a(\v k)-p^+\v k^2\,\rho^a(\v k)\rho^b(-\v p)}{p^+\v k^2-k^+\v p^2}\,.
\end{align}
With the decomposition~\eqref{eq:decomp}, $\rho^a(\v k)\rho^b(-\v p)=\frac12\{\rho^a(\v k),\rho^b(-\v p)\}+\frac i2f^{abc}\rho^c(\v k-\v p)$ and the opposite sign of the last term for the opposite order, the energy denominator cancels in the anticommutator part, which gives Eqs.~\eqref{eq:rrsum}--\eqref{eq:psi_comm}.
In coordinate space, with Eqs.~\eqref{eq:FT_K} and~\eqref{eq:FT_i2} and $\xi=p^+/k^+$,
\begin{align}
\ket{(\Psi^a_i(k^+,\v z'))_{g\{\rho,\rho\}}}&=-\frac{g^2}{8\pi^3}\int_{\Lambda/k^+}^{\vee/k^+}\frac{d\xi}{\sqrt\xi}\int_{x,y,z}\frac{X'^i\,Y^j}{\v X'^2\,\v Y^2}\,\big\{\rho^a(\v x),\rho^b(\v y)\big\}\ket{g^b_j(\xi k^+,\v z)}\,,\label{eq:psi_antiX}\\
\ket{(\Psi^a_i(k^+,\v z'))_{g[\rho,\rho]}}&=-\frac{ig^2f^{abc}}{8\pi^3}\int_{\Lambda/k^+}^{\vee/k^+}\frac{d\xi}{\sqrt\xi}\int_{x,z}\frac{X'^i\,X^j}{\v X'^2\,\v X^2}\,\frac{\v X'^2+\xi\v X^2}{\xi\v X^2-\v X'^2}\,\rho^c(\v x)\ket{g^b_j(\xi k^+,\v z)}\,.\label{eq:psi_commX}
\end{align}


\subsection{Two outgoing gluons}

In coordinate space
\begin{align}\label{eq:psiggrhoX}
\ket{(\Psi^a_i(k^+,\v z'))_{gg\rho}}=\frac{ig}{2\pi^{3/2}}\int_\Lambda^\vee\frac{dp^+}{\sqrt{p^+}}\int_{x,z}\frac{X^j}{\v X^2}\,\rho^b(\v x)\,\ket{g^a_i(k^+,\v z')\,g^b_j(p^+,\v z)}\,.
\end{align}
For the $\rho$-independent component~\eqref{eq:psigg} the coordinate-space form is contained in the coefficient $\mf C_1$, Eq.~\eqref{eq:C1X}.

\subsection{Normalisation of the dressed one-gluon state}\label{app:norm}

We verify Eq.~\eqref{eq:normcond} to order $g^2$. The terms of order $g^2$ in the overlap are the overlaps of two components of order $g$ and the overlaps of $\ket g$ with the components of order $g^2$.

\paragraph{Terms with two charges.} From Eqs.~\eqref{eq:psi_rho} and~\eqref{eq:psiggrho}, with the normalisation~\eqref{eq:norms} of the two-gluon states,
\begin{align}
\braket{(\Psi^{a'}_{i'}(k'))_\rho|(\Psi^a_i(k))_\rho}&=\frac{g^2}{4\pi^3\sqrt{k^+k'^+}}\,\frac{k^ik'^{i'}}{\v k^2\v k'^2}\,\rho^{a'}(-\v k')\,\rho^a(\v k)\,,\label{eq:ov1}\\
\braket{(\Psi^{a'}_{i'}(k'))_{gg\rho}|(\Psi^a_i(k))_{gg\rho}}&=\frac{g^2}{4\pi^3\sqrt{k^+k'^+}}\,\frac{k^ik'^{i'}}{\v k^2\v k'^2}\,\rho^a(\v k)\,\rho^{a'}(-\v k')\notag\\
&\quad+\delta^{aa'}\delta_{ii'}\delta^{(3)}(k-k')\,\frac{g^2}{4\pi^3}\int_\Lambda^\vee\frac{dp^+}{p^+}\int d^2p\,\frac{\rho^b(\v p)\,\rho^b(-\v p)}{\v p^2}\,,\label{eq:ov2}
\end{align}
where the first term of Eq.~\eqref{eq:ov2} comes from the exchange term of Eq.~\eqref{eq:norms}. From Eq.~\eqref{eq:psi_anti},
\begin{align}\label{eq:ov3}
\braket{g^{a'}_{i'}(k')|(\Psi^a_i(k))_{g\{\rho,\rho\}}}+\braket{(\Psi^{a'}_{i'}(k'))_{g\{\rho,\rho\}}|g^a_i(k)}=-\frac{g^2}{4\pi^3\sqrt{k^+k'^+}}\,\frac{k^ik'^{i'}}{\v k^2\v k'^2}\,\big\{\rho^a(\v k),\rho^{a'}(-\v k')\big\}\,.
\end{align}
The terms that are not proportional to $\delta^{(3)}(k-k')$ cancel in the sum of Eqs.~\eqref{eq:ov1}--\eqref{eq:ov3}, for non-commuting charges. The remaining term of Eq.~\eqref{eq:ov2} is $-2\mf N_2\,\delta^{aa'}\delta_{ii'}\delta^{(3)}(k-k')$ and is cancelled by the term $2\mf N_2$ from the first term of Eq.~\eqref{eq:psi_sum}, with $\mf N_2$ of Eq.~\eqref{eq:N2mom}. The coefficient of the identity in $\Omega$ is thus the same in the vacuum and in the one-gluon sector, as it must be.

\paragraph{Terms linear in $\rho$.} Consider a component of the form $\int d^3p\,\ket{g^b_j(p)}\,M^{ba}_{ji}(k,p)/(E_g(k)-E_g(p))$. Its overlap with $\bra{g^{a'}_{i'}(k')}$, plus the conjugate term, is
\begin{align}
\frac{M^{a'a}_{i'i}(k,k')-\big[M^{aa'}_{ii'}(k',k)\big]^\dagger}{E_g(k)-E_g(k')}\,.
\end{align}
It vanishes if $M$ is the matrix element of a Hermitian operator; this is the case for the instantaneous terms~\eqref{eq:psi_inst}. The commutator term~\eqref{eq:psi_comm} cancels against its conjugate in the same way, because its kernel is odd under the exchange of the two gluons. The terms with the three-gluon vertex, Eqs.~\eqref{eq:psi1} and~\eqref{eq:psi2}, do not cancel by themselves; they cancel against the overlaps of the two-gluon components $\ket{(\Psi)_{gg}}$ and $\ket{(\Psi)_{gg\rho}}$,
\begin{align}\label{eq:ov4}
&\braket{g^{a'}_{i'}(k')|(\Psi^a_i(k))^1_{g\rho}}+\braket{g^{a'}_{i'}(k')|(\Psi^a_i(k))^2_{g\rho}}\notag\\
&\quad+\braket{(\Psi^{a'}_{i'}(k'))^1_{g\rho}|g^a_i(k)}+\braket{(\Psi^{a'}_{i'}(k'))^2_{g\rho}|g^a_i(k)}\notag\\
&\quad+\braket{(\Psi^{a'}_{i'}(k'))_{gg\rho}|(\Psi^a_i(k))_{gg}}+\braket{(\Psi^{a'}_{i'}(k'))_{gg}|(\Psi^a_i(k))_{gg\rho}}=0\,.
\end{align}
Only the sum of the six terms vanishes. Eq.~\eqref{eq:ov4} is the triple-vertex part of the unitarity relation~\eqref{eq:B3unit} below, where it is derived; we have verified it numerically with Eqs.~\eqref{eq:B3mom}, \eqref{eq:C1mom} and~\eqref{eq:A1mom}.

\paragraph{Terms independent of $\rho$.} From Eq.~\eqref{eq:psigg}, summing over colours and polarisations,
\begin{align}\label{eq:ov5}
\braket{(\Psi^{a'}_{i'}(k'))_{gg}|(\Psi^a_i(k))_{gg}}&=\delta^{aa'}\delta_{ii'}\delta^{(3)}(k-k')\,\frac{g^2N_c}{8\pi^3}\int\frac{d^2\tilde p}{\tilde{\v p}^2}\notag\\
&\quad\times\int_{\Lambda/k^+}^{1-\Lambda/k^+}d\xi\,\Big[\xi(1-\xi)+\frac{1-\xi}{\xi}+\frac{\xi}{1-\xi}\Big]\notag\\
&=\delta^{aa'}\delta_{ii'}\delta^{(3)}(k-k')\,\frac{g^2N_c}{8\pi^3}\,\Big[2\log\frac{k^+}{\Lambda}-\frac{11}{6}\Big]\int\frac{d^2\tilde p}{\tilde{\v p}^2}\,,
\end{align}
with $\tilde{\v p}=\v p-\xi\v k$ and up to terms that vanish for $\Lambda/k^+\to0$. This is the probability for the incoming gluon to split. It has to be compensated by a $\rho$-independent term in the component along $\ket{g^a_i(k)}$. Such a term is not part of $\mf N$, which is fixed by the vacuum sector; it is a contribution to $\mf B_3$ proportional to $\delta^{(3)}(k-p)$,
\begin{align}\label{eq:B3rho0}
\Big[\mf B^{ba}_{3ji}(k,p)\Big]_{\rho^0}=-\frac12\,(2\pi)^3\,\delta^{ab}\delta_{ij}\,\delta^{(3)}(k-p)\,\frac{g^2N_c}{8\pi^3}\,\Big[2\log\frac{k^+}{\Lambda}-\frac{11}{6}\Big]\int\frac{d^2\tilde p}{\tilde{\v p}^2}\,,
\end{align}
which is the second term of Eq.~\eqref{eq:psi_sum} and the term $-\frac12\sum_mV_{im}V_{mi}/(E_i-E_m)^2$ of Eq.~\eqref{eq:pt2} with two-gluon intermediate states. It is the wave-function renormalisation of the soft gluon. The transverse integral $\int d^2\tilde p/\tilde{\v p}^2$ has no scale. It diverges logarithmically at small and at large $\tilde{\v p}$, and it vanishes in dimensional regularisation, where the two divergences cancel each other. With this rule Eq.~\eqref{eq:B3rho0} is zero, the $\rho$-independent term is absent from $\mf B_3$, Eq.~\eqref{eq:B3mom}, and
\begin{align}\label{eq:Nresult}
\mf N=1-\frac{g^2}{8\pi^3}\,\log\frac{\vee}{\Lambda}\int d^2p\,\frac{\rho^a(-\v p)\,\rho^a(\v p)}{\v p^2}
\end{align}
is the complete coefficient of $\ket{g^a_i(k)}$ at order $g^2$.

% ======================================================================
\section{Unitarity relations: derivation}\label{app:unitarity}
% ======================================================================

In this appendix the relations among the coefficients of $\Omega$ that follow from $\Omega^\dagger\Omega=1$ are derived order by order in $g$.

The conjugate of Eq.~\eqref{eq:OmegaExp} is
\begin{align}\label{eq:OmegaDag}
\Omega^\dagger = \ &\mf{N}^\dagger
+\int\frac{d^3p}{(2\pi)^3}\Big[\mf{A}^{\dagger b}_{1j}(p)\,a^{\dagger b}_j(p)+\mf{A}^{\dagger b}_{2j}(p)\,a^{b}_j(p)\Big]
\notag\\&
+\int\frac{d^3p}{(2\pi)^3}\frac{d^3q}{(2\pi)^3}\Big[\mf{B}^{\dagger cb}_{1kj}(p,q)\,a^{\dagger b}_j(p)a^{\dagger c}_k(q)+\mf{B}^{\dagger cb}_{2kj}(p,q)\,a^{b}_j(p)a^{c}_k(q)
\notag\\&\hspace{3.4cm}+\mf{B}^{\dagger cb}_{3kj}(p,q)\,a^{\dagger b}_j(p)a^{c}_k(q)\Big]
\notag\\&
+\int\frac{d^3p}{(2\pi)^3}\frac{d^3q}{(2\pi)^3}\Big[\mf{C}^{\dagger dcb}_{1lkj}(p,q)\,a^{\dagger b}_j(p)a^{c}_k(p-q)a^{d}_l(q)
\notag\\&\hspace{3.4cm}+\mf{C}^{\dagger dcb}_{2lkj}(p,q)\,a^{\dagger d}_l(q)a^{\dagger c}_k(p-q)a^{b}_j(p)\Big]+\dots,
\end{align}
where the dagger on a coefficient is the Hermitian conjugate on $\mathcal H_{\rm val}$ and the indices and arguments are those of the coefficient that is conjugated. The product $\Omega^\dagger\Omega$ is brought to normal order with the commutator~\eqref{eq:aadagger}. The coefficients commute with $a$ and $a^\dagger$, but not with each other: in each product the coefficient that comes from $\Omega^\dagger$ stands to the left of the one from $\Omega$. The normal-ordered products of $a^\dagger$ and $a$ are linearly independent, so that the coefficient of each of them has to vanish, except for the identity. We give the relations in the order in which they arise; the coefficients of the conjugate structures give the conjugate relations.

\paragraph{Order $g$.} The terms of order $g$ in $\Omega^\dagger\Omega$ are linear in the coefficients $\mf A^{(1)}_{1,2}$ and $\mf C_{1,2}$. The coefficient of $a^{\dagger b}_j(p)$ is $\mf A^{(1)b}_{2j}(p)+\mf A^{(1)\dagger b}_{1j}(p)$, and the coefficient of $a^{\dagger d}_l(q)\,a^{\dagger c}_k(p-q)\,a^b_j(p)$ is $\mf C^{dcb}_{1lkj}(p,q)+\mf C^{\dagger dcb}_{2lkj}(p,q)$. Hence
\begin{align}
\mf A^{(1)\dagger b}_{1j}(p)&=-\mf A^{(1)b}_{2j}(p)\,,\label{eq:A11unit}\\
\mf C^{\dagger dcb}_{2lkj}(p,q)&=-\mf C^{dcb}_{1lkj}(p,q)\,.\label{eq:C2unit}
\end{align}
At this order $\Omega=1+\Omega^{(1)}$ with $\Omega^{(1)}$ anti-Hermitian. Eq.~\eqref{eq:A11unit} is satisfied by the results of the diagrammatic calculation, Eqs.~\eqref{eq:A1mom} and~\eqref{eq:A11mom}. Eq.~\eqref{eq:C2unit} determines $\mf C_2$, the amplitude for the merging of two gluons, which we have not computed directly.

\paragraph{Order $g^2$.} At this order the terms linear in $\mf N_2$ and $\mf B_{1,2,3}$ are accompanied by the products of two coefficients of order $g$, normal ordered. The coefficient of the identity receives $\mf N_2+\mf N_2^\dagger$ and the contraction of $\mf A^{(1)\dagger}_2a$ with $\mf A^{(1)}_2a^\dagger$:
\begin{align}\label{eq:Nunit}
\mf N_2+\mf N_2^\dagger=-\int\frac{d^3p}{(2\pi)^3}\,\mf A^{(1)\dagger b}_{2j}(p)\,\mf A^{(1)b}_{2j}(p)\,.
\end{align}
Unitarity fixes the Hermitian part of $\mf N_2$ only. With the canonical normalisation of Section~\ref{sec:lcpt}, $\mf N_2$ is Hermitian, and
\begin{align}\label{eq:N2unit}
\mf N_2=-\frac12\int\frac{d^3p}{(2\pi)^3}\,\mf A^{(1)\dagger b}_{2j}(p)\,\mf A^{(1)b}_{2j}(p)
=-\frac{g^2}{8\pi^3}\int_\Lambda^\vee\frac{dp^+}{p^+}\int d^2p\,\frac{\rho^b(\v p)\,\rho^b(-\v p)}{\v p^2}\,,
\end{align}
where Eq.~\eqref{eq:A1mom} was inserted. This reproduces Eq.~\eqref{eq:N2mom}: the coefficient of the identity is minus one half of the probability to find one soft gluon. The coefficient of $a^{\dagger b}_j(p)\,a^{\dagger c}_k(q)$, symmetrised according to Eq.~\eqref{eq:symm}, gives
\begin{align}\label{eq:B1unit}
\mf B^{\dagger cb}_{1kj}(p,q)=\ &-\mf B^{cb}_{2kj}(p,q)+\mf C^{cbd}_{1kjl}(p+q,q)\,\mf A^{(1)d}_{2l}(p+q)\notag\\
&+\frac12\Big[\mf A^{(1)b}_{2j}(p)\,\mf A^{(1)c}_{2k}(q)+\mf A^{(1)c}_{2k}(q)\,\mf A^{(1)b}_{2j}(p)\Big]\,,
\end{align}
where Eqs.~\eqref{eq:A11unit} and~\eqref{eq:C2unit} were used. It determines $\mf B_1$, the amplitude for the absorption of two gluons. The second term is the absorption of one gluon by the valence charge after the two gluons have merged; the last term is the symmetrised product of two absorptions. Finally, the coefficient of $a^{\dagger b}_j(p)\,a^c_k(q)$ gives
\begin{align}\label{eq:B3unit}
&\mf B^{\dagger cb}_{3kj}(p,q)+\mf B^{bc}_{3jk}(q,p)+\mf A^{(1)b}_{2j}(p)\,\mf A^{(1)\dagger c}_{2k}(q)+\mf A^{(1)\dagger c}_{2k}(q)\,\mf A^{(1)b}_{2j}(p)\notag\\
&\quad+2\,\mf A^{(1)\dagger d}_{2l}(q-p)\,\mf C^{bdc}_{1jlk}(q,p)+2\,\mf C^{\dagger cdb}_{1klj}(p,q)\,\mf A^{(1)d}_{2l}(p-q)\notag\\
&\quad+2\,(2\pi)^3\delta^{(3)}(p-q)\int\frac{d^3r}{(2\pi)^3}\,\mf C^{\dagger edb}_{1mlj}(p,r)\,\mf C^{edc}_{1mlk}(p,r)=0\,.
\end{align}
The first term of the second line is present for $q^+>p^+$ and the second for $p^+>q^+$. Eq.~\eqref{eq:B3unit} contains only coefficients that are fixed by the diagrammatic calculation. It is a check of $\mf B_3$, and it is satisfied by Eq.~\eqref{eq:B3mom} term by term:
\begin{itemize}
\item The instantaneous term and the commutator term of Eq.~\eqref{eq:B3mom} cancel between $\mf B_3$ and $\mf B_3^\dagger$, as shown in Section~\ref{sec:summarywf}.
\item The anticommutator term gives
\begin{align}
\mf B^{bc}_{3jk}(q,p)+\mf B^{\dagger cb}_{3kj}(p,q)=-\frac{2g^2\,p^jq^k}{\v p^2\v q^2\sqrt{p^+q^+}}\,\big\{\rho^b(-\v p),\rho^c(\v q)\big\}\,,
\end{align}
which is minus the sum of the two products of $\mf A^{(1)}_2$ in the first line, by Eq.~\eqref{eq:A1mom}.
\item The terms with the three-gluon vertex cancel the second line. This fixes, in particular, the relative sign of the term $\propto\Theta_>$ of Eq.~\eqref{eq:B3mom} and $\mf C_1$.
\item The last line is cancelled by the $\rho$-independent term~\eqref{eq:B3rho0}.
\end{itemize}
The origin of these cancellations is an identity for the second-order terms of Eq.~\eqref{eq:pt2}. For two states $i\neq f$ with $E_i\neq E_f$ and an intermediate state $m$,
\begin{align}\label{eq:adaggera_identity}
\frac{V_{fm}V_{mi}}{(E_i-E_f)(E_i-E_m)}+\Big[\frac{V_{im}V_{mf}}{(E_f-E_i)(E_f-E_m)}\Big]^\dagger=-\Big[\frac{V_{mf}}{E_f-E_m}\Big]^\dagger\,\frac{V_{mi}}{E_i-E_m}\,,
\end{align}
which follows from $V_{mf}^\dagger=V_{fm}$ and $\frac{1}{E_i-E_m}-\frac{1}{E_f-E_m}=\frac{E_f-E_i}{(E_i-E_m)(E_f-E_m)}$, with the order of the operators kept. The left-hand side is the contribution of the state $m$ to $\mf B_3+\mf B_3^\dagger$, the right-hand side is minus the product of two first-order amplitudes. With the vacuum as intermediate state the right-hand side is the product of two absorption amplitudes, with a two-gluon intermediate state it is the product of two emission amplitudes or of an emission and a splitting amplitude. The pole at $E_i=E_f$ cancels on the left-hand side. Eq.~\eqref{eq:adaggera_identity} holds point by point for $E_i\neq E_f$ and therefore for the principal-value prescription~\eqref{eq:PV}.

\paragraph{Order $g^3$.} At this order we need only the coefficient of $a^{\dagger b}_j(p)$. It contains $\mf A^{(3)}_2+\mf A^{(3)\dagger}_1$, the products of $\mf N_2$ with the coefficients of order $g$, and the contractions of $\mf A^{(1)\dagger}_2a$ with $\mf B_2a^\dagger a^\dagger$, of $\mf B_3^\dagger a^\dagger a$ with $\mf A^{(1)}_2a^\dagger$, and of $\mf C_1^\dagger a^\dagger aa$ with $\mf B_2a^\dagger a^\dagger$; the result is Eq.~\eqref{eq:A31unit} of the main text.
The commutator arises as follows: the two terms with $\mf N_2$ are $\mf N_2^\dagger\,\mf A^{(1)}_2$, from $\mf N^\dagger$ in $\Omega^\dagger$ and $\mf A_2a^\dagger$ in $\Omega$, and $\mf A^{(1)\dagger}_1\,\mf N_2=-\mf A^{(1)}_2\,\mf N_2$, from $\mf A_1^\dagger a^\dagger$ in $\Omega^\dagger$ and $\mf N$ in $\Omega$. Eq.~\eqref{eq:A31unit} determines $\mf A^{(3)}_1$, the amplitude for the absorption of a gluon at order $g^3$, in terms of coefficients that are fixed by the diagrammatic calculation. In the remaining relations of order $g^3$ the coefficients denoted by dots in Eq.~\eqref{eq:OmegaExp} appear; they are not needed.

\paragraph{The order of the operators.} For commuting charges the commutator in Eq.~\eqref{eq:A31unit} is absent. Here it is not: $\mf N_2$ is bilinear in the charges at two different points, and with Eqs.~\eqref{eq:rhoalgebra}, \eqref{eq:N2mom} and~\eqref{eq:A1mom}
\begin{align}\label{eq:commN2A}
\big[\mf N_2,\mf A^{(1)a}_{2i}(k)\big]=-\log\frac{\vee}{\Lambda}\ \frac{\sqrt2\,g^3\,k^i}{8\pi^3\sqrt{k^+}\,\v k^2}\int d^2p\ \frac{if^{abc}\,\{\rho^b(-\v p),\rho^c(-\v k+\v p)\}}{\v p^2}\,.
\end{align}
This is a term with two charges, of the same order in $g$ as the terms with three charges from which it originates. Comparing with Eq.~\eqref{eq:A3tworho}, it is equal to the term $-2k^i\log(\vee/\Lambda)$ of $[\mf A^{(3)}_2]_{\rho\rho}$. Two consequences follow. First, the order of the charges in the three-$\rho$ term~\eqref{eq:A3threerho} is not a matter of convention: writing it as $\mf N_2\,\mf A^{(1)}_2$ instead of $\mf A^{(1)}_2\,\mf N_2$ would change $\mf A^{(3)}_2$ by the two-$\rho$ term~\eqref{eq:commN2A}. Second, with the order of Eq.~\eqref{eq:A3threerho} the first two terms of Eq.~\eqref{eq:A31unit} combine into Eq.~\eqref{eq:N2Acomb},
so that the three-$\rho$ parts of the emission and of the absorption amplitude are $\mf A^{(1)}_2\mf N_2$ and $-(\mf N_2\,\mf A^{(1)}_2)^\dagger=\mf A^{(1)}_1\,\mf N_2$, respectively. In both, the amplitude of order $g$ stands to the left of the norm factor, in agreement with Eq.~\eqref{eq:pt3} and its analogue for the one-gluon state.


\paragraph{Verification.} The relations of this appendix were verified by a symbolic normal ordering of $\Omega^\dagger\Omega$ with non-commuting coefficients for three discrete modes (script \texttt{scripts/check\_unitarity.py}): after the substitution of Eqs.~\eqref{eq:A11unit}, \eqref{eq:C2unit}, \eqref{eq:N2unit}, \eqref{eq:B1unit} and~\eqref{eq:A31unit} the coefficients of $1$, $a^\dagger$, $a^\dagger a^\dagger$ and, at order $g$, $a^\dagger a^\dagger a$ vanish identically, and the coefficient of $a^\dagger a$ is Eq.~\eqref{eq:B3unit}. The triple-vertex part of Eq.~\eqref{eq:B3unit} was checked numerically with Eqs.~\eqref{eq:B3mom}, \eqref{eq:C1mom} and~\eqref{eq:A1mom} (relative residual $10^{-16}$; with the opposite sign of the term $\propto\Theta_>$ the residual is of order one), and Eq.~\eqref{eq:commN2A} with $SU(2)$ charges on a lattice (script \texttt{scripts/check\_misc.py}).

% ======================================================================
\section{The coefficients of \texorpdfstring{$\Omega$}{Omega}: complete expressions}\label{app:coefficients}
% ======================================================================

This appendix collects the coefficients that are not displayed in full in the main text, in momentum space, and all coefficients in transverse coordinate space.

\subsection{Momentum space}

The coefficients $\mf A^{(1)}_2$, $\mf N_2$, $[\mf A^{(3)}_2]_\rho$, $[\mf A^{(3)}_2]_{\rho\rho\rho}$, $\mf A^{(1)}_1$, $\mf B_3$ and $\mf C_1$ are given in Eqs.~\eqref{eq:A1mom}, \eqref{eq:N2mom}, \eqref{eq:A3onerho}, \eqref{eq:A3threerho}, \eqref{eq:A11mom}, \eqref{eq:B3mom} and~\eqref{eq:C1mom}. The two-gluon coefficient is~\cite{Lublinsky:2016meo}
\begin{align}\label{eq:B2mom}
\mf B^{ba}_{2ji}(k,p)=\ &\frac{ig^2f^{abc}\,p^jk^i\rho^c(-\v p-\v k)\sqrt{p^+k^+}}{2\,(p^+\v k^2+k^+\v p^2)}\Big[\frac{1}{k^+\v p^2}-\frac{1}{p^+\v k^2}\Big]\,\Theta(\vee-k^+-p^+)
\notag\\&
+\frac{g^2p^jk^i\,\{\rho^b(-\v p),\rho^a(-\v k)\}}{2\,\v p^2\v k^2\sqrt{p^+k^+}}\,\Theta(\vee-k^+)\,\Theta(\vee-p^+)
\notag\\&
-\frac{ig^2f^{abc}\rho^c(-\v p-\v k)\,(p+k)^l\,(p^mk^+-p^+k^m)}{(\v p+\v k)^2\sqrt{p^+k^+}\,(p^++k^+)\big(\frac{\v p^2}{p^+}+\frac{\v k^2}{k^+}\big)}\notag\\&\quad\times\Big[\frac{\delta_{lm}\delta_{ji}}{p^++k^+}-\frac{\delta_{im}\delta_{jl}}{k^+}-\frac{\delta_{jm}\delta_{il}}{p^+}\Big]\Theta(\vee-k^+-p^+)
\notag\\&
-\frac{ig^2f^{abc}\rho^c(-\v p-\v k)\,(p^+-k^+)\,\delta_{ji}}{2\sqrt{p^+k^+}\,(p^++k^+)^2\big(\frac{\v k^2}{k^+}+\frac{\v p^2}{p^+}\big)}\,\Theta(\vee-k^+-p^+)\,,
\end{align}
for $k^+,p^+>\Lambda$. The first term is the commutator part of the two emissions, the second the anticommutator part, the third the splitting and the fourth the instantaneous emission. The supports in the longitudinal momenta follow from the kinematic restrictions of Section~\ref{sec:lcpt}: in the splitting the parent gluon of momentum $k^++p^+$ is a soft mode, and the same support is used for the instantaneous term; two independently emitted gluons are each soft modes, with no condition on $k^++p^+$.

The part of $\mf A^{(3)}_2$ with two charges is
\begin{align}\label{eq:A3tworho}
\Big[\mf A^{(3)a}_{2i}(k)\Big]_{\rho\rho}=\ &\int d^2p\,\frac{ig^3f^{abc}\{\rho^b(-\v p),\rho^c(-\v k+\v p)\}}{8\sqrt2\,\pi^3\,\v p^2\,\v k^2\sqrt{k^+}}
\Bigg[\log\frac{\Lambda}{k^+}\,\frac{\v k^2p^i-(\v k\cdot\v p)\,k^i}{(\v k-\v p)^2}
-2k^i\log\frac{\vee}{\Lambda}
\notag\\&
-2k^i\log\frac{k^+\v p^2}{\vee\,\v k^2+k^+\v p^2}
-\frac{\v p^2k^i+\v k^2p^i}{(\v k-\v p)^2}\log\frac{\vee-k^+}{\Lambda}-k^i\log\frac{\vee}{k^+}
\notag\\&
+\frac{(2\v p^2-2\v k\cdot\v p+\v k^2)\,k^i+\v k^2p^i}{(\v k-\v p)^2}\log\Big[1+\Big(\frac{\vee}{k^+}-1\Big)\frac{\v k^2}{(\v k-\v p)^2}\Big]\Bigg]\,.
\end{align}
It collects the anticommutator parts of the diagrams with two and three valence vertices; the limit $\Lambda\to0$ has been taken wherever possible.

\subsection{Transverse coordinate space}

In the cross section the coefficients are used in transverse coordinate space, where the eikonal scattering on the target is diagonal. The transverse arguments are defined by the Fourier conventions~\eqref{eq:FTconventions}: $\mf A_2(p^+,-\v x)$ and $\mf A_1(p^+,\v x)$ multiply $a^\dagger(p^+,\v x)$ and $a(p^+,\v x)$, and in general an argument conjugate to the momentum of a creation operator appears with a minus sign. Eq.~\eqref{eq:OmegaFinal} becomes
\begin{align}\label{eq:OmegaX}
\Omega = \ &\mf N
+\int\frac{dp^+}{2\pi}\int_x\Big[\mf A^{(1)b}_{1j}(p^+,\v x)\,a^b_j(p^+,\v x)-\mf A^{(1)\dagger b}_{1j}(p^+,\v x)\,a^{\dagger b}_j(p^+,\v x)
\notag\\&\hspace{2.6cm}
+\mf A^{(3)b}_{1j}(p^+,\v x)\,a^b_j(p^+,\v x)+\mf A^{(3)b}_{2j}(p^+,-\v x)\,a^{\dagger b}_j(p^+,\v x)\Big]
\notag\\&
+\int\frac{dp^+}{2\pi}\frac{dq^+}{2\pi}\int_{x,y}\Big[\mf B^{cb}_{1kj}(p^+,\v x;q^+,\v y)\,a^b_j(p^+,\v x)a^c_k(q^+,\v y)
\notag\\&\hspace{3.2cm}
+\mf B^{cb}_{2kj}(p^+,-\v x;q^+,-\v y)\,a^{\dagger b}_j(p^+,\v x)a^{\dagger c}_k(q^+,\v y)
\notag\\&\hspace{3.2cm}
+\mf B^{cb}_{3kj}(p^+,\v x;q^+,-\v y)\,a^{\dagger c}_k(q^+,\v y)a^b_j(p^+,\v x)\Big]
\notag\\&
+\int\frac{dp^+}{2\pi}\frac{dq^+}{2\pi}\int_{x,y,z}\Big[\mf C^{dcb}_{1lkj}(p^+,\v z-\v y;q^+,\v y-\v x)
\notag\\&\hspace{3.2cm}\times a^{\dagger d}_l(q^+,\v x)a^{\dagger c}_k(p^+-q^+,\v y)a^b_j(p^+,\v z)
\notag\\&\hspace{3.2cm}
-\mf C^{\dagger dcb}_{1lkj}(p^+,\v z-\v y;q^+,\v y-\v x)
\notag\\&\hspace{3.2cm}\times a^{\dagger b}_j(p^+,\v z)a^c_k(p^+-q^+,\v y)a^d_l(q^+,\v x)\Big]+\dots
\end{align}
The coefficients fixed by the diagrammatic calculation are the following. The normalisation and the amplitude of order $g$ are
\begin{align}
\mf N_2&=-\frac{g^2}{8\pi^3}\,\log\frac{\vee}{\Lambda}\int_{x,y,z}K(\v x,\v y,\v z)\,\rho^a(\v x)\,\rho^a(\v y)\,,\label{eq:N2X}\\
\mf A^{(1)a}_{2i}(k^+,-\v z)&=-\mf A^{(1)\dagger a}_{1i}(k^+,\v z)=\frac{ig}{\sqrt2\,\pi\sqrt{k^+}}\int_x\Kv^i(\v x-\v z)\,\rho^a(\v x)\,.\label{eq:A1X}
\end{align}
The amplitude of order $g^3$ for the emission of one gluon has the three parts
\begin{align}\label{eq:A3onerhoX}
\Big[\mf A^{(3)a}_{2i}(k^+,-\v z)\Big]_{\rho}=\ &-\int_x\frac{ig^3N_c\,\rho^a(\v x)\,\Kv^i(\v x-\v z)}{16\sqrt2\,\pi^3\sqrt{k^+}}\Bigg\{\Big[-\frac2\epsilon-2\gamma_E-\log\frac{(\v x-\v z)^2\mu^2_{\MS}}{4}\Big]
\notag\\&\hspace{0.5cm}\times\Big[\frac{11}{3}-2\log\frac{k^+}{\Lambda}+\log\Big(1-\frac{k^+}{\vee}\Big)\Big]
\notag\\&\hspace{0.5cm}
-\frac{67}{9}+\frac{\pi^2}{3}+\log^2\Big(1-\frac{k^+}{\vee}\Big)+\log\Big(1-\frac{k^+}{\vee}\Big)\log\frac{\vee}{k^+}\Bigg\},
\end{align}
where $\gamma_E$ is the Euler constant,
\begin{align}\label{eq:A3tworhoX}
\Big[\mf A^{(3)a}_{2i}(k^+,-\v z)\Big]_{\rho\rho}=\ &\int d^2p\,d^2k\int_{x,y}e^{i\v k\cdot(\v z-\v y)}e^{i\v p\cdot(\v y-\v x)}\,\notag\\&\times\frac{ig^3f^{abc}\{\rho^b(\v x),\rho^c(\v y)\}}{32\sqrt{2k^+}\,\pi^5\,\v p^2\,\v k^2}\,\Big[\ \cdots\ \Big]\,,
\end{align}
with the same square bracket as in Eq.~\eqref{eq:A3tworho} ($\v k$ and $\v p$ are integration variables here), and
\begin{align}\label{eq:A3threerhoX}
\Big[\mf A^{(3)a}_{2i}(k^+,-\v z)\Big]_{\rho\rho\rho}=\ &-\log\frac{\vee}{\Lambda}\int_{w,x,y,u}\rho^a(\v w)\,\rho^b(\v x)\,\rho^b(\v y)\,K(\v x,\v y,\v u)\,\Kv^i(\v w-\v z)\notag\\&\times\frac{ig^3}{8\sqrt2\,\pi^4\sqrt{k^+}}=\mf A^{(1)a}_{2i}(k^+,-\v z)\,\mf N_2\,.
\end{align}
The amplitude for two gluons with longitudinal momenta $k^+$, $p^+$ at the transverse positions $\v w$, $\v z$ is
\begin{align}\label{eq:B2X}
\mf B^{bc}_{2ji}(k^+,-\v w;p^+,-\v z)=\ &\int_x\frac{ig^2f^{cba}\,\rho^a(\v x)\,\sqrt{p^+k^+}}{4\pi^2\big[p^+(\v x-\v z)^2+k^+(\v x-\v w)^2\big](p^++k^+)}
\notag\\&\times\Bigg[\frac{\delta^{ij}\big[(\v x-\v z)^2-(\v x-\v w)^2\big]}{2(\v z-\v w)^2}
\notag\\&\quad
+\frac{p^++k^+}{k^+}\Big(\frac{(x-z)^j(z-w)^i}{(\v z-\v w)^2}-\frac{(x-z)^j(x-w)^i}{2(\v x-\v w)^2}\Big)
\notag\\&\quad
+\frac{p^++k^+}{p^+}\Big(\frac{(x-w)^i(z-w)^j}{(\v z-\v w)^2}+\frac{(x-z)^j(x-w)^i}{2(\v x-\v z)^2}\Big)\Bigg]
\notag\\&\quad\times\Theta(\vee-k^+-p^+)
\notag\\&
-\int_{x,y}\frac{g^2\,\{\rho^c(\v y),\rho^b(\v x)\}\,\Kv^i(\v y-\v w)\,\Kv^j(\v x-\v z)}{8\pi^2\sqrt{p^+k^+}}
\notag\\&\quad\times\Theta(\vee-k^+)\,\Theta(\vee-p^+)\,.
\end{align}
The amplitude for the transition of a gluon $(k^+,\v z,a,i)$ into a gluon $(p^+,\v u,b,j)$ on the valence charges is
\begin{align}\label{eq:B3X}
\mf B^{ba}_{3ji}(k^+,\v z;p^+,-\v u)=\ &\frac{ig^2f^{abc}}{(2\pi)^2}\int_x\frac{\sqrt{p^+k^+}\,(p^++k^+)}{(k^+-p^+)^2}\,
\notag\\&\qquad\times\frac{\rho^c(\v x)\,\delta_{ij}}{p^+(\v x-\v u)^2-k^+(\v x-\v z)^2}\,\big[\Theta_>+\Theta_<\big]
\notag\\&
+\frac{ig^2f^{abc}}{2\pi^2k^+}\int_x\frac{\sqrt{k^+p^+}}{k^+-p^+}\,\rho^c(\v x)
\notag\\&\times\frac{\big(k^+(x-z)^m-p^+(x-u)^m\big)\big((x-z)^l-(x-u)^l\big)}{k^+(\v x-\v z)^2-p^+(\v x-\v u)^2}
\notag\\&\qquad\times\Big(\delta_{il}\delta_{jm}-\frac{k^+}{p^+}\delta_{jl}\delta_{im}-\frac{k^+}{k^+-p^+}\delta_{ij}\delta_{lm}\Big)
\notag\\&\qquad\times\Bigg[-\frac{\Theta_>}{(\v u-\v z)^2}-\frac{p^+k^+\,\Theta_<}{\big(k^+(\v x-\v z)-p^+(\v x-\v u)\big)^2}\Bigg]
\notag\\&
-\frac{g^2}{4\pi^2\sqrt{p^+k^+}}\int_{x,y}\Kv^i(\v x-\v z)\,\Kv^j(\v y-\v u)\,\{\rho^a(\v x),\rho^b(\v y)\}
\notag\\&
-\frac{ig^2f^{abc}}{4\pi^2\sqrt{p^+k^+}}\int_x\Kv^i(\v x-\v z)\,\Kv^j(\v x-\v u)
\notag\\&\qquad\times\frac{k^+(\v x-\v z)^2+p^+(\v x-\v u)^2}{p^+(\v x-\v u)^2-k^+(\v x-\v z)^2}\,\rho^c(\v x)\,,
\end{align}
with principal values for the denominators $p^+(\v x-\v u)^2-k^+(\v x-\v z)^2$. It follows from Eqs.~\eqref{eq:psi_instX}, \eqref{eq:psi1X}, \eqref{eq:psi2X}, \eqref{eq:psi_antiX} and~\eqref{eq:psi_commX} with $\v X'=\v x-\v z$, $\v X=\v x-\v u$, $\xi=p^+/k^+$ and the factor $2\pi/k^+$ between the coefficient of $\ket{g^b_j(\xi k^+,\v u)}$ under $\int d\xi\int_u$ and $\mf B_3$. The amplitude for the splitting of a gluon $(k^+,a,i)$ into the gluons $(p^+,b,j)$ and $(k^+-p^+,c,k)$ is
\begin{align}\label{eq:C1X}
\mf C^{bca}_{1jki}(k^+,\v z;p^+,\v x)=\ &-\frac{gf^{abc}}{2\pi\sqrt{2k^+}}\,\frac{\sqrt{p^+(k^+-p^+)}}{k^+}\,\delta^{(2)}\Big(\v z+\frac{p^+}{k^+}\v x\Big)\,\Kv^m(\v x)
\notag\\&\times\Big[\delta_{im}\delta_{jk}-\frac{k^+}{p^+}\delta_{jm}\delta_{ik}-\frac{k^+}{k^+-p^+}\delta_{km}\delta_{ij}\Big]\Theta(k^+-p^+-\Lambda)\,,
\end{align}
where the arguments are those of Eq.~\eqref{eq:OmegaX}: $\v z$ is the position of the incoming gluon relative to the outgoing gluon $(c,k)$, and $\v x$ the position of the latter relative to the outgoing gluon $(b,j)$. The $\delta$-function expresses that the incoming gluon sits at the centre of the two outgoing gluons weighted with their longitudinal momenta.

The coefficients fixed by unitarity are, in coordinate space,
\begin{align}\label{eq:unitX}
\mf C^{\dagger dcb}_{2lkj}(p^+,\v z;q^+,\v x)&=-\mf C^{dcb}_{1lkj}(p^+,-\v z;q^+,-\v x)\,,\notag\\
\mf B^{\dagger cb}_{1kj}(p^+,\v x;q^+,\v y)&=-\mf B^{cb}_{2kj}(p^+,-\v x;q^+,-\v y)\notag\\
&\quad+\frac12\Big[\mf A^{(1)b}_{2j}(p^+,-\v x)\,\mf A^{(1)c}_{2k}(q^+,-\v y)\notag\\&\qquad\quad+\mf A^{(1)c}_{2k}(q^+,-\v y)\,\mf A^{(1)b}_{2j}(p^+,-\v x)\Big]
\notag\\&\quad
+\int_z\mf C^{cbd}_{1kjl}(q^++p^+,\v z-\v x;q^+,\v x-\v y)\,\mf A^{(1)d}_{2l}(q^++p^+,-\v z)\,,\notag\\
\mf A^{(3)\dagger b}_{1j}(p^+,\v x)&=-\mf A^{(3)b}_{2j}(p^+,-\v x)-\big[\mf N_2,\mf A^{(1)b}_{2j}(p^+,-\v x)\big]
\notag\\&\quad+\int\frac{dq^+}{2\pi}\int_y\Big[2\,\mf A^{(1)c}_{1k}(q^+,\v y)\,\mf B^{cb}_{2kj}(p^+,-\v x;q^+,-\v y)
\notag\\&\qquad
-\mf B^{\dagger cb}_{3kj}(p^+,\v x;q^+,-\v y)\,\mf A^{(1)c}_{2k}(q^+,-\v y)
\notag\\&\qquad-2\int_z\mf C^{\dagger dcb}_{1lkj}(p^+,\v x-\v z;q^+,\v z-\v y)\,\mf B^{dc}_{2lk}(p^+-q^+,-\v z;q^+,-\v y)\Big]\,.
\end{align}

\NOTE{Source of this appendix: the coordinate-space expressions~\eqref{eq:A3onerhoX}--\eqref{eq:C1X} and~\eqref{eq:unitX} are taken from the checked calculation file; the transforms from Eq.~\eqref{eq:B2mom} to Eq.~\eqref{eq:B2X} and from Eq.~\eqref{eq:C1mom} to Eq.~\eqref{eq:C1X}, and the transforms~\eqref{eq:FT_i4} and~\eqref{eq:FT_i5}, were not re-derived here.}


% ======================================================================
\section{The rows of the cross section in terms of the coefficients}\label{app:rows}
% ======================================================================

This appendix lists all rows of Table~\ref{tab:rows} in terms of the coefficients of $\Omega$, and describes the verification of the amplitudes. All coefficients are in transverse coordinate space, Appendix~\ref{app:coefficients}; the conjugated factor stands to the left, at the primed point $\v w'$; $\langle\dots\rangle$ is the average over the valence state.

\subsection{Group I}

We abbreviate the conjugate of the leading-order amplitude~\eqref{eq:ALO} at the point $\v w'$ by
\begin{align}\label{eq:Ldag}
\mathcal L^{\dagger a}_i(\v w')\equiv\mf A^{\dagger b'}_i(k^+,-\v w')\,U^{ab'}(\v w')-\bar{\mf A}^{\dagger a}_i(k^+,-\v w')\,.
\end{align}
Each row of Group~I has the form
\begin{align}\label{eq:GroupIform}
\frac{d^3N^{\rm(I.n)}}{d^3k}=\frac{1}{(2\pi)^3}\int_{w,w'}e^{-i\v k\cdot(\v w-\v w')}\,\Big\langle\mathcal L^{\dagger a}_i(\v w')\,\big[\mathcal A_{\rm n}\big]^a_i(\v w)\Big\rangle+{\rm c.c.}\,,
\end{align}
with the following operators $\mathcal A_{\rm n}$:
\begin{align}
\text{I.I:}\quad&-\bar{\mf A}^a_i(k^+,-\v w)\,\mf N_2+U^{ab}(\v w)\,\bar{\mf N}_2\,\mf A^b_i(k^+,-\v w)\notag\\&+U^{ab}(\v w)\,\mf A^b_i(k^+,-\v w)\,\mf N_2-\bar{\mf N}_2\,\bar{\mf A}^a_i(k^+,-\v w)\,,\label{eq:rowI_I}\\
\text{I.II:}\quad&U^{ab}(\v w)\,\Big[\mf A^{(3)b}_i(k^+,-\v w)\Big]_{\rho+\rho\rho}-\Big[\bar{\mf A}^{(3)a}_i(k^+,-\v w)\Big]_{\rho+\rho\rho}\,,\label{eq:rowI_II}\\
\text{I.III:}\quad&\int\frac{dp^+}{2\pi}\int_x\bar{\mf B}^{\dagger ba}_{3ji}(k^+,\v w;p^+,-\v x)\Big[U^{bc}(\v x)\,\mf A^c_j(p^+,-\v x)-\bar{\mf A}^b_j(p^+,-\v x)\Big]\,,\label{eq:rowI_IIIb}\\
\text{I.IV:}\quad&2\int\frac{dp^+}{2\pi}\int_x\Big[U^{ad}(\v w)U^{bc}(\v x)\,\bar{\mf A}^{\dagger b}_j(p^+,-\v x)\,\mf B^{cd}_{2ji}(k^+,-\v w;p^+,-\v x)\notag\\
&\hspace{2.2cm}-\bar{\mf A}^{\dagger b}_j(p^+,-\v x)\,\bar{\mf B}^{ba}_{2ji}(k^+,-\v w;p^+,-\v x)\Big]\,,\label{eq:rowI_IV}\\
\text{I.V:}\quad&2\int\frac{dp^+}{2\pi}\int_{x,z}\notag\\&\times\Big[\bar{\mf C}^{\dagger dca}_{1jki}(k^+,\v w-\v z;p^+,\v z-\v x)\,U^{db}(\v x)U^{ce}(\v z)\,\mf B^{be}_{2jk}(k^+-p^+,-\v z;p^+,-\v x)\notag\\
&\quad-\bar{\mf C}^{\dagger bca}_{1jki}(k^+,\v w-\v z;p^+,\v z-\v x)\,\bar{\mf B}^{bc}_{2jk}(k^+-p^+,-\v z;p^+,-\v x)\Big]\,.\label{eq:rowI_V}
\end{align}
In Eq.~\eqref{eq:rowI_II} the subscript denotes the parts of $\mf A^{(3)}$ with one and two charges, Eqs.~\eqref{eq:A3onerhoX} and~\eqref{eq:A3tworhoX}; the part with three charges, $\mf A\,\mf N_2$, is the third term of Eq.~\eqref{eq:rowI_I} and, for the barred coefficient, is combined with the commutator of $\mathcal A_N$ into the fourth term, Eq.~\eqref{eq:Nterms}. In $\bar{\mf B}_3^\dagger$ the first pair of arguments and the second pair of indices, $(k^+,\v w;a,i)$, belong to the measured gluon, which is the incoming gluon of $\bar{\mf B}_3$, Eq.~\eqref{eq:B3X}.

\subsection{Group II}

The conjugates of the four parts~\eqref{eq:A2parts} of the two-gluon amplitude, at the point $\v w'$ and with $(q^+,\v v)\to(p^+,\v z)$, are
\begin{align}\label{eq:A2partsdag}
\mathcal A_B^\dagger&=\mf B^{\dagger bc'}_{2ji}(k^+,-\v w';p^+,-\v z)\,U^{ac'}(\v w')-\bar{\mf B}^{\dagger c'a}_{2ji}(k^+,-\v w';p^+,-\v z)\,U^{c'b}(\v z)\,,\notag\\
\mathcal A^{(1)\dagger}_{AA}&=\tfrac12\Big[\bar{\mf A}^{\dagger a}_i(k^+,-\v w')\,\bar{\mf A}^{\dagger c'}_j(p^+,-\v z)\,U^{c'b}(\v z)\notag\\&\hspace{1.2cm}-\mf A^{\dagger c'}_i(k^+,-\v w')\,U^{d'b}(\v z)\,\bar{\mf A}^{\dagger d'}_j(p^+,-\v z)\,U^{ac'}(\v w')\Big]\,,\notag\\
\mathcal A^{(2)\dagger}_{AA}&=\tfrac12\Big[\bar{\mf A}^{\dagger c'}_j(p^+,-\v z)\,\bar{\mf A}^{\dagger a}_i(k^+,-\v w')\,U^{c'b}(\v z)-\mf A^{\dagger b}_j(p^+,-\v z)\,\bar{\mf A}^{\dagger a}_i(k^+,-\v w')\Big]\,,\notag\\
\mathcal A_{CA}^\dagger&=\int_{y'}\Big[\mf A^{\dagger e'}_{k'}(p^++k^+,-\v y')\,U^{c'e'}(\v y')\,U^{d'b}(\v z)\,\bar{\mf C}^{\dagger d'ac'}_{1jik'}(p^++k^+,\v y'-\v w';p^+,\v w'-\v z)\notag\\
&\hspace{1.2cm}-\bar{\mf A}^{\dagger c'}_{k'}(p^++k^+,-\v y')\,\bar{\mf C}^{\dagger d'ac'}_{1jik'}(p^++k^+,\v y'-\v w';p^+,\v w'-\v z)\,U^{d'b}(\v z)\Big]\,.
\end{align}
With these, and with the parts~\eqref{eq:A2parts} at the point $\v w$ (in $\mathcal A_{CA}$ the position of the parent gluon is denoted by $\v y$), the rows of Group~II are
\begin{align}\label{eq:GroupIIrows}
\frac{d^3N^{\rm(II.n)}}{d^3k}=\frac{4}{(2\pi)^3}\int\frac{dp^+}{2\pi}\int_{z,w,w'}e^{-i\v k\cdot(\v w-\v w')}\,\Big\langle\mathcal X^{\dagger}_{\rm n}\,\mathcal X_{\rm n}\Big\rangle\,,\qquad
\mathcal X_{\rm n}=\mathcal A_B,\ \mathcal A^{(1)}_{AA},\ \mathcal A^{(2)}_{AA},\ \mathcal A_{CA}
\end{align}
for ${\rm n}={\rm I},\dots,{\rm IV}$, with the sum over $a,b,i,j$ understood. Explicitly:
\begin{align}
\text{II.I:}\quad&\Big[\mf B^{\dagger bc'}_{2ji}(\v w';\v z)\,U^{ac'}(\v w')-\bar{\mf B}^{\dagger c'a}_{2ji}(\v w';\v z)\,U^{c'b}(\v z)\Big]\notag\\&\times\Big[U^{ac}(\v w)\,\mf B^{bc}_{2ji}(\v w;\v z)-U^{cb}(\v z)\,\bar{\mf B}^{ca}_{2ji}(\v w;\v z)\Big]\,,\notag\\
\text{II.II:}\quad&\tfrac14\Big[\bar{\mf A}^{\dagger a}_i(\v w')\,\bar{\mf A}^{\dagger c'}_j(\v z)\,U^{c'b}(\v z)-\mf A^{\dagger c'}_i(\v w')\,U^{d'b}(\v z)\,\bar{\mf A}^{\dagger d'}_j(\v z)\,U^{ac'}(\v w')\Big]\notag\\
&\times\Big[U^{cb}(\v z)\,\bar{\mf A}^c_j(\v z)\,\bar{\mf A}^a_i(\v w)-U^{ac}(\v w)\,\bar{\mf A}^d_j(\v z)\,U^{db}(\v z)\,\mf A^c_i(\v w)\Big]\,,\notag\\
\text{II.III:}\quad&\tfrac14\Big[\bar{\mf A}^{\dagger c'}_j(\v z)\,\bar{\mf A}^{\dagger a}_i(\v w')\,U^{c'b}(\v z)-\mf A^{\dagger b}_j(\v z)\,\bar{\mf A}^{\dagger a}_i(\v w')\Big]\notag\\&\times\Big[U^{cb}(\v z)\,\bar{\mf A}^a_i(\v w)\,\bar{\mf A}^c_j(\v z)-\bar{\mf A}^a_i(\v w)\,\mf A^b_j(\v z)\Big]\,,\notag\\
\text{II.IV:}\quad&\int_{y,y'}\Big[\mf A^{\dagger e'}_{k'}(\v y')\,U^{c'e'}(\v y')\,U^{d'b}(\v z)\,\bar{\mf C}^{\dagger d'ac'}_{1jik'}(\v y'-\v w';\v w'-\v z)\notag\\&\qquad\quad-\bar{\mf A}^{\dagger c'}_{k'}(\v y')\,\bar{\mf C}^{\dagger d'ac'}_{1jik'}(\v y'-\v w';\v w'-\v z)\,U^{d'b}(\v z)\Big]\notag\\
&\times\Big[\bar{\mf C}^{dac}_{1jik}(\v y-\v w;\v w-\v z)\,U^{db}(\v z)\,U^{ce}(\v y)\,\mf A^e_k(\v y)\notag\\&\qquad\quad-U^{db}(\v z)\,\bar{\mf C}^{dac}_{1jik}(\v y-\v w;\v w-\v z)\,\bar{\mf A}^c_k(\v y)\Big]\,.\label{eq:GroupIIexplicit}
\end{align}
Here and in the next subsection the longitudinal arguments are suppressed: $\mf A_i(\v w)\equiv\mf A_i(k^+,-\v w)$ for the measured gluon, $\mf A_j(\v z)\equiv\mf A_j(p^+,-\v z)$ for the second gluon, $\mf A_k(\v y)\equiv\mf A_k(p^++k^+,-\v y)$ for the parent gluon, $\mf B_{2ji}(\v w;\v z)\equiv\mf B_{2ji}(k^+,-\v w;p^+,-\v z)$ and $\bar{\mf C}_{1jik}(\v y-\v w;\v w-\v z)\equiv\bar{\mf C}_{1jik}(p^++k^+,\v y-\v w;p^+,\v w-\v z)$.

\subsection{Group III}

The rows of Group~III are the six interferences,
\begin{align}\label{eq:GroupIIIrows}
\frac{d^3N^{\rm(III.n)}}{d^3k}=\frac{4}{(2\pi)^3}\int\frac{dp^+}{2\pi}\int_{z,w,w'}e^{-i\v k\cdot(\v w-\v w')}\,\Big\langle\mathcal X^{\dagger}_{\rm n}\,\mathcal Y_{\rm n}\Big\rangle+{\rm c.c.}\,,
\end{align}
with
\begin{align}\label{eq:GroupIIIlist}
&\text{III.I:}\ \ \mathcal A^{(1)\dagger}_{AA}\,\mathcal A_B\,,&&\text{III.II:}\ \ \mathcal A^{(2)\dagger}_{AA}\,\mathcal A_B\,,&&\text{III.III:}\ \ -\mathcal A^{\dagger}_{CA}\,\mathcal A_B\,,\notag\\
&\text{III.IV:}\ \ \mathcal A^{(2)\dagger}_{AA}\,\mathcal A^{(1)}_{AA}\,,&&\text{III.V:}\ \ -\mathcal A^{\dagger}_{CA}\,\mathcal A^{(1)}_{AA}\,,&&\text{III.VI:}\ \ -\mathcal A^{\dagger}_{CA}\,\mathcal A^{(2)}_{AA}\,.
\end{align}
Explicitly, in the notation of Eq.~\eqref{eq:GroupIIexplicit}:
\begin{align}
\text{III.I:}\quad&\tfrac12\Big[\bar{\mf A}^{\dagger a}_i(\v w')\,\bar{\mf A}^{\dagger c'}_j(\v z)\,U^{c'b}(\v z)-\mf A^{\dagger c'}_i(\v w')\,U^{d'b}(\v z)\,\bar{\mf A}^{\dagger d'}_j(\v z)\,U^{ac'}(\v w')\Big]\notag\\
&\times\Big[U^{ac}(\v w)\,\mf B^{bc}_{2ji}(\v w;\v z)-U^{cb}(\v z)\,\bar{\mf B}^{ca}_{2ji}(\v w;\v z)\Big]\,,\notag\\
\text{III.II:}\quad&\tfrac12\Big[\bar{\mf A}^{\dagger c'}_j(\v z)\,\bar{\mf A}^{\dagger a}_i(\v w')\,U^{c'b}(\v z)-\mf A^{\dagger b}_j(\v z)\,\bar{\mf A}^{\dagger a}_i(\v w')\Big]\notag\\&\times\Big[U^{ac}(\v w)\,\mf B^{bc}_{2ji}(\v w;\v z)-U^{cb}(\v z)\,\bar{\mf B}^{ca}_{2ji}(\v w;\v z)\Big]\,,\notag\\
\text{III.III:}\quad&-\int_{y'}\Big[\mf A^{\dagger e'}_{k'}(\v y')\,U^{c'e'}(\v y')\,U^{d'b}(\v z)\,\bar{\mf C}^{\dagger d'ac'}_{1jik'}(\v y'-\v w';\v w'-\v z)\notag\\&\qquad\quad-\bar{\mf A}^{\dagger c'}_{k'}(\v y')\,\bar{\mf C}^{\dagger d'ac'}_{1jik'}(\v y'-\v w';\v w'-\v z)\,U^{d'b}(\v z)\Big]\notag\\
&\times\Big[U^{ac}(\v w)\,\mf B^{bc}_{2ji}(\v w;\v z)-U^{cb}(\v z)\,\bar{\mf B}^{ca}_{2ji}(\v w;\v z)\Big]\,,\notag\\
\text{III.IV:}\quad&\tfrac14\Big[\bar{\mf A}^{\dagger c'}_j(\v z)\,\bar{\mf A}^{\dagger a}_i(\v w')\,U^{c'b}(\v z)-\mf A^{\dagger b}_j(\v z)\,\bar{\mf A}^{\dagger a}_i(\v w')\Big]\notag\\
&\times\Big[U^{cb}(\v z)\,\bar{\mf A}^c_j(\v z)\,\bar{\mf A}^a_i(\v w)-U^{ac}(\v w)\,\bar{\mf A}^d_j(\v z)\,U^{db}(\v z)\,\mf A^c_i(\v w)\Big]\,,\notag\\
\text{III.V:}\quad&-\tfrac12\int_{y'}\Big[\mf A^{\dagger e'}_{k'}(\v y')\,U^{c'e'}(\v y')\,U^{d'b}(\v z)\,\bar{\mf C}^{\dagger d'ac'}_{1jik'}(\v y'-\v w';\v w'-\v z)\notag\\&\qquad\quad-\bar{\mf A}^{\dagger c'}_{k'}(\v y')\,\bar{\mf C}^{\dagger d'ac'}_{1jik'}(\v y'-\v w';\v w'-\v z)\,U^{d'b}(\v z)\Big]\notag\\
&\times\Big[U^{cb}(\v z)\,\bar{\mf A}^c_j(\v z)\,\bar{\mf A}^a_i(\v w)-U^{ac}(\v w)\,\bar{\mf A}^d_j(\v z)\,U^{db}(\v z)\,\mf A^c_i(\v w)\Big]\,,\notag\\
\text{III.VI:}\quad&-\tfrac12\int_{y'}\Big[\mf A^{\dagger e'}_{k'}(\v y')\,U^{c'e'}(\v y')\,U^{d'b}(\v z)\,\bar{\mf C}^{\dagger d'ac'}_{1jik'}(\v y'-\v w';\v w'-\v z)\notag\\&\qquad\quad-\bar{\mf A}^{\dagger c'}_{k'}(\v y')\,\bar{\mf C}^{\dagger d'ac'}_{1jik'}(\v y'-\v w';\v w'-\v z)\,U^{d'b}(\v z)\Big]\notag\\
&\times\Big[U^{cb}(\v z)\,\bar{\mf A}^a_i(\v w)\,\bar{\mf A}^c_j(\v z)-\bar{\mf A}^a_i(\v w)\,\mf A^b_j(\v z)\Big]\,.\label{eq:GroupIIIexplicit}
\end{align}
The sum of Eqs.~\eqref{eq:GroupIIrows} and~\eqref{eq:GroupIIIrows} over all rows is $\frac{4}{(2\pi)^3}\int\frac{dp^+}{2\pi}$ times the square of the sum $\mathcal A_B+\mathcal A^{(1)}_{AA}+\mathcal A^{(2)}_{AA}-\mathcal A_{CA}$: all sixteen products are present. The rows were compared one by one with the expressions from which the results of Sections~\ref{sec:rapidity}--\ref{sec:finite} were derived (weights $1$, $1$, $2$, $2$ for Rows~I.II--I.V; $4$ and $4\cdot\frac14$ for Group~II; $2$, $2$, $-4$, $4\cdot\frac14$, $-2$, $-2$ for Group~III, all in units of $(2\pi)^{-3}\int\frac{dp^+}{2\pi}$; placement of the Wilson lines; phase; ranges of $p^+$). They agree. In those expressions the three-$\rho$ part of $\mathcal A_{A^{(3)}}$ is kept in the second row of Group~I; here it is assigned to Row~I.I, Eq.~\eqref{eq:Nterms}, which does not change the sum.

\subsection{Verification of the amplitudes}\label{app:rows_check}

The amplitudes~\eqref{eq:A3parts} and~\eqref{eq:A2parts}, from which all rows are built, were verified by an independent symbolic expansion (script \texttt{scripts/check\_amplitudes.py}). In a model with three discrete modes, $[a_p,a^\dagger_q]=\delta_{pq}$, the operator $\Omega$ is written as in Eq.~\eqref{eq:OmegaExp} with formal non-commuting coefficients, and $\bar\Omega$ with a second, independent set of barred coefficients and with every $a$, $a^\dagger$ rotated by a real orthogonal matrix $U$ on the mode index; $\mf C_{1,2}$ are the same in both. The product $\bar\Omega^\dagger\Omega\ket0$ is brought to normal order, the unitarity relations of Appendix~\ref{app:unitarity} are applied to the barred coefficients, and the projections on the one-gluon states at order $g^3$ and on the two-gluon states at order $g^2$, multiplied by $U$ for the measured gluon, are compared with Eqs.~\eqref{eq:A3parts} and~\eqref{eq:A2parts}, word by word in the non-commuting coefficients. The one-gluon amplitude contains 255 and the two-gluon amplitude 504 distinct operator words (summed over the external indices); the maximal deviation of a numerical coefficient is $4\times10^{-16}$ and $2\times10^{-16}$, respectively. The check covers the signs, the factors $2$ and $\frac12$, the placement of the Wilson lines, the order of the coefficients, and the commutator term of $\mathcal A_N$. It does not cover the continuum normalisation factors $(2\pi)^n$ and the transverse arguments of the coefficients, which were checked by hand against Eqs.~\eqref{eq:OmegaX} and~\eqref{eq:states}, nor the assembly of the amplitudes into rows, which is the multiplication in Eqs.~\eqref{eq:Groups}.


\bibliographystyle{JHEP}
\bibliography{refs}

\end{document}
